% ============================================================================
% Fichier     : CASI_main.tex
% Titre       : Contrôle et Audit des Systèmes d'Information
% Auteur      : MINKA MI NGUIDJOÏ Thierry Emmanuel
% Institution : PSSFP — Université de Yaoundé II
% Version     : 0.1 -- novembre 2026 (squelette)
% ============================================================================
% UE CAS3401, Master professionnel en finances publiques, spécialité Audit et
% Contrôle, niveau 5. 8 jours de 4 h ; 30 h effectives (les 2 dernières
% heures du J08 sont offertes). Prologue + C01 à C14, numérotation continue.
% ============================================================================

\documentclass[12pt, twoside, openright, a4paper]{book}

\usepackage{Style_CASI}
\raggedbottom

\addbibresource{bibliographie/CASI_references.bib}

\hypersetup{
    pdftitle={Contrôle et Audit des Systèmes d'Information},
    pdfauthor={MINKA MI NGUIDJOÏ Thierry Emmanuel},
    pdfsubject={CAS3401 · Master professionnel en finances publiques · PSSFP, Université de Yaoundé II},
    pdfkeywords={contrôle interne, évaluation du contrôle interne, audit des SI, finances publiques, CONSUPE, COSO, INTOSAI}
}

\begin{document}

\frontmatter

    % ==============================================================================
% Fichier : front/00_page_titre.tex
% Titre   : Page de titre - Contrôle et Audit des Systèmes d'Information
% Auteur  : MINKA MI NGUIDJOÏ Thierry Emmanuel
% Version : 0.1 -- novembre 2026 (squelette)
% ==============================================================================
% Logos : Université de Yaoundé II à gauche, PSSFP à droite.
% ==============================================================================

\begin{titlepage}
\thispagestyle{empty}

% --- Fond décoratif : bandeaux bleu marine et filets or ---
\begin{tikzpicture}[remember picture, overlay]
    \fill[CouleurPrimaire]
        (current page.north west) rectangle
        ([yshift=-3.6cm] current page.north east);
    \fill[CouleurAccent]
        ([yshift=-3.6cm] current page.north west) rectangle
        ([yshift=-4.0cm] current page.north east);
    \fill[CouleurPrimaire]
        (current page.south west) rectangle
        ([yshift=2.5cm] current page.south east);
    \fill[CouleurAccent]
        ([yshift=2.5cm] current page.south west) rectangle
        ([yshift=2.9cm] current page.south east);
    % En-tête institutionnel, posé dans le bandeau supérieur
    \node[align=center, text=white, text width=12.5cm] at ([yshift=-1.8cm] current page.north) {%
        {\fontsize{11}{13}\selectfont
        \textsc{République du Cameroun}}\\[0.3ex]
        {\color{CouleurAccent!60!white}\rule{0.4\paperwidth}{0.4pt}}\\[0.6ex]
        {\fontsize{10}{12}\selectfont\textbf{\tutelle}}\\[0.2ex]
        {\fontsize{9}{11}\selectfont \institutionprimaire{} (PSSFP)}\\[0.2ex]
        {\fontsize{8}{10}\selectfont \unitepedago}};
    % Logos : Université de Yaoundé II à gauche, PSSFP à droite
    \node[fill=white, rounded corners=4pt, inner sep=2pt]
        at ([xshift=1.75cm, yshift=-1.8cm] current page.north west)
        % uy2_logo.jpeg est un bandeau (MINFI, MINESUP, UY2) : on n'en garde que le sceau de l'UY2
        {\includegraphics[height=2.6cm, trim=805 172 118 12, clip]{figures/logos/uy2_logo.jpeg}};
    \node[fill=white, rounded corners=4pt, inner sep=2pt]
        at ([xshift=-1.75cm, yshift=-1.8cm] current page.north east)
        {\includegraphics[height=2.6cm]{figures/logos/pssfp_logo.jpg}};
    % Année et version, posées dans le bandeau inférieur
    \node[align=center] at ([yshift=1.25cm] current page.south) {%
        {\color{white}\bfseries\large Année académique \anneescolaire}\\[0.3ex]
        {\color{CouleurAccent!70!white}\small
        Version 0.1 | Document pédagogique de l'UE \codeUE}};
\end{tikzpicture}

\vspace*{0.6cm}

% --- Titre principal (sur deux lignes) ---
\begin{center}
    {\color{CouleurPrimaire}\rule{0.85\textwidth}{1.2pt}}\\[1.2ex]
    {\fontsize{26}{32}\selectfont\bfseries\color{CouleurPrimaire}
        Contrôle et Audit des\\[0.4ex]
        Systèmes d'Information}\\[0.8ex]
    {\color{CouleurAccent}\rule{0.85\textwidth}{2pt}}\\[0.5ex]
    {\color{CouleurPrimaire}\rule{0.85\textwidth}{0.5pt}}
\end{center}

\vspace{0.8cm}

\begin{center}
    {\fontsize{13}{16}\selectfont\itshape\color{CouleurSecondaire}
        Évaluer le contrôle interne des finances publiques\\
        dans un environnement numérique}
\end{center}

\vspace{0.5cm}

% --- Cartouche sigle et code UE ---
\begin{center}
    \begin{tikzpicture}
        \node[draw=CouleurPrimaire, fill=CouleurPrimaire, rounded corners=6pt,
              inner xsep=14pt, inner ysep=6pt]{%
            \color{white}\bfseries\large \sigleCours{} - UE~\codeUE
        };
    \end{tikzpicture}
\end{center}

\vspace{1.0cm}

% --- Schéma de la discipline : l'évaluation du système de contrôle interne ---
\begin{center}
\begin{tikzpicture}[
    niv/.style = {
        draw=CouleurBordDef, fill=CouleurDef, rounded corners=3pt,
        minimum width=2.7cm, minimum height=0.75cm, align=center,
        font=\footnotesize\bfseries\color{CouleurPrimaire}, inner sep=2pt
    },
    arrow/.style = {-{Stealth[length=4pt]}, thick, color=CouleurSecondaire}
]
\foreach \t [count=\i] in {Mandat, Décrire, Questionner, Cheminer} {
    \node[niv] (n\i) at (3.3*\i, 0) {\t};
}
\foreach \t [count=\i from 5] in {Conception, Fonctionnement, Déficiences, Conclusion} {
    \node[niv] (n\i) at (3.3*\i - 13.2, -1.25) {\t};
}
\foreach \i [evaluate=\i as \j using int(\i+1)] in {1,2,3,5,6,7}
    \draw[arrow] (n\i) -- (n\j);
\draw[arrow] (n4.south) -- ++(0,-0.25) -| (n5.north);
\node[draw=CouleurPrimaire, fill=CouleurPrimaire, rounded corners=4pt,
      align=center, font=\footnotesize\bfseries\color{white}, inner xsep=8pt]
      (mes) at (8.25, -2.65)
      {\OUI{} / \NON{} / \NEK{} \quad$\longrightarrow$\quad
       le contrôle interne est-il efficace~?};
\draw[arrow] (n7.south) -- (n7.south |- mes.north);
\end{tikzpicture}
\end{center}

\vspace{0.6cm}

\begin{center}
    {\large\bfseries\color{CouleurPrimaire}
        \auteurcours}\\[0.5ex]
    {\normalsize\color{CouleurSecondaire}
       Enseignant Chercheur | \laboratoire{} | ENSPY}\\[0.2ex]
    {\small\color{CouleurBordInfo}
        Intervenant au PSSFP · \tutelle}\\[0.2ex]
    {\small\color{CouleurBordInfo}
        \href{mailto:maletyon@proton.me}{maletyon@proton.me}}
\end{center}

\end{titlepage}

    \cleardoublepage

    % ============================================================================
% Fichier  : front/01_avant_propos.tex
% Titre    : Avant-propos
% Auteur   : MINKA MI NGUIDJOÏ Thierry Emmanuel
% Version  : 0.1 -- novembre 2026
% ============================================================================
% 8 jours de 4 h (32 h au programme, 30 h effectives : les 2 dernières heures
% du J08 sont offertes). Seuls A0 et A1 (HAL) sont cités comme travaux de
% l'auteur.
% ============================================================================

\chapter*{Avant-propos}
\addcontentsline{toc}{chapter}{Avant-propos}
\markboth{Avant-propos}{Avant-propos}

\epigraph{« Le contrôle interne ne disparaît pas quand l'administration
s'informatise~; il change de place, et l'auditeur doit le suivre. »}%
{\textit{Principe du cours}}

Ce cours s'adresse aux auditeurs et aux contrôleurs des finances publiques
du Programme Supérieur de Spécialisation en Finances Publiques. Ils
connaissent déjà le contrôle sur pièces~: le mandat, la dépense, la
régularité, le contradictoire. Ils vont désormais exercer devant des
applications de budget, de dépense, de paie et de recette, où la pièce est
un enregistrement, la signature un droit d'accès, et le visa un contrôle
automatique.

Il repose sur une conviction~: l'informatique ne supprime pas le contrôle
interne, elle le déplace. La séparation entre l'ordonnateur et le comptable
devient une matrice de droits~; la vérification du service fait devient un
contrôle de saisie~; l'archive devient une sauvegarde. Un auditeur qui
continue de ne regarder que le dossier papier ne voit plus qu'une partie du
système. Le cœur du cours est donc le \textbf{système de contrôle interne}
et son \textbf{évaluation}~: le décrire, le tester, classer ses
déficiences, conclure sur son efficacité.

\bigskip

\textbf{Le cadre camerounais.} Au Cameroun, dans le secteur public, c'est
le Contrôle Supérieur de l'État qui a la compétence explicite du contrôle
des systèmes d'information~: son texte organique, le décret n°~2013/287
du 4~septembre~2013, la range parmi les cinq contrôles qu'il effectue.
L'ANTIC conduit les audits de sécurité des réseaux et systèmes
d'information au titre de la loi n°~2010/012. La Chambre des Comptes, juge
des comptes publics, peut s'y retrouver par extension, à travers les
comptes que ces systèmes produisent. Le Chapitre~\ref{chap:institutions}
précise ces compétences.

\bigskip

\textbf{Ce que ce cours n'est pas.} Il n'est pas un cours d'informatique,
ni une préparation à une certification, ni un cours de sécurité offensive.
Il ne demande aucun prérequis technique~: un tableur suffit pour les labs.

\bigskip

\textbf{Ce que ce cours est.} Un traité de quinze chapitres, en six
parties précédées d'un prologue, conçu pour trente heures effectives en
huit journées de quatre heures. La théorie y est lisible~; chaque notion
sert un geste de mission. La Partie~I pose le contrôle interne et la
manière de l'évaluer. La Partie~II présente les institutions et les
normes. La Partie~III conduit la mission, de l'ordre de mission au cycle
suivant~: principes, décomposition, instrument \OUI{}~/ \NON{}~/ \NEK{},
mesures de la conformité et de l'incertitude, rapport et suivi. La
Partie~IV évalue le contrôle interne domaine par domaine~: accès,
exploitation, applications, données, projets. La Partie~V réunit tout le
cours dans un cas intégrateur traité en salle. La Partie~VI ouvre sur la
recherche et propose des sujets de mémoire.

\bigskip

\textbf{Une méthode fondée sur une recherche.} La démarche de mission est
celle du cours d'audit des systèmes d'information de l'auteur à l'École
Nationale Supérieure Polytechnique de Yaoundé, elle-même instanciation,
pour le praticien, du programme de recherche de
l'\emph{auditologie}~\cite{minka2026auditologie}. L'instrument de réponse
\OUI{}~/ \NON{}~/ \NEK{} et ses mesures viennent de l'article
\cite{minka2026nek}. Ces deux textes sont des prépublications déposées sur
HAL, et le cours les cite comme telles.

\bigskip

\textbf{Trois entités fictives.} Le cours suit trois dossiers. Au
\emph{MINÉQUIR}, ministère imaginaire de l'Équipement rural, une facture
est payée deux fois par l'application de dépense. À la \emph{CAPNO},
caisse publique de prévoyance de Garoua, le fichier de paie contient des
agents qui n'existent pas. À la recette des finances de Bafoussam, le
projet de dématérialisation \emph{e-RECETTES} a été livré sans que
personne n'en ait vérifié les sauvegardes. Les noms, les dates et les
montants sont fictifs~; les mécanismes sont exacts.

\bigskip

\textbf{Comment travailler.} Entre deux journées, faites la fiche
d'auto-évaluation de chaque chapitre et les labs de NDAP SUKULU qui s'y
rattachent. Prenez surtout l'habitude que le cours veut transmettre~:
séparer ce que l'on a constaté de ce que l'on en pense, et écrire ce que
l'on ne sait pas.

\bigskip

\begin{flushright}
    \IfFileExists{../../Logos/signature_avant_propos.png}{\includegraphics[height=1.6cm]{../../Logos/signature_avant_propos.png}}{\textit{MINKA MI NGUIDJOÏ Thierry Emmanuel}}\\
    \textit{Yaoundé, novembre 2026}
\end{flushright}

    \cleardoublepage

    % ============================================================================
% Fichier  : front/02_mode_emploi.tex
% Titre    : Mode d'emploi du cours
% Auteur   : MINKA MI NGUIDJOÏ Thierry Emmanuel
% Version  : 0.1 -- novembre 2026
% ============================================================================
% 8 jours de 4 h ; les 2 dernières heures du J08 sont offertes (30 h
% effectives). Fiches et labs hors de ces heures.
% ============================================================================

\chapter*{Mode d'emploi du cours}
\addcontentsline{toc}{chapter}{Mode d'emploi du cours}
\markboth{Mode d'emploi}{Mode d'emploi}

\section*{Structure générale}

Le cours compte un prologue et six parties, soit quinze chapitres. Il
est donné en huit journées de quatre heures~; les deux dernières heures de
la huitième journée sont offertes aux auditeurs, ce qui fait trente heures
effectives. Les fiches d'auto-évaluation et les laboratoires de
60~minutes de NDAP SUKULU sont des travaux à faire hors de ces heures.

\begin{tableaucours}{L{1.6cm}L{4.4cm}L{4.4cm}L{1.8cm}}{%
    \tableauHeader{Partie} &
    \tableauHeader{Thème} &
    \tableauHeader{Compétences visées} &
    \tableauHeader{Niveau}
}
Prologue & Le paiement du 28 décembre & Voir ce que le dossier papier ne
montre pas & \niveaubase \\
I & Fondements et contrôle interne & Situer l'audit d'un SI public,
décrire et évaluer un système de contrôle interne & \niveaubase{} à
\niveaumoyen \\
II & Cadre institutionnel et normatif & Savoir qui contrôle quoi, choisir
ses normes et ses critères & \niveaubase \\
III & Démarche de mission & Préparer et décomposer, coter \OUI{}~/
\NON{}~/ \NEK{}, mesurer $\mu$, $U$, $H$, conclure, rapporter, suivre les
cycles & \niveaumoyen \\
IV & Domaines d'audit & Évaluer les accès, l'exploitation, les
applications, les données et les projets & \niveaumoyen{} à
\niveauexpert \\
V & Synthèse & Conduire une mission complète en équipe & \niveauexpert \\
VI & Ouverture sur la recherche & Situer l'auditologie, formuler un sujet
de mémoire falsifiable & \niveauexpert \\
\end{tableaucours}

\section*{Les huit journées}

\begin{tableaucours}{L{1.4cm}L{8.6cm}L{2.2cm}}{%
    \tableauHeader{Jour} &
    \tableauHeader{Contenu} &
    \tableauHeader{Partie}
}
J01 & Prologue~; contrôler et auditer un SI public
(Ch.~\ref{chap:audit-public}) & Prologue, I \\
J02 & Le système de contrôle interne (Ch.~\ref{chap:sci})~; évaluer le
système de contrôle interne (Ch.~\ref{chap:eval-sci}) & I \\
J03 & Qui contrôle les SI publics (Ch.~\ref{chap:institutions})~; normes
et référentiels (Ch.~\ref{chap:normes}) & II \\
J04 & La mission~: principes, préparation, décomposition
(Ch.~\ref{chap:preparer})~; collecter la preuve~: \OUI{}, \NON{}, \NEK{}
(Ch.~\ref{chap:preuve}) & III \\
J05 & Mesurer, conclure, rapporter, suivre (Ch.~\ref{chap:constater})~;
accès et habilitations (Ch.~\ref{chap:acces}) & III, IV \\
J06 & Changements, exploitation et continuité
(Ch.~\ref{chap:exploitation})~; contrôles d'application
(Ch.~\ref{chap:applications}) & IV \\
J07 & Analyse de données et fraude (Ch.~\ref{chap:donnees-fraude})~;
gouvernance, projets et prestataires (Ch.~\ref{chap:gouvernance}) & IV \\
J08 & Cas intégrateur en salle (Ch.~\ref{chap:cas-integrateur}) et
ouverture sur la recherche (Ch.~\ref{chap:recherche}), deux heures~; les
deux heures suivantes sont offertes & V, VI \\
\end{tableaucours}

\section*{Les encadrés}

\begin{boxdef}
\textbf{Définition}. Introduit un terme avec sa définition~; les termes
définis sont indexés et repris dans le glossaire (Annexe~\ref{ann:glossaire}).
\end{boxdef}

\begin{boxscenario}{Titre du scénario}
\textbf{Scénario}. Un épisode d'un des trois dossiers du cours~: MINÉQUIR,
CAPNO, e-RECETTES.
\end{boxscenario}

\begin{boxobjet}{Code · Intitulé}
\textbf{Point de contrôle}. Présente un contrôle sous une forme
constante~: objectif, risque couvert, test de conception, test de
fonctionnement, preuve attendue.
\end{boxobjet}

\begin{boxprocedure}
\textbf{Papier de travail}. Un modèle de feuille de mission, prêt à
remplir~: ce que l'auditeur fait et ce que le dossier doit contenir.
\end{boxprocedure}

\begin{boxlecture}{Titre de la lecture}
\textbf{Lecture auditologique}. Des réponses \OUI{}~/ \NON{}~/ \NEK{},
avec la constatation, l'analyse et son niveau de confiance, et, selon le
cas, le décompte des contrôles cotés ou les trois mesures du domaine.

\medskip
\hfill\conformite{3}{8}{0}\quad\lecture{0{,}55}{0{,}30}{1{,}41}
\end{boxlecture}

\begin{boxmetier}{Vérificateur}
\textbf{Application métier}. Traduit le chapitre en gestes concrets pour
l'un des métiers visés~: vérificateur du Contrôle Supérieur de l'État,
juge des comptes, inspecteur, contrôleur financier, auditeur interne.
\end{boxmetier}

\begin{boxoutils}{Titre de l'outil}
\textbf{Outils}. Une formule de tableur, une requête ou une commande, à
refaire sur les données du cas.
\end{boxoutils}

\begin{boxjuris}
Cite un texte de loi ou de règlement, lu dans sa version disponible, puis
en tire la portée pour l'audit.
\end{boxjuris}

\begin{boxalert}
\textbf{Violation typique} et \textbf{règle}~: une erreur fréquente, et ce
qu'il faut faire à la place.
\end{boxalert}

\begin{boxresume}
\textbf{À retenir}. Les points essentiels du chapitre.
\end{boxresume}

\begin{boxexo}
\textbf{Exercice}. Gradués en trois niveaux~: \niveaubase\ fondamental,
\niveaumoyen\ intermédiaire, \niveauexpert\ expert. Les corrigés figurent
dans le livret de correction.
\end{boxexo}

\section*{Conventions}

\begin{itemize}
    \item Chaque contrôle testé reçoit une réponse~: \OUI{} (le contrôle
          existe et fonctionne), \NON{} (il manque ou ne fonctionne pas),
          \NEK{} (connaissance insuffisante pour conclure). Un contrôle sans
          objet est écarté, raison écrite.
    \item L'absence d'une trace que le contrôle devait produire vaut \NON{},
          et non \NEK{}.
    \item $\mu$ est le taux de conformité, $U$ le taux d'incertitude, $H$
          l'entropie en bits (au plus $\log_2 3 \approx 1{,}585$), selon
          l'article A1~\cite{minka2026nek}~; la règle de décision examine $U$
          d'abord, $\mu$ ensuite, $H$ pour contrôler. Les seuils du quadrant
          ($\mu = 0{,}5$~; $U = 0{,}20$) sont des conventions du cours.
    \item Les axiomes de l'auditologie portent la numérotation de l'article
          A0 (Chapitre~\ref{chap:recherche}).
    \item Chaque analyse indique un niveau de confiance~: faible, moyenne ou
          élevée.
    \item Les \termecle{termes clés} apparaissent en gras~; les normes
          s'écrivent \norme{ISSAI~100}, les textes camerounais
          \loicam{décret n°~2013/287}.
\end{itemize}

\section*{Matériel d'accompagnement}

\begin{itemize}
    \item \textbf{Annexes}~: glossaire, outils, papiers de travail types,
          textes, bibliographie, contacts (Annexes~\ref{ann:glossaire}
          à~\ref{ann:contacts}).
    \item \textbf{Livret de correction}~: corrigés des exercices et des
          fiches, réservé au formateur.
    \item \textbf{NDAP SUKULU}~: fiches d'auto-évaluation et laboratoires,
          présentés dans la section qui suit.
\end{itemize}

    \cleardoublepage

    % ==============================================================================
% Fichier : front/03_ndap_sukulu.tex
% Titre   : NDAP SUKULU, l'environnement compagnon du cours
% Auteur  : MINKA MI NGUIDJOÏ Thierry Emmanuel
% Version : 1.0 -- novembre 2026
% ==============================================================================
% Texte commun aux cours de l'auteur ; seul le tableau des labs est propre au
% cours (fichier front/03b_labs.tex).
% ==============================================================================

\chapter*{NDAP SUKULU, la salle de classe compagnon}
\addcontentsline{toc}{chapter}{NDAP SUKULU, la salle de classe compagnon}
\markboth{NDAP SUKULU}{NDAP SUKULU}

\begin{center}
    \includegraphics[width=0.62\textwidth]{figures/logos/logo_ndap_sukulu.png}
\end{center}

\section*{Un nom venu du pays bassa}

\begin{quote}\itshape
Le soir tombe sur Éséka. Un enfant revient des champs avec son grand-père et
lui demande pourquoi il doit marcher si loin, chaque matin, pour aller à
l'école. Le vieil homme s'arrête au bord du chemin et répond, en bassa~:
« \textup{Ndap sukulu}, la salle de classe, ce ne sont pas les murs. Ce sont
les murs qui viennent après. La salle de classe, c'est l'endroit où quelqu'un
t'attend pour t'apprendre ce qu'il sait. » L'enfant insiste~: « Et qui
m'attend~? » Le grand-père sourit~: « \textup{Malet yon}. Ton enseignant. »
\end{quote}

En bassa, la langue paternelle de l'auteur, \textbf{\textit{ndap sukulu}}
désigne la salle de classe et \textbf{\textit{malet yon}} signifie « ton
enseignant ». L'environnement compagnon de ce cours porte ce nom parce qu'il
prolonge la salle de classe au-delà des heures de cours~: il attend
l'étudiant le soir, le week-end, à distance, avec des questions, des pièces à
examiner et une correction immédiate. Il est édité par \textbf{MaletYon},
« ton enseignant », et mis à la disposition des étudiants et des praticiens
qui souhaitent se perfectionner (\href{mailto:maletyon@proton.me}{maletyon@proton.me}).

\section*{Ce que l'on y trouve}

\begin{itemize}
    \item \textbf{Une fiche d'auto-évaluation par chapitre.} Chaque fiche
          puise vingt questions dans une banque de quarante (QCM, vrai ou
          faux, textes à trous). Elle dure trente minutes~; la copie validée
          est verrouillée et notée aussitôt. Le seuil de réussite est de
          60~\%.
    \item \textbf{Des laboratoires pratiques.} Chaque lab place l'étudiant
          devant les pièces d'une affaire du fil rouge (tableurs, journaux,
          contrats, procès-verbaux) dans un bureau d'analyste virtuel,
          ouvert dans le navigateur. Il dure soixante minutes, se termine
          par des points de contrôle notés automatiquement et renvoie aux
          sections du livre qu'il mobilise.
    \item \textbf{Un suivi pour l'enseignant}~: rapport de classe, notes par
          fiche et par lab, journal d'activité.
\end{itemize}

Les fiches et les labs sont des travaux à faire hors des heures de cours.
Ils ne remplacent ni le livre ni les exercices de fin de chapitre~: ils en
vérifient la maîtrise.

\section*{Comment y travailler}

\begin{enumerate}
    \item Se connecter au portail NDAP SUKULU avec ses identifiants et le
          code de cours \textbf{\codeUE}.
    \item Pour un lab, commencer par vérifier l'empreinte SHA-256 des pièces
          (\texttt{sha256sum -c MANIFESTE\_SHA256.txt}), puis travailler sur
          la copie du dossier \texttt{travail}~: les pièces originales sont
          en lecture seule et l'activité est journalisée, comme dans une
          vraie mission.
    \item Lire le rappel de cours affiché par le lab et relire les sections
          citées avant de répondre.
\end{enumerate}

\begin{boxinfo}
\textbf{Si le portail n'est pas encore ouvert.} En attendant le déploiement
de NDAP SUKULU, les fiches de questions sont fournies au format texte dans le
dossier \texttt{fiches} du cours, prêtes à être distribuées. Les corrections
des fiches figurent dans le livret de correction, à la suite des corrigés des
exercices de chaque chapitre.
\end{boxinfo}

\section*{Les laboratoires du cours}

% Tableau des labs du cours CASI (CAS3401) dans NDAP SUKULU (propre au cours).
Le cours compte dix laboratoires, disponibles sous le code \codeUE. Ils se
font à la maison, hors des heures de cours.

\begin{tableaucours}{L{2.6cm}L{6.0cm}L{1.6cm}L{2.0cm}}{%
    \tableauHeader{Lab} &
    \tableauHeader{Sujet} &
    \tableauHeader{Chapitre} &
    \tableauHeader{Dossier}
}
CAS3401-LAB-01 & Prise de connaissance et risques de SIDÉPENSE &
\ref{chap:preparer} & MINÉQUIR \\
CAS3401-LAB-02 & Évaluer le contrôle interne~: questionnaire, cheminement,
déficiences & \ref{chap:eval-sci} & MINÉQUIR \\
CAS3401-LAB-03 & Habilitations~: séparation des tâches et comptes des
agents partis & \ref{chap:acces} & MINÉQUIR \\
CAS3401-LAB-04 & Chronologie du double paiement dans les journaux &
\ref{chap:preuve} & MINÉQUIR \\
CAS3401-LAB-05 & Intégrité d'une extraction & \ref{chap:preuve} &
CAPNO \\
CAS3401-LAB-06 & Agents fictifs et doublons dans la paie &
\ref{chap:donnees-fraude} & CAPNO \\
CAS3401-LAB-07 & Loi de Benford et fractionnement des marchés &
\ref{chap:donnees-fraude} & MINÉQUIR \\
CAS3401-LAB-08 & Changements et sauvegardes & \ref{chap:exploitation} &
e-RECETTES \\
CAS3401-LAB-09 & Tester un contrôle d'application &
\ref{chap:applications} & MINÉQUIR \\
CAS3401-LAB-10 & Rédiger un constat et classer la déficience &
\ref{chap:constater} & les trois \\
\end{tableaucours}
\vspace{-0.8em}
\begin{center}{\small\itshape Les laboratoires du cours dans NDAP SUKULU}\end{center}


    \cleardoublepage

    \tableofcontents
    \cleardoublepage

    \listoffigures
    \listoftables
    \cleardoublepage

\mainmatter

    % ======================================================================
    % PROLOGUE OPÉRATIONNEL (J01)
    % « Le dossier était en ordre ; le système, non. »
    % Standards de référence : aucun ; un cas qui échoue.
    % Niveau : introductif
    % ======================================================================
    \part*{Prologue Opérationnel}
    \addcontentsline{toc}{part}{Prologue Opérationnel}

        \setcounter{chapter}{-1}
        \IfFileExists{partie0/P00_paiement_28_decembre.tex}{% ============================================================================
% Fichier  : partie0/P00_paiement_28_decembre.tex
% Titre    : Prologue -- Le Paiement du 28 Décembre
% Auteur   : MINKA MI NGUIDJOÏ Thierry Emmanuel
% Version  : 0.1 -- novembre 2026
% ============================================================================
% Faits du dossier MINÉQUIR à garder identiques dans tout le cours :
%  - marché n° 0412/M/MINÉQUIR/2025, BATIRURAL SARL (Mbalmayo), pistes rurales ;
%    décompte n° 3, facture F-2025-117, 187 450 000 FCFA ;
%  - mandat 2025-08817 (22/12/2025, circuit normal), payé le 24/12 sur le
%    compte n° 1 (banque A) ;
%  - 27/12/2025 21 h 52 : RIB modifié par le compte mbarga.j (compte n° 2,
%    banque B, au nom de « BATIRURAL SERVICES ») ;
%  - 28/12/2025 23 h 14 : mandat 2025-09342, facture « F-2025-117 bis »,
%    liquidé et ordonnancé par mbarga.j ; payé le 30/12 ;
%  - MBARGA, ex-chef du service de la liquidation, muté à Ebolowa le
%    14/08/2025 ; compte jamais désactivé ;
%  - janvier 2026 : inspection sur pièces, 40 dossiers papier de décembre,
%    « régularité satisfaisante » ;
%  - 16/03/2026 : ordre de mission du CONSUPE (vérificatrice NGO BIYONG) ;
%    extraction de 9 611 mandats 2025 ; 3 paires de doublons, dont 1 seule
%    avec RIB différent ;
%  - lecture des 8 contrôles attendus : 2 OUI, 5 NON, 1 NEK.
% Personnages : EBOUMBOU (DAF), ATANGANA (chef de la cellule informatique),
% FOTSO (chef du poste comptable), NGO BIYONG (vérificatrice du CONSUPE).
% ============================================================================

\chapter[Prologue : Le Paiement du 28 Décembre]{Prologue : Le Paiement du 28 Décembre}
\label{chap:prologue}
\chaptermark{Prologue}

\epigraph{%
    « Le dossier était en ordre~; le système, non. »%
}{Note de synthèse d'une vérificatrice}

\niveauChapitre{Introductif}{%
    Aucun prérequis informatique. Il suffit de connaître le circuit de la
    dépense publique~: engagement, liquidation, ordonnancement, paiement.
    Ce prologue se lit avant toute définition~: il part d'un contrôle qui
    n'a rien vu.%
}

\begin{boxobjectifs}
\begin{itemize}
    \item Reconstituer, à partir des pièces et des traces informatiques, la
          chronologie d'un double paiement (dossier MINÉQUIR).
    \item Expliquer pourquoi un contrôle sur pièces, bien conduit, peut ne
          pas voir une irrégularité qui se trouve dans le système.
    \item Distinguer une constatation, appuyée sur une preuve conservée, de
          l'analyse que l'on en tire.
    \item Identifier les contrôles internes qui auraient dû empêcher ou
          détecter le double paiement, et les coter \OUI{}, \NON{} ou
          \NEK{}.
    \item Citer le texte qui donne au Contrôle Supérieur de l'État la
          compétence du contrôle des systèmes d'information.
    \item Énoncer les trois réflexes de l'auditeur public face à un système
          d'information~: le mandat, la preuve, le contradictoire.
\end{itemize}
\end{boxobjectifs}

\bigskip

% ============================================================================
\section*{Avant de commencer : un contrôle qui n'a rien vu}
% ============================================================================

Yaoundé, lundi 16~mars~2026. Mme~NGO~BIYONG, vérificatrice au Contrôle
Supérieur de l'État, reçoit son ordre de mission~: vérifier l'exécution du
budget 2025 de la Direction des Affaires Financières du ministère de
l'Équipement rural, le MINÉQUIR\footnote{Le MINÉQUIR, ses agents, ses
fournisseurs, les dates et les montants de ce dossier sont fictifs. Les
mécanismes du circuit de la dépense, les contrôles et les textes cités sont
exacts.}. Dans le dossier permanent, une pièce attire son attention~: le
rapport d'une inspection conduite en janvier sur les dépenses de décembre.
Il examine quarante dossiers de paiement et conclut à une «~régularité
satisfaisante~».

Trois semaines plus tard, la même vérificatrice établit qu'une facture a
été payée deux fois en décembre, la seconde fois sur un compte bancaire qui
n'appartient pas au fournisseur. Le montant en jeu est de
\textbf{187\,450\,000~francs CFA}.
L'inspection de janvier ne s'était pas trompée sur les dossiers qu'elle a
examinés~: ils étaient complets et réguliers. Elle n'a simplement jamais vu
le second paiement, parce qu'il n'avait pas de dossier.

\medskip

Ce prologue reconstitue ce dossier. Il montre ce qu'un contrôle sur pièces
peut voir et ce qu'il ne peut pas voir, ce que la trace informatique
révèle, et quels contrôles internes ont manqué. Il en tire les trois
réflexes qui structurent tout le cours. Les notions y sont nommées~; elles
sont définies avec précision dans les chapitres suivants.

\begin{boxalert}
\textbf{Violation typique.} Tirer l'échantillon d'un contrôle dans les
classeurs de dossiers papier, alors que les paiements sont produits par une
application~: ce qui n'a pas de dossier ne peut pas être tiré.

\smallskip
\textbf{Règle~:} la population d'un contrôle est l'ensemble des opérations
enregistrées dans le système, pas l'ensemble des dossiers rangés.
\end{boxalert}

\separateur

% ============================================================================
\section{Ce qui s'est passé}
\label{sec:pro-chronologie}
% ============================================================================

\subsection{Le marché et le circuit}

Le MINÉQUIR a passé en 2025 le marché n°~0412/M/MINÉQUIR/2025 avec la
société BATIRURAL~SARL, de Mbalmayo, pour l'aménagement de pistes rurales.
Le décompte n°~3 de ce marché, facturé sous le numéro F-2025-117, s'élève
à 187\,450\,000~FCFA.

Depuis 2024, la dépense du ministère passe par une application, SIDÉPENSE.
Chaque étape du circuit y est un écran et un profil d'utilisateur~:
l'engagement réserve le crédit, la liquidation constate le service fait et
arrête le montant, l'ordonnancement émet le mandat, puis le poste comptable
reçoit le mandat par voie électronique et paie. Les pièces justificatives
sont scannées et jointes au mandat~; le dossier papier est classé à la
Direction des Affaires Financières.

\begin{tableaucours}{L{4.0cm}L{8.2cm}}{%
    \tableauHeader{Élément} &
    \tableauHeader{Description}
}
Entité & Direction des Affaires Financières (DAF) du MINÉQUIR, Yaoundé \\
Application & SIDÉPENSE~: engagement, liquidation, ordonnancement,
transmission des mandats au poste comptable, fichier des fournisseurs \\
Volume 2025 & 9\,611 mandats \\
Acteurs & M.~EBOUMBOU, directeur des affaires financières~; M.~ATANGANA,
chef de la cellule informatique~; M.~FOTSO, chef du poste comptable \\
Contrôle & Inspection sur pièces (janvier 2026)~; mission du CONSUPE
(mars 2026), conduite par Mme~NGO~BIYONG \\
\end{tableaucours}
\vspace{-0.8em}
\begin{center}{\small\itshape Tableau~P.1. Le dossier MINÉQUIR}\end{center}

\subsection{Reconstituer avant de juger}

Une vérificatrice ne commence pas par chercher un responsable~: elle
reconstitue la suite des événements, et rattache chaque événement à une
pièce. Ce qu'on lui raconte sans pièce reste une déclaration, à vérifier.
Le Tableau~P.2 donne la chronologie établie par Mme~NGO~BIYONG.

\begin{tableaucours}{L{2.6cm}L{6.4cm}L{3.2cm}}{%
    \tableauHeader{Date} &
    \tableauHeader{Événement} &
    \tableauHeader{Pièce}
}
14 août 2025 & M.~MBARGA, chef du service de la liquidation, est muté à
Ebolowa & Décision de mutation \\
22 déc. 2025 & Mandat 2025-08817 (décompte n°~3) émis par le circuit
normal & Mandat, dossier papier complet \\
24 déc. 2025 & Paiement sur le compte n°~1 de BATIRURAL (banque~A) &
Relevé du poste comptable \\
27 déc. 2025, 21~h~52 & Le RIB de BATIRURAL est remplacé, dans le fichier
des fournisseurs, par un compte n°~2 (banque~B) au nom de «~BATIRURAL
SERVICES~» & Journal de SIDÉPENSE \\
28 déc. 2025, 23~h~14 & Mandat 2025-09342, facture «~F-2025-117 bis~»,
même montant, liquidé et ordonnancé par le compte \texttt{mbarga.j} &
Journal de SIDÉPENSE, mandat électronique \\
30 déc. 2025 & Paiement du second mandat sur le compte n°~2 & Relevé du
poste comptable \\
Janv. 2026 & Inspection de 40 dossiers papier de décembre~: «~régularité
satisfaisante~» & Rapport d'inspection \\
16 mars 2026 & Ordre de mission du CONSUPE & Ordre de mission \\
23 mars 2026 & Extraction des 9\,611 mandats de 2025 et du journal des
opérations & Fichiers extraits, empreintes \\
\end{tableaucours}
\vspace{-0.8em}
\begin{center}{\small\itshape Tableau~P.2. Chronologie du double paiement
(reconstitution de la vérificatrice)}\end{center}

Ce tableau ne contient que des faits. Il ne dit pas encore qui a agi, ni
pourquoi~: le compte \texttt{mbarga.j} a servi, mais rien n'établit à ce
stade que M.~MBARGA l'ait utilisé lui-même. Séparer ce qui est établi de ce
que l'on en déduit est la première discipline de l'auditeur.

\begin{boxdef}
\termecle{Constatation}\index{constatation} (\textit{finding})~: fait
vérifiable, établi par l'auditeur et appuyé sur une preuve conservée. Un
tiers compétent qui refait la même démarche sur les mêmes pièces aboutit au
même résultat.

\medskip
\termecle{Analyse}\index{analyse}~: interprétation que l'auditeur tire
d'une ou plusieurs constatations, au regard d'un critère, avec un niveau de
confiance déclaré (faible, moyen ou élevé).
\end{boxdef}

\separateur

% ============================================================================
\section{Ce que le contrôle sur pièces a vu}
\label{sec:pro-pieces}
% ============================================================================

L'inspection de janvier a fait correctement ce qu'elle a fait. Elle a tiré
quarante dossiers dans les classeurs de décembre, vérifié pour chacun
l'engagement, le service fait, la facture, le visa et l'imputation. Les
quarante dossiers étaient complets. Le mandat 2025-08817 en faisait partie.

Son erreur n'est pas dans l'examen, elle est dans la \textbf{population}.
L'échantillon a été tiré dans les classeurs, c'est-à-dire parmi les
mandats qui avaient un dossier papier. Le mandat 2025-09342 n'en avait pas~:
il a été créé directement dans l'application, avec un scan de la facture du
premier mandat, renommée «~bis~». Aucun tirage dans les classeurs ne
pouvait le faire sortir.

\begin{boxscenario}{Dossier MINÉQUIR. Deux populations}
\textbf{Constatation.} Les classeurs de décembre contiennent 1\,184
dossiers~; l'application enregistre 1\,191 mandats émis en décembre. Sept
mandats n'ont pas de dossier papier, dont le mandat 2025-09342.

\medskip
\textbf{Analyse} (confiance élevée). L'inspection a conclu sur la
population des dossiers rangés, pas sur celle des mandats émis. Sa
conclusion vaut pour les 1\,184 dossiers~; elle ne dit rien des sept
autres. Un rapprochement entre les deux populations, fait avant le tirage,
aurait suffi à les faire apparaître.
\end{boxscenario}

Cette différence entre deux populations est la première leçon du
cours~: dès que l'administration s'informatise, le dossier papier n'est
plus le système, il n'en est qu'une sortie.

\separateur

% ============================================================================
\section{Ce que la trace informatique révèle}
\label{sec:pro-trace}
% ============================================================================

\subsection{L'extraction}

Mme~NGO~BIYONG ne part pas des classeurs, mais de l'application. Elle
demande par écrit à M.~ATANGANA une extraction de tous les mandats de 2025
et du journal des opérations, assiste à l'extraction, calcule l'empreinte
des fichiers et la fait contresigner. Elle vérifie ensuite que le total des
mandats extraits est égal au total des paiements du compte administratif.
Le Chapitre~\ref{chap:preuve} détaille ces précautions~: sans elles, une
extraction n'est qu'un fichier remis par l'audité.

\subsection{Le test des doublons}

Sur les 9\,611 mandats, elle recherche ceux qui portent le même
fournisseur et le même montant. Le test, fait dans un tableur, renvoie
trois paires. Deux s'expliquent par des pièces (deux livraisons identiques
d'un même fournisseur). La troisième est celle de BATIRURAL~: même
fournisseur, même montant, même numéro de facture au suffixe près, et deux
comptes bancaires différents.

\begin{boxoutils}{Recherche des doublons dans un tableur}
Dans la feuille des mandats (colonnes~: B fournisseur, D montant), une
colonne de clé et une colonne de comptage suffisent~:
\begin{lstlisting}
F2 : =B2&"|"&D2
G2 : =NB.SI($F$2:$F$9612;F2)
\end{lstlisting}
Un filtre sur \texttt{G~>~1} affiche les mandats qui partagent le même
fournisseur et le même montant. Le test ne prouve rien à lui seul~: il
désigne les mandats à examiner, pièce par pièce.
\end{boxoutils}

\subsection{Le journal}

Le journal de SIDÉPENSE enregistre, pour chaque opération sensible,
l'utilisateur, la date, l'heure et l'action. Il permet de reconstituer
les deux opérations du Tableau~P.3.

\begin{lstlisting}[basicstyle=\ttfamily\scriptsize]
2025-12-27 21:52:07 mbarga.j FOURNISSEUR MODIF_RIB   BATIRURAL  cpte1->cpte2
2025-12-28 23:09:41 mbarga.j MANDAT      LIQUIDATION 2025-09342 187450000
2025-12-28 23:14:18 mbarga.j MANDAT      ORDONNANCE  2025-09342 187450000
2025-12-30 10:02:55 pc.fotso MANDAT      PAIEMENT    2025-09342 cpte2
\end{lstlisting}
\vspace{-0.6em}
\begin{center}{\small\itshape Tableau~P.3. Extrait du journal de SIDÉPENSE
(fichier extrait le 23~mars~2026, empreinte au dossier)}\end{center}

\begin{boxscenario}{Dossier MINÉQUIR. Ce que dit le journal}
\textbf{Constatations.}
\begin{enumerate}[(a)]
    \item Le compte \texttt{mbarga.j} était actif le 28~décembre~2025,
          quatre mois après la mutation de son titulaire.
    \item Le même compte a modifié un RIB, liquidé et ordonnancé le même
          mandat~; aucune deuxième personne n'est intervenue.
    \item Les deux opérations ont eu lieu un samedi soir et un dimanche
          soir, hors des heures de service.
    \item SIDÉPENSE a accepté un second mandat de même fournisseur, même
          montant et même numéro de facture au suffixe «~bis~» près.
\end{enumerate}

\medskip
\textbf{Analyse} (confiance moyenne). Le double paiement n'est pas une
erreur de saisie~: la modification du RIB la veille, sur le seul compte
du fournisseur concerné, et l'heure des opérations désignent une
manœuvre. L'identité de la personne qui a utilisé le compte n'est pas
établie~: un compte resté actif peut être utilisé par quelqu'un qui en
connaît le mot de passe.
\end{boxscenario}

\separateur

% ============================================================================
\section{Les contrôles internes qui ont manqué}
\label{sec:pro-controles}
% ============================================================================

La question de l'auditeur n'est pas seulement «~qui a fait cela~?~», qui
relève d'autres procédures, mais «~qu'est-ce qui aurait dû l'empêcher, ou
le détecter, et pourquoi cela n'a-t-il pas fonctionné~?~». C'est la
question du \termecle{contrôle interne}\index{contrôle interne}~: l'ensemble
des dispositifs qui donnent une assurance raisonnable que les opérations
sont régulières, les informations fiables et les biens protégés. Le
Chapitre~\ref{chap:sci} le définit selon le COSO, et le
Chapitre~\ref{chap:eval-sci} apprend à l'évaluer.

Mme~NGO~BIYONG dresse la liste des contrôles qu'elle s'attendait à trouver
sur ce circuit, et cote chacun. Elle répond \OUI{} si le contrôle existe et
a fonctionné, \NON{} s'il manque ou n'a pas fonctionné, \NEK{}
(\textit{Not Enough Knowledge}) si elle n'a pas encore de quoi conclure.

\begin{tableaucours}{L{0.7cm}L{6.6cm}L{1.2cm}L{3.6cm}}{%
    \tableauHeader{N°} &
    \tableauHeader{Contrôle attendu} &
    \tableauHeader{Cote} &
    \tableauHeader{Preuve}
}
1 & Les comptes des agents mutés sont désactivés à la date de la mutation
& \NON{} & Liste des comptes actifs, décision de mutation \\
2 & Un même profil ne peut pas liquider et ordonnancer & \NON{} & Matrice
des profils de SIDÉPENSE \\
3 & Toute modification de RIB exige une seconde validation & \NON{} &
Journal~: une seule intervention \\
4 & L'application bloque un mandat de même fournisseur, même montant et
même facture & \NON{} & Mandat 2025-09342 accepté \\
5 & Le poste comptable vérifie les pièces jointes avant paiement & \NEK{} &
Pièces du poste comptable demandées, non reçues \\
6 & Les opérations sensibles sont journalisées & \OUI{} & Journal
extrait \\
7 & Les paiements de fin de gestion font l'objet d'une revue & \NON{} &
Aucune revue produite \\
8 & Les journaux sont conservés et protégés & \OUI{} & Journal intact
depuis 2024 \\
\end{tableaucours}
\vspace{-0.8em}
\begin{center}{\small\itshape Tableau~P.4. Les contrôles attendus sur le
circuit de la dépense}\end{center}

\begin{boxlecture}{Circuit de la dépense, MINÉQUIR}
\textbf{Constatation.} Sur huit contrôles attendus, deux existent et
fonctionnent, cinq manquent ou n'ont pas fonctionné, un ne peut pas encore
être coté.

\medskip
\textbf{Analyse} (confiance élevée). Les deux contrôles qui fonctionnent
sont des contrôles de \emph{détection} a posteriori~: ils ont permis de
reconstituer les faits, pas de les empêcher. Tous les contrôles de
\emph{prévention} ont manqué. La cote \NEK{} du contrôle~5 n'est pas une
faiblesse de la mission~: elle dit honnêtement ce qui reste à établir, et
désigne la prochaine pièce à obtenir.

\medskip
\hfill\conformite{2}{8}{1}
\end{boxlecture}

\begin{boxinfo}
\textbf{Pourquoi \NEK{} et pas \NON{}~?} Si le poste comptable devait
produire une trace de sa vérification et ne l'a pas produite, la réponse
est \NON{}~: l'absence d'une trace exigée est une constatation. Ici, la
vérificatrice n'a pas encore reçu les pièces qu'elle a demandées~: elle ne
sait pas. Elle l'écrit, et relance.
\end{boxinfo}

\separateur

% ============================================================================
\section{Qui peut contrôler ce système}
\label{sec:pro-mandat}
% ============================================================================

Mme~NGO~BIYONG n'a pas eu à demander la permission d'auditer SIDÉPENSE.
Le texte organique du Contrôle Supérieur de l'État le lui permet
expressément.

\begin{boxjuris}
\textbf{\loicam{Décret n°~2013/287 du 4~septembre~2013} portant
organisation des services du Contrôle supérieur de l'État, article~2,
alinéa~3}~\cite{decret2013287consupe}~: \textit{« Dans l'exercice de leurs
attributions, les services du Contrôle supérieur de l'Etat effectuent~:
un contrôle de conformité et de régularité~; un contrôle financier~; un
contrôle de performance~; un contrôle de l'environnement~; un contrôle des
systèmes d'information. »}

\medskip
\textbf{Article~2, alinéa~4}~: \textit{« Les contrôles sont effectués par
les vérificateurs du Contrôle supérieur de l'Etat sur pièces et sur place.
A cet effet, et nonobstant les contrôles effectués par les autres organes
de contrôle, les responsables des structures auditées sont tenus de
présenter les pièces relatives à la période auditée. »}

\medskip
\textbf{Portée pour l'audit.} Le contrôle des systèmes d'information est
une compétence écrite du CONSUPE, au même titre que le contrôle de
régularité. L'obligation de présenter les pièces couvre les pièces
électroniques~: une extraction, un journal, une matrice des profils sont
des pièces relatives à la période auditée.
\end{boxjuris}

Deux autres institutions regardent les systèmes d'information publics,
chacune sous un angle différent. L'\sigle{ANTIC} conduit les audits de
sécurité des réseaux et des systèmes d'information prévus par la
\loicam{loi n°~2010/012}. La Chambre des Comptes de la Cour suprême juge
les comptes des comptables publics~; elle rencontre les systèmes
d'information par extension, parce que ces comptes en sont issus. Le
Chapitre~\ref{chap:institutions} précise qui fait quoi, et comment ces
interventions s'articulent.

\separateur

% ============================================================================
\section{Trois réflexes}
\label{sec:pro-reflexes}
% ============================================================================

Le dossier MINÉQUIR fait apparaître les trois réflexes que le cours
installe, chapitre après chapitre.

\begin{enumerate}
    \item \textbf{Le mandat.} Savoir au nom de quoi l'on contrôle, sur quel
          périmètre et selon quels critères. C'est l'ordre de mission qui
          a permis d'exiger l'extraction~; c'est le critère («~un mandat
          ne peut pas être payé deux fois~») qui a permis de coter
          (Chapitres~\ref{chap:institutions} et~\ref{chap:preparer}).
    \item \textbf{La preuve.} Partir du système, pas du classeur~;
          obtenir la donnée soi-même ou y assister, en garantir
          l'intégrité, et conserver tout ce qui fonde une cote
          (Chapitre~\ref{chap:preuve}).
    \item \textbf{Le contradictoire.} Soumettre les constatations à
          l'entité avant de conclure. Le directeur des affaires
          financières et le chef de la cellule informatique recevront le
          Tableau~P.4 et pourront le contester pièce par pièce
          (Chapitre~\ref{chap:constater}).
\end{enumerate}

\begin{boxprocedure}
\textbf{Le contenu minimal d'un papier de travail}, pour chaque contrôle
coté~:
\begin{itemize}
    \item le contrôle attendu et le critère qui le fonde~;
    \item le test effectué, la date et la personne qui l'a fait~;
    \item la pièce obtenue (nom du fichier, empreinte, source)~;
    \item la cote (\OUI{}, \NON{} ou \NEK{}) et la constatation qui la
          justifie~;
    \item pour un \NEK{}, la pièce manquante et la relance.
\end{itemize}
\end{boxprocedure}

\begin{boxmetier}{Vérificateur}
Avant tout tirage d'échantillon, rapprochez le nombre et le montant des
opérations du système avec ceux de la source que vous allez échantillonner
(classeurs, bordereaux). Un écart, même de quelques unités, désigne les
opérations qu'il faut examiner en premier.
\end{boxmetier}

\begin{boxmetier}{Inspecteur}
Un rapport qui conclut «~régularité satisfaisante~» doit dire sur quelle
population il porte. Écrivez-le~: «~sur les 1\,184 dossiers classés de
décembre~», et non «~sur les dépenses de décembre~».
\end{boxmetier}

\begin{boxmetier}{Contrôleur financier}
Les paiements de fin de gestion concentrent le risque~: volume élevé,
délais courts, personnel réduit. Une revue des mandats émis hors des
heures de service et des RIB modifiés dans les derniers jours de
l'exercice est un contrôle simple, et il aurait arrêté le second paiement.
\end{boxmetier}

\separateur

% ============================================================================
\section{Le chemin du cours}
\label{sec:pro-chemin}
% ============================================================================

Le cours se déroule en huit journées de quatre heures. Les deux dernières
heures de la huitième journée sont offertes.

\begin{tableaucours}{L{1.3cm}L{8.6cm}L{2.3cm}}{%
    \tableauHeader{Jour} &
    \tableauHeader{Question traitée} &
    \tableauHeader{Chapitres}
}
J01 & Qu'est-ce qu'auditer un système d'information public~? & P00,
\ref{chap:audit-public} \\
J02 & Qu'est-ce qu'un système de contrôle interne, et comment
l'évaluer~? & \ref{chap:sci}, \ref{chap:eval-sci} \\
J03 & Qui contrôle quoi, et selon quelles normes~? &
\ref{chap:institutions}, \ref{chap:normes} \\
J04 & Comment préparer la mission et obtenir la preuve~? &
\ref{chap:preparer}, \ref{chap:preuve} \\
J05 & Comment conclure~? Que vaut le contrôle des accès~? &
\ref{chap:constater}, \ref{chap:acces} \\
J06 & Que valent l'exploitation et les contrôles des applications~? &
\ref{chap:exploitation}, \ref{chap:applications} \\
J07 & Que disent les données~? Que valent les projets et les
prestataires~? & \ref{chap:donnees-fraude}, \ref{chap:gouvernance} \\
J08 & Une mission complète, en équipe & \ref{chap:cas-integrateur} \\
\end{tableaucours}
\vspace{-0.8em}
\begin{center}{\small\itshape Tableau~P.5. Les huit journées}\end{center}

Trois dossiers accompagnent le cours. Le MINÉQUIR revient dans les
chapitres sur le contrôle interne, la mission, les accès et les
applications. La \emph{CAPNO}, caisse publique de prévoyance de Garoua,
sert aux chapitres sur la preuve et l'analyse des données~: son fichier de
paie contient des agents qui n'existent pas. Le projet \emph{e-RECETTES},
à Bafoussam, sert aux chapitres sur l'exploitation, les projets et au cas
intégrateur.

\separateur

\begin{boxresume}
\begin{itemize}
    \item \textbf{Le dossier n'est pas le système.} Dans une administration
          informatisée, le classeur n'est qu'une sortie de l'application~;
          la population d'un contrôle se prend dans le système.
    \item \textbf{Reconstituer avant de juger.} Chaque événement est
          rattaché à une pièce~; ce qui n'a pas de pièce reste une
          déclaration.
    \item \textbf{Constatation et analyse se séparent.} La constatation est
          vérifiable par un tiers~; l'analyse porte un niveau de
          confiance.
    \item \textbf{La question du contrôle interne.} Ce qui aurait dû
          empêcher ou détecter l'irrégularité, et pourquoi cela n'a pas
          fonctionné.
    \item \textbf{Trois réponses.} \OUI{}, \NON{}, \NEK{}~; l'absence d'une
          trace exigée vaut \NON{}, une pièce non encore obtenue vaut
          \NEK{}.
    \item \textbf{Une compétence écrite.} Le décret n°~2013/287 range le
          contrôle des systèmes d'information parmi les cinq contrôles du
          CONSUPE.
    \item \textbf{Trois réflexes.} Le mandat, la preuve, le contradictoire.
\end{itemize}
\end{boxresume}

\begin{boxexo}[title={\ding{45}\quad Exercice~P.1. Constatation ou analyse~? \niveaubase}]
Classez chacun des énoncés suivants et justifiez en une ligne.
\begin{enumerate}[(a)]
    \item «~Le compte \texttt{mbarga.j} a ordonnancé le mandat 2025-09342
          le 28~décembre~2025 à 23~h~14 (journal extrait le 23~mars,
          empreinte au dossier).~»
    \item «~M.~MBARGA a détourné 187\,450\,000~FCFA.~»
    \item «~Sept mandats de décembre n'ont pas de dossier papier.~»
    \item «~Le contrôle des modifications de RIB n'est pas efficace.~»
    \item «~La cellule informatique semble débordée.~»
\end{enumerate}
\end{boxexo}

\begin{boxexo}[title={\ding{45}\quad Exercice~P.2. Coter un contrôle \niveaubase}]
Pour chacune des situations suivantes, donnez la cote (\OUI{}, \NON{} ou
\NEK{}) et justifiez.
\begin{enumerate}[(a)]
    \item La procédure exige un procès-verbal mensuel de revue des comptes
          actifs~; la DAF n'en produit aucun pour 2025.
    \item Le chef de la cellule informatique affirme que les mots de passe
          expirent tous les 90~jours~; le paramétrage n'a pas encore été
          montré.
    \item La matrice des profils, extraite par la vérificatrice, ne donne
          à aucun profil à la fois la liquidation et l'ordonnancement.
    \item Le test de doublon a été paramétré en juin 2026, après les
          faits.
\end{enumerate}
\end{boxexo}

\begin{boxexo}[title={\ding{45}\quad Exercice~P.3. Deux populations \niveaumoyen}]
L'inspection de janvier a examiné 40 dossiers parmi 1\,184.
\begin{enumerate}[(a)]
    \item Quelle était la probabilité que le mandat 2025-09342 figure dans
          son échantillon~? Justifiez.
    \item Proposez, en trois étapes, la démarche qui aurait permis de le
          détecter sans augmenter la taille de l'échantillon.
    \item Rédigez la phrase de conclusion que l'inspection aurait dû
          écrire, compte tenu de sa population réelle.
\end{enumerate}
\end{boxexo}

\begin{boxexo}[title={\ding{45}\quad Exercice~P.4. Le premier constat \niveauexpert}]
À partir du Tableau~P.4 et du journal (Tableau~P.3)~:
\begin{enumerate}[(a)]
    \item Rédigez, en séparant constatation et analyse, le constat portant
          sur le contrôle~2 (séparation entre liquidation et
          ordonnancement).
    \item Indiquez le critère que vous citez, et la pièce qui fonde la
          cote.
    \item Formulez une recommandation adressée au directeur des affaires
          financières, et une adressée au chef de la cellule
          informatique.
    \item Dites quelles réponses vous attendez de chacun pendant la
          procédure contradictoire, et ce qui pourrait vous faire changer
          la cote.
\end{enumerate}
\end{boxexo}

\begin{boxinfo}
\textbf{Transition.} Le prologue a montré un contrôle qui n'a rien vu et
une vérificatrice qui a vu, parce qu'elle est partie du système. Le
Chapitre~\ref{chap:audit-public} pose le vocabulaire~: ce que signifient
contrôler, auditer, vérifier, et ce qui change quand l'objet du contrôle
est un système d'information public.
\end{boxinfo}
}{}

    % ======================================================================
    % PARTIE 1 — FONDEMENTS ET CONTRÔLE INTERNE (J01, J02)
    % « On n'audite pas un système, on évalue la maîtrise qu'on en a. »
    % Standards de référence : ISSAI 100, COSO 2013, INTOSAI GOV 9100 ;
    %   décret n° 2013/287 (art. 2).
    % Niveau : fondamental à intermédiaire
    % ======================================================================
    \part{Fondements et Contrôle Interne}
    \label{part:fondements}

        \IfFileExists{partie1/P01_C01_controler_auditer_si_public.tex}{% ============================================================================
% Fichier  : partie1/P01_C01_controler_auditer_si_public.tex
% Titre    : Contrôler et Auditer un SI Public
% Auteur   : MINKA MI NGUIDJOÏ Thierry Emmanuel
% Version  : 0.1 -- novembre 2026
% ============================================================================
% Sources : INF3226 C01 (resserré) ; ISSAI 100 (via le manuel IDI) ;
% GUID 5100 ; décret n° 2013/287, art. 2. Seuls A0 et A1 (HAL) sont cités
% comme travaux de l'auteur.
% ============================================================================

\chapter{Contrôler et Auditer un SI Public}
\label{chap:audit-public}

\epigraph{%
    « Auditer, c'est donner à un tiers une assurance sur ce qu'il ne peut
    pas vérifier lui-même. »%
}{Principe du cours}

\niveauChapitre{Fondamental}{Prologue. Aucune connaissance informatique
préalable.}

\begin{boxobjectifs}
\begin{itemize}
    \item Distinguer contrôle, audit, inspection et vérification, et situer
          chacun dans l'administration financière camerounaise.
    \item Expliquer la différence entre assurance raisonnable et assurance
          limitée.
    \item Reconnaître les trois types d'audit du secteur public (financier,
          de conformité, de performance) et ce que chacun demande au
          système d'information.
    \item Situer le contrôle des systèmes d'information parmi les cinq
          contrôles du CONSUPE.
    \item Cartographier les systèmes d'information des finances publiques et
          leurs points de contrôle.
    \item Expliquer ce que l'informatique change à la preuve.
\end{itemize}
\end{boxobjectifs}

\bigskip

% ============================================================================
\section*{Pourquoi ce chapitre}
% ============================================================================

Les auditeurs du PSSFP connaissent déjà le contrôle des finances
publiques. Ce chapitre ne le leur réapprend pas~: il fixe un vocabulaire
commun et montre ce qui change quand l'objet du contrôle devient un système
d'information. La question directrice est simple~: \emph{que faut-il
regarder, et dans quel ordre, quand la dépense passe par une
application~?}

\separateur

% ============================================================================
\section{Contrôler, auditer, inspecter, vérifier}
\label{sec:c01-notions}
% ============================================================================

Les mots du contrôle sont employés indifféremment dans les services. Le
cours les distingue, parce que chacun engage une démarche et une
responsabilité différentes.

\begin{boxdef}
\termecle{Contrôle}\index{contrôle}~: dispositif ou action, intégré au
fonctionnement de l'entité, qui vise à empêcher ou à détecter une anomalie
(un visa, un rapprochement, un blocage automatique). Il est exercé par les
services eux-mêmes.

\medskip
\termecle{Audit}\index{audit}~: examen méthodique, indépendant et
documenté, qui compare une situation à des critères connus à l'avance et
qui conduit à une opinion motivée, communiquée à un tiers.

\medskip
\termecle{Inspection}\index{inspection}~: vérification ponctuelle de la
régularité d'opérations ou du fonctionnement d'un service, souvent
déclenchée par un événement, au nom de l'autorité hiérarchique.

\medskip
\termecle{Vérification}\index{vérification}~: dans la langue du décret
n°~2013/287, l'activité des vérificateurs du Contrôle Supérieur de l'État,
conduite sur pièces et sur place~; elle recouvre l'audit externe du
secteur public.
\end{boxdef}

La différence tient en deux questions. \textbf{Qui agit~?} Le contrôle est
exercé par l'entité~; l'audit et la vérification par quelqu'un qui n'a pas
participé aux opérations. \textbf{Pour qui~?} Le contrôle sert l'entité~;
l'audit sert un tiers (le Président de la République pour le CONSUPE, le
juge des comptes, le parlement, un conseil d'administration) qui ne peut
pas vérifier lui-même.

\subsection{Les éléments d'un audit}

Tout audit, quel que soit son objet, manipule les mêmes cinq éléments. Le
cours les emploie dans ce sens précis.

\begin{tableaucours}{L{3.0cm}L{9.2cm}}{%
    \tableauHeader{Élément} &
    \tableauHeader{Définition et exemple (dossier MINÉQUIR)}
}
Critère & Exigence de référence, connue avant le test~: «~un mandat ne peut
pas être payé deux fois~» \\
Preuve & Information vérifiable que l'auditeur obtient et conserve~:
l'extraction des 9\,611 mandats, son empreinte \\
Constatation & Comparaison de la preuve au critère~: «~le mandat 2025-09342
reprend la facture du mandat 2025-08817~» \\
Conclusion & Résultat de l'ensemble des constatations sur un objectif~:
«~le contrôle des doublons n'existe pas dans SIDÉPENSE~» \\
Opinion & Forme que prend la conclusion dans le rapport, avec son niveau
d'assurance \\
\end{tableaucours}
\vspace{-0.8em}
\begin{center}{\small\itshape Tableau~1.1. Les cinq éléments d'un
audit}\end{center}

\begin{boxalert}
\textbf{Violation typique.} Choisir le critère après avoir vu les faits~:
«~les droits de M.~MBARGA étaient excessifs~», sans dire quelle règle les
limitait.

\smallskip
\textbf{Règle~:} le critère est écrit avant le test, et cité dans le
constat. Un constat sans critère est une opinion personnelle.
\end{boxalert}

\separateur

% ============================================================================
\section{Assurance raisonnable, assurance limitée}
\label{sec:c01-assurance}
% ============================================================================

Aucun auditeur n'examine tout. Il donne donc une assurance, jamais une
certitude.

\begin{boxdef}
\termecle{Assurance raisonnable}\index{assurance raisonnable}~: niveau de
confiance élevé, mais non absolu, obtenu par des travaux étendus. Elle
s'exprime sous une forme positive~: «~les contrôles examinés fonctionnent
efficacement~».

\medskip
\termecle{Assurance limitée}\index{assurance limitée}~: niveau de
confiance plus faible, obtenu par des travaux réduits. Elle s'exprime sous
une forme négative~: «~nous n'avons rien relevé qui laisse penser que…~».
\end{boxdef}

Le choix se fait au moment du mandat, pas à la fin de la mission. Une
inspection de deux jours ne peut donner qu'une assurance limitée~; si son
rapport conclut sous une forme positive, il promet plus qu'il n'a vérifié.
C'est exactement le défaut du rapport de janvier dans le prologue
(Chapitre~\ref{chap:prologue})~: «~régularité satisfaisante~» sur une
population incomplète.

\separateur

% ============================================================================
\section{Trois types d'audit dans le secteur public}
\label{sec:c01-types}
% ============================================================================

Les normes de l'INTOSAI distinguent trois types d'audit du secteur
public~\cite{idi2021conformite}. Chacun interroge le système d'information
d'une manière différente.

\begin{tableaucours}{L{2.6cm}L{4.6cm}L{5.0cm}}{%
    \tableauHeader{Type} &
    \tableauHeader{Question} &
    \tableauHeader{Ce qu'il demande au SI}
}
Audit financier & Les états financiers et les comptes sont-ils fiables~? &
Les données produites par l'application sont-elles exactes et
exhaustives~? \\
Audit de conformité & Les opérations respectent-elles les lois, les
règlements et les règles internes~? & L'application impose-t-elle les
règles (séparation des tâches, plafonds, visas)~? \\
Audit de performance & Les ressources sont-elles employées avec économie,
efficience et efficacité~? & Le système atteint-il ses objectifs, à quel
coût, et ses données permettent-elles de le savoir~? \\
\end{tableaucours}
\vspace{-0.8em}
\begin{center}{\small\itshape Tableau~1.2. Les trois types d'audit et le
système d'information}\end{center}

L'audit des systèmes d'information n'est pas un quatrième type à côté des
trois autres. Il les sert tous les trois~: on audite un SI pour savoir si
l'on peut se fier à ses données (financier), s'il fait respecter les règles
(conformité), ou s'il rend le service attendu (performance). Les lignes
directrices GUID~5100 de l'INTOSAI le disent ainsi~: un audit des SI peut
être une mission autonome, de performance ou de conformité, ou une partie
d'une mission plus large, financière, de conformité ou de
performance~\cite{intosai2019guid5100}.

\subsection{Les cinq contrôles du CONSUPE}

Le décret n°~2013/287 range le contrôle des systèmes d'information parmi
les cinq contrôles qu'effectuent les services du Contrôle supérieur de
l'État~\cite{decret2013287consupe}.

\begin{boxjuris}
\textbf{\loicam{Décret n°~2013/287}, article~2, alinéa~1}~: \textit{« Les
services du Contrôle supérieur de l'Etat constituent l'institution
supérieure de contrôle des finances publiques (ISC) du Cameroun. Ils sont
chargés de l'audit externe. »}

\medskip
\textbf{Article~2, alinéa~3}~: \textit{« Dans l'exercice de leurs
attributions, les services du Contrôle supérieur de l'Etat effectuent~:
un contrôle de conformité et de régularité~; un contrôle financier~; un
contrôle de performance~; un contrôle de l'environnement~; un contrôle des
systèmes d'information. »}

\medskip
\textbf{Portée pour l'audit.} Le décret reprend les trois types de
l'INTOSAI (conformité, financier, performance) et leur ajoute deux
contrôles, dont celui des systèmes d'information. Le vérificateur peut
donc conduire une mission dont l'objet principal est un SI, ou intégrer un
volet SI à une mission de régularité ou de performance.
\end{boxjuris}

\separateur

% ============================================================================
\section{Les systèmes d'information des finances publiques}
\label{sec:c01-cartographie}
% ============================================================================

\subsection{Une chaîne de systèmes}

Les finances publiques reposent sur une chaîne d'applications, de la
préparation du budget à la reddition des comptes. Le cours ne décrit pas
les applications réelles de l'État camerounais, qui évoluent et dont le
détail n'est pas public~: il raisonne sur des fonctions, que l'auditeur
retrouve sous des noms différents d'une administration à l'autre.

\begin{tableaucours}{L{2.6cm}L{4.8cm}L{4.8cm}}{%
    \tableauHeader{Fonction} &
    \tableauHeader{Ce que fait le système} &
    \tableauHeader{Risque typique}
}
Budget & Prépare, vote, notifie et modifie les crédits & Crédits modifiés
sans acte, virements non autorisés \\
Dépense & Engage, liquide, ordonnance, transmet au comptable & Double
paiement, service fait fictif, fractionnement \\
Recette & Liquide, émet et recouvre les titres de recettes & Titres
annulés sans motif, recouvrements non reversés \\
Paie & Calcule et verse la solde & Agents fictifs, éléments de solde
indus \\
Trésorerie & Exécute les paiements, tient les comptes de disponibilités
& Paiement sur un compte bancaire modifié \\
Comptabilité & Produit les comptes et les états financiers & États qui ne
correspondent pas aux opérations \\
\end{tableaucours}
\vspace{-0.8em}
\begin{center}{\small\itshape Tableau~1.3. Les fonctions des systèmes
financiers publics et leurs risques}\end{center}

\subsection{Les points de contrôle de la chaîne}

Dans une chaîne de systèmes, les anomalies naissent souvent aux
\textbf{jointures}~: là où une application transmet des données à une
autre (de la dépense vers la trésorerie, de la paie vers la comptabilité),
et là où une donnée de référence est modifiée (un RIB, un barème, un
plafond). L'auditeur y place ses premiers tests.

\begin{boxscenario}{Dossier MINÉQUIR. Trois jointures}
\textbf{Constatation.} Dans SIDÉPENSE, trois données conditionnent tous
les paiements~: le fichier des fournisseurs (et leurs RIB), la table des
profils d'utilisateurs, et le fichier de transmission des mandats au poste
comptable.

\medskip
\textbf{Analyse} (confiance élevée). Le double paiement du prologue a
utilisé deux de ces trois jointures~: un RIB modifié et un profil resté
actif. Une mission qui teste ces trois points en premier couvre une grande
part du risque de la chaîne, pour un effort limité.
\end{boxscenario}

\separateur

% ============================================================================
\section{Ce que l'informatique change à la preuve}
\label{sec:c01-preuve}
% ============================================================================

Le contrôle sur pièces repose sur une évidence~: la pièce existe, elle est
signée, elle est classée. Dans une administration informatisée, chacune de
ces évidences change de forme.

\begin{tableaucours}{L{3.2cm}L{4.4cm}L{4.6cm}}{%
    \tableauHeader{Sur papier} &
    \tableauHeader{Dans le système} &
    \tableauHeader{Ce que l'auditeur doit vérifier}
}
La pièce & Un enregistrement dans une base & Qu'il n'a pas été modifié
après coup \\
La signature & Un droit d'accès utilisé & Que le titulaire du compte est
bien celui qui a agi \\
Le visa & Un contrôle automatique & Que le paramétrage n'a pas été changé
\\
Le classeur & Une extraction & Qu'elle est complète, et qu'il l'a obtenue
lui-même \\
L'archive & Une sauvegarde & Qu'elle peut être restaurée \\
\end{tableaucours}
\vspace{-0.8em}
\begin{center}{\small\itshape Tableau~1.4. Ce que l'informatique change à
la preuve}\end{center}

Deux conséquences en découlent. D'abord, la \textbf{population} d'un test
est dans le système, pas dans les classeurs. Ensuite, une preuve
électronique ne vaut que si l'auditeur peut en établir l'intégrité~: une
extraction remise sur une clé, sans empreinte ni rapprochement, est une
déclaration de l'audité. Le Chapitre~\ref{chap:preuve} donne la manière de
faire.

\begin{boxinfo}
\textbf{Une réponse de plus.} Parce que la preuve électronique peut manquer
ou être incertaine, le cours admet trois réponses pour chaque contrôle
testé~: \OUI{}, \NON{} et \NEK{} (connaissance insuffisante). Elles
viennent d'un programme de recherche sur l'audit, l'auditologie, qui
cherche à mesurer ce que l'auditeur ne sait
pas~\cite{minka2026auditologie,minka2026nek}. Le cours n'en garde que
l'essentiel~: dire ce que l'on ne sait pas, plutôt que de le cacher sous un
\OUI{} ou un \NON{}.
\end{boxinfo}

\begin{boxmetier}{Vérificateur}
Avant la mission, demandez la liste des applications qui produisent les
opérations de votre périmètre, avec, pour chacune, le service
propriétaire et l'interlocuteur technique. Cette liste est la première
pièce de votre dossier permanent.
\end{boxmetier}

\begin{boxmetier}{Juge des comptes}
Un compte produit par une application n'est fiable que si les contrôles
de l'application le sont. Lorsque vous relevez une anomalie dans un
compte, demandez quel contrôle de l'application aurait dû l'empêcher~:
la réponse oriente l'instruction.
\end{boxmetier}

\begin{boxmetier}{Auditeur interne}
Votre rapport n'est pas lu que par votre direction~: l'article~5 du
décret n°~2013/287 le fait transmettre au CONSUPE
(Chapitre~\ref{chap:institutions}). Rédigez vos constats SI comme s'ils
devaient être repris par un vérificateur.
\end{boxmetier}

\separateur

\begin{boxresume}
\begin{itemize}
    \item \textbf{Contrôler n'est pas auditer.} Le contrôle est exercé par
          l'entité~; l'audit est un examen indépendant au profit d'un
          tiers.
    \item \textbf{Cinq éléments.} Critère, preuve, constatation,
          conclusion, opinion~; le critère est écrit avant le test.
    \item \textbf{Une assurance, pas une certitude.} Raisonnable ou limitée,
          choisie au mandat, et formulée en conséquence.
    \item \textbf{Trois types d'audit, un même SI.} L'audit des SI sert
          l'audit financier, de conformité et de performance.
    \item \textbf{Une compétence écrite.} Le décret n°~2013/287 compte le
          contrôle des SI parmi les cinq contrôles du CONSUPE.
    \item \textbf{Les jointures d'abord.} Les données de référence et les
          transmissions entre applications concentrent le risque.
    \item \textbf{La preuve change de forme.} Population dans le système,
          intégrité à établir, trois réponses possibles.
\end{itemize}
\end{boxresume}

\begin{boxexo}[title={\ding{45}\quad Exercice~1.1. Qui fait quoi~? \niveaubase}]
Pour chacune des situations suivantes, dites s'il s'agit d'un contrôle,
d'un audit, d'une inspection ou d'une vérification, et justifiez.
\begin{enumerate}[(a)]
    \item Le chef du service de la liquidation vérifie chaque matin la
          liste des mandats saisis la veille.
    \item Un vérificateur du CONSUPE examine l'exécution du budget 2025 de
          la DAF du MINÉQUIR.
    \item L'inspecteur général du ministère se rend à la DAF après un
          article de presse.
    \item SIDÉPENSE refuse un engagement supérieur au crédit disponible.
\end{enumerate}
\end{boxexo}

\begin{boxexo}[title={\ding{45}\quad Exercice~1.2. Les cinq éléments \niveaubase}]
Dans le dossier CAPNO, le fichier de paie de décembre compte 4\,312
agents, dont 12 partagent le même compte bancaire avec au moins un autre
agent. Identifiez, dans cette phrase et dans la suite que vous imaginerez,
un critère, une preuve, une constatation et une conclusion.
\end{boxexo}

\begin{boxexo}[title={\ding{45}\quad Exercice~1.3. Choisir l'assurance \niveaumoyen}]
On vous confie trois jours pour examiner la paie de la CAPNO.
\begin{enumerate}[(a)]
    \item Quel niveau d'assurance pouvez-vous donner~? Pourquoi~?
    \item Rédigez la phrase de conclusion correspondante, dans les deux
          cas~: vous n'avez rien trouvé~; vous avez trouvé 37 agents qui
          ne figurent pas au fichier du personnel.
\end{enumerate}
\end{boxexo}

\begin{boxexo}[title={\ding{45}\quad Exercice~1.4. Cartographier \niveauexpert}]
Pour le projet e-RECETTES de Bafoussam (dématérialisation des recettes),
dressez le tableau des fonctions du système, des jointures avec les
autres systèmes (trésorerie, comptabilité), et du risque principal à
chaque jointure. Dites quelles trois jointures vous testeriez en premier,
et pourquoi.
\end{boxexo}

\begin{boxinfo}
\textbf{Transition.} Ce chapitre a posé le vocabulaire et désigné
l'objet~: les systèmes qui portent les finances publiques. Le
Chapitre~\ref{chap:sci} présente ce que l'auditeur y cherche avant
tout~: le système de contrôle interne.
\end{boxinfo}
}{}
        \IfFileExists{partie1/P01_C02_systeme_controle_interne.tex}{% ============================================================================
% Fichier  : partie1/P01_C02_systeme_controle_interne.tex
% Titre    : Le Système de Contrôle Interne en Environnement Informatisé
% Auteur   : MINKA MI NGUIDJOÏ Thierry Emmanuel
% Version  : 0.1 -- novembre 2026
% ============================================================================
% Sources : COSO 2013, traduction IFACI (Normes) ; INF3226 C03 (resserré) ;
% décret n° 2013/287 (art. 28). INTOSAI GOV 9100 cité sans citation
% textuelle (texte absent du dossier Normes).
% Faits MINÉQUIR : profils de SIDÉPENSE (P00) ; matrice de C09 à garder
% cohérente : ENG (engagement), LIQ (liquidation), ORD (ordonnancement),
% FRS (fichier des fournisseurs), ADM (administration des comptes).
% ============================================================================

\chapter{Le Système de Contrôle Interne en Environnement Informatisé}
\label{chap:sci}
\chaptermark{Le système de contrôle interne}

\epigraph{%
    « L'organisation est-elle sous contrôle~? »%
}{Question à laquelle le Référentiel COSO veut permettre de répondre}

\niveauChapitre{Fondamental}{Chapitre~\ref{chap:audit-public}.}

\begin{boxobjectifs}
\begin{itemize}
    \item Définir le contrôle interne selon le COSO~2013 et en expliquer
          les quatre idées clés.
    \item Énoncer les trois catégories d'objectifs et les illustrer dans
          une administration financière.
    \item Présenter les cinq composantes et les dix-sept principes, et dire
          ce que chacun exige du système d'information.
    \item Classer un contrôle selon son moment (préventif, détectif,
          correctif) et son mode (manuel, automatisé, semi-automatisé).
    \item Distinguer les contrôles généraux informatiques des contrôles
          d'application, et expliquer pourquoi les premiers conditionnent
          les seconds.
    \item Traduire la séparation entre ordonnateur et comptable, et entre
          les étapes de la dépense, en profils d'une application.
    \item Énoncer les limites du contrôle interne.
\end{itemize}
\end{boxobjectifs}

\bigskip

% ============================================================================
\section*{Pourquoi ce chapitre}
% ============================================================================

Le prologue a posé une question~: qu'est-ce qui aurait dû empêcher le
double paiement~? La réponse porte un nom, le \emph{système de contrôle
interne}. Ce chapitre dit ce qu'il est~; le Chapitre~\ref{chap:eval-sci}
dit comment l'évaluer. Ensemble, ils forment le cœur du cours~: tous les
chapitres de la Partie~IV appliquent la même grille, domaine par domaine.

\separateur

% ============================================================================
\section{Qu'est-ce que le contrôle interne~?}
\label{sec:c02-definition}
% ============================================================================

\subsection{La définition du COSO}

Le cadre de référence le plus utilisé au monde est celui du \sigle{COSO}
(\textit{Committee of Sponsoring Organizations of the Treadway
Commission}), dans sa version de 2013, traduite en français par
l'IFACI~\cite{coso2013ic,coso2013icfr}.

\begin{boxdef}
\termecle{Contrôle interne}\index{contrôle interne}~: processus mis en
œuvre par l'organe de gouvernance, la direction et le personnel, destiné à
fournir une assurance raisonnable quant à la réalisation d'objectifs liés
aux opérations, à l'information (\textit{reporting}) et à la conformité.

\medskip
\termecle{Système de contrôle interne}\index{système de contrôle interne}
(\sigle{SCI})~: l'ensemble organisé des dispositifs de contrôle interne
d'une entité, considéré comme un tout~; c'est lui que l'auditeur évalue.
\end{boxdef}
{\small\itshape Définition du contrôle interne reformulée d'après le
COSO~2013.}

\medskip

La définition contient quatre idées, que l'auditeur public retrouve dans
son propre métier.
\begin{enumerate}
    \item \textbf{Un processus, pas un document.} Un manuel de procédures
          que personne n'applique n'est pas un contrôle interne.
    \item \textbf{L'affaire de tous.} Il est porté par l'ordonnateur, le
          chef de service, l'agent de saisie, l'informaticien, et pas
          seulement par le contrôleur.
    \item \textbf{Une assurance raisonnable.} Le coût d'un contrôle ne doit
          pas dépasser le risque qu'il traite~; aucun dispositif n'empêche
          tout.
    \item \textbf{Des objectifs.} On ne contrôle pas «~en général~», mais
          en vue d'objectifs précis~: sans objectif, pas de risque, et sans
          risque, pas de contrôle à justifier.
\end{enumerate}

\subsection{Trois catégories d'objectifs}

\begin{tableaucours}{L{2.6cm}L{4.6cm}L{5.0cm}}{%
    \tableauHeader{Catégorie} &
    \tableauHeader{Question} &
    \tableauHeader{Exemple au MINÉQUIR}
}
Opérations & Les activités sont-elles efficaces, les ressources
protégées~? & Chaque facture n'est payée qu'une fois, au bon
créancier \\
Information & L'information produite est-elle fiable~? & Le compte
administratif reflète les mandats réellement émis \\
Conformité & Les lois et règlements sont-ils respectés~? & Le circuit
engagement, liquidation, ordonnancement, paiement est respecté \\
\end{tableaucours}
\vspace{-0.8em}
\begin{center}{\small\itshape Tableau~2.1. Les objectifs du contrôle
interne}\end{center}

\subsection{Le contrôle interne dans le secteur public}

L'INTOSAI a adapté le cadre du COSO au secteur public dans ses lignes
directrices INTOSAI~GOV~9100. Elles ajoutent deux préoccupations propres à
l'administration~: la \textbf{protection des ressources publiques} contre
la perte, le gaspillage et l'usage abusif, et la conduite des opérations
dans le respect de l'\textbf{éthique} du service public. Le cours retient
le COSO~2013 comme cadre de travail, parce qu'il est le plus récent et le
seul dont le texte est dans le dossier des normes, et lit chaque principe
avec ces deux préoccupations.

Au Cameroun, l'évaluation des structures de contrôle interne est une
mission écrite du CONSUPE.

\begin{boxjuris}
\textbf{\loicam{Décret n°~2013/287}, article~28, alinéa~1}~\cite{decret2013287consupe}~:
les sections des divisions d'inspection sont chargées, chacune dans son
domaine de compétence, \textit{« de la préparation et de l'exécution des
opérations de vérification~; du suivi et de l'évaluation des activités
des structures de contrôle interne~; de la tenue de la documentation
technique et du fichier. »}

\medskip
\textbf{Portée pour l'audit.} Le vérificateur ne se contente pas de
relever des irrégularités~: il évalue les dispositifs qui auraient dû les
empêcher. L'article~27 place par ailleurs les systèmes d'information dans
le champ d'une section de la division chargée des administrations
publiques.
\end{boxjuris}

\separateur

% ============================================================================
\section{Cinq composantes, dix-sept principes}
\label{sec:c02-coso}
% ============================================================================

Le COSO organise le contrôle interne en \textbf{cinq composantes},
détaillées en \textbf{dix-sept principes}. Un dispositif n'est efficace
que si les cinq composantes sont mises en place, fonctionnent, et
fonctionnent conjointement~: une excellente activité de contrôle ne
compense pas une direction qui la contourne.

\begin{tableaucours}{L{3.2cm}cL{8.0cm}}{%
    \tableauHeader{Composante} &
    \tableauHeader{N°} &
    \tableauHeader{Principe (reformulé)}
}
Environnement de contrôle & 1 & S'engager en faveur de l'intégrité et de
l'éthique \\
 & 2 & Faire surveiller le contrôle interne par une instance indépendante
de la direction \\
 & 3 & Définir les structures, les pouvoirs et les responsabilités \\
 & 4 & Développer et retenir des personnes compétentes \\
 & 5 & Faire rendre compte chacun de ses responsabilités de contrôle \\
Évaluation des risques & 6 & Fixer des objectifs assez clairs pour en
identifier les risques \\
 & 7 & Identifier et analyser les risques \\
 & 8 & Évaluer le risque de fraude \\
 & 9 & Identifier et analyser les changements significatifs \\
Activités de contrôle & 10 & Choisir et mettre en place des activités de
contrôle qui réduisent les risques \\
 & 11 & Choisir et mettre en place des contrôles généraux informatiques \\
 & 12 & Déployer les contrôles par des politiques et des procédures \\
Information et communication & 13 & Produire et utiliser une information
pertinente et de qualité \\
 & 14 & Communiquer en interne sur le contrôle interne \\
 & 15 & Communiquer avec les tiers \\
Pilotage & 16 & Conduire des évaluations continues ou ponctuelles \\
 & 17 & Évaluer les déficiences et les communiquer à temps \\
\end{tableaucours}
\vspace{-0.8em}
\begin{center}{\small\itshape Tableau~2.2. Les composantes et les principes
du COSO~2013 (reformulation du cours)}\end{center}

\subsection{Les principes vus depuis le système d'information}

L'auditeur des SI lit chaque principe avec une question~: \emph{qu'est-ce
que ce principe exige du système d'information, et quelle pièce le
prouve~?} Le Tableau~2.3 en donne la réponse pour une administration
financière. C'est la base du questionnaire de contrôle interne du
Chapitre~\ref{chap:eval-sci}.

\begin{tableaucours}{cL{5.6cm}L{6.0cm}}{%
    \tableauHeader{N°} &
    \tableauHeader{Ce que le principe exige du SI} &
    \tableauHeader{Pièce que l'auditeur peut demander}
}
1 & Une charte d'usage des outils informatiques, connue des agents &
Charte, accusés de réception \\
2 & Une instance qui examine les risques informatiques & Comptes rendus du
comité de direction ou d'un comité dédié \\
3 & Des rôles écrits pour la cellule informatique et les propriétaires
d'application & Organigramme, décisions de nomination, délégations \\
4 & Des compétences informatiques entretenues & Plan de formation \\
5 & Des propriétaires d'application qui valident les droits & Revues des
habilitations signées \\
6 & Des objectifs écrits de disponibilité et de sécurité & Schéma
directeur, politique de sécurité \\
7 & Une analyse des risques informatiques & Cartographie des risques \\
8 & Une analyse des fraudes possibles par le système & Scénarios de
fraude (Chapitre~\ref{chap:donnees-fraude}) \\
9 & L'analyse des changements majeurs (nouvelle application, migration)
& Dossier de mise en service \\
10 & Des contrôles dans les applications & Paramétrage, contrôles
bloquants \\
11 & Des contrôles généraux informatiques & Accès, changements,
exploitation (section~\ref{sec:c02-cgi}) \\
12 & Des procédures informatiques écrites et appliquées & Procédures,
preuves d'exécution \\
13 & Des données fiables pour le pilotage & Rapprochements, contrôles de
qualité \\
14 & Une communication interne sur les règles & Notes de service \\
15 & Des échanges fiables avec les tiers (banques, trésor, prestataires)
& Conventions d'échange, contrats \\
16 & Une surveillance des contrôles & Tableaux de bord, revues
périodiques \\
17 & Un suivi des déficiences & Plan d'action, suivi des
recommandations \\
\end{tableaucours}
\vspace{-0.8em}
\begin{center}{\small\itshape Tableau~2.3. Les dix-sept principes vus
depuis le système d'information}\end{center}

\begin{boxinfo}
\textbf{Le principe~11 concerne directement l'auditeur des SI.} Le COSO
y demande de choisir et de mettre en place des contrôles généraux sur les
technologies~: les accès, les changements, l'exploitation. Les
Chapitres~\ref{chap:acces} et~\ref{chap:exploitation} les évaluent.
\end{boxinfo}

\separateur

% ============================================================================
\section{Classer les contrôles}
\label{sec:c02-types}
% ============================================================================

\subsection{Selon le moment}

\begin{boxdef}
\termecle{Contrôle préventif}\index{contrôle préventif}~: empêche
l'anomalie de se produire (SIDÉPENSE refuse un engagement supérieur au
crédit disponible).

\medskip
\termecle{Contrôle détectif}\index{contrôle détectif}~: repère l'anomalie
après coup (la revue mensuelle des mandats émis hors des heures de
service).

\medskip
\termecle{Contrôle correctif}\index{contrôle correctif}~: corrige
l'anomalie et ses effets (le recouvrement d'un paiement indu).
\end{boxdef}

Un dispositif solide combine les trois. Le préventif est préférable,
parce qu'il évite le dommage~; mais il peut être contourné, et seul le
détectif révèle alors le contournement. Au MINÉQUIR, tous les contrôles
préventifs du circuit ont manqué, et seuls deux contrôles détectifs ont
fonctionné (Chapitre~\ref{chap:prologue}).

\subsection{Selon le mode d'exécution}

\begin{tableaucours}{L{2.8cm}L{4.6cm}L{4.8cm}}{%
    \tableauHeader{Mode} &
    \tableauHeader{Forces} &
    \tableauHeader{Faiblesses}
}
Manuel (un visa) & Jugement humain, adaptable & Variable, oublié en fin
de gestion, preuve souvent absente \\
Automatisé (un blocage) & Constant, exhaustif & Ne vaut que si le
paramétrage est protégé~; une erreur se répète à chaque opération \\
Semi-automatisé (une liste d'anomalies examinée par un agent) & Combine
exhaustivité et jugement & Dépend de l'exactitude de la liste et de la
diligence de l'agent \\
\end{tableaucours}
\vspace{-0.8em}
\begin{center}{\small\itshape Tableau~2.4. Les modes d'exécution d'un
contrôle}\end{center}

\subsection{Selon la portée}

\begin{boxdef}
\termecle{Contrôles généraux informatiques}\index{contrôles généraux
informatiques} (\sigle{CGI}, \textit{IT general controls})~: contrôles
qui s'appliquent à l'environnement dans lequel les applications sont
exploitées et maintenues~: gestion des accès, gestion des changements,
exploitation (sauvegardes, traitements, incidents).

\medskip
\termecle{Contrôles d'application}\index{contrôles d'application}~:
contrôles intégrés à une application, qui agissent sur les opérations~:
contrôles de saisie, de traitement, d'interface et de sortie.
\end{boxdef}

\begin{boxalert}
\textbf{Violation typique.} Constater que SIDÉPENSE refuse tout mandat
supérieur au crédit disponible, et conclure que le contrôle budgétaire
est efficace. Or un compte d'administration peut modifier le crédit
directement dans la base, sans acte et sans trace.

\smallskip
\textbf{Règle~:} un contrôle d'application ne vaut que si les contrôles
généraux qui le protègent (accès, changements) sont efficaces.
L'auditeur les teste en premier.
\end{boxalert}

\separateur

% ============================================================================
\section{Les contrôles généraux informatiques}
\label{sec:c02-cgi}
% ============================================================================

\begin{tableaucours}{L{2.6cm}L{5.0cm}L{4.6cm}}{%
    \tableauHeader{Famille} &
    \tableauHeader{Objectif} &
    \tableauHeader{Contrôles typiques}
}
Accès & Seules les personnes autorisées accèdent aux applications et aux
données, avec les seuls droits nécessaires & Création et retrait des
comptes, revue des droits, comptes à privilèges \\
Changements & Aucune modification de programme ou de paramètre n'est mise
en service sans être demandée, testée et approuvée & Demande écrite,
tests, approbation, journal des mises en service \\
Exploitation & Les traitements s'exécutent comme prévu, les incidents sont
traités, les données sont sauvegardées & Planification, suivi des
incidents, sauvegardes testées \\
Acquisition et développement & Les nouvelles applications répondent aux
besoins et intègrent les contrôles & Cahier des charges, recette,
réception \\
\end{tableaucours}
\vspace{-0.8em}
\begin{center}{\small\itshape Tableau~2.5. Les quatre familles de
contrôles généraux informatiques}\end{center}

L'ordre compte. Si les accès ne sont pas maîtrisés, n'importe qui peut
modifier un paramètre~; si les changements ne le sont pas, un contrôle
automatique peut être désactivé sans que personne ne le sache~; dans les
deux cas, tester les contrôles d'application ne prouve plus rien. Le
cours suit donc cet ordre~: accès (Chapitre~\ref{chap:acces}),
changements et exploitation (Chapitre~\ref{chap:exploitation}), puis
contrôles d'application (Chapitre~\ref{chap:applications}).

\separateur

% ============================================================================
\section{Séparer les tâches dans une application}
\label{sec:c02-separation}
% ============================================================================

\subsection{Du principe au profil}

Le droit des finances publiques sépare l'ordonnateur, qui engage,
liquide et ordonnance, du comptable, qui paie après contrôle~; il sépare
aussi les étapes de la dépense~\cite{minfi2023regimefinancier}. Dans une
administration informatisée, ces séparations n'existent que si
l'application les impose.

\begin{boxdef}
\termecle{Séparation des tâches}\index{séparation des tâches}~: répartition
des étapes d'une opération entre des personnes différentes, de sorte
qu'aucune ne puisse à elle seule réaliser et dissimuler une irrégularité.

\medskip
\termecle{Profil}\index{profil}~: ensemble de droits attribué dans une
application (créer un engagement, liquider, ordonnancer, modifier un
RIB). Un utilisateur peut avoir plusieurs profils~: c'est leur
combinaison qui crée, ou non, un conflit.
\end{boxdef}

\subsection{La matrice des profils de SIDÉPENSE}

SIDÉPENSE connaît cinq profils~: ENG (engagement), LIQ (liquidation), ORD
(ordonnancement), FRS (gestion du fichier des fournisseurs, dont les RIB)
et ADM (administration des comptes). Le Tableau~2.6 indique les
combinaisons qui ne doivent jamais être réunies sur un même compte.

\begin{tableaucours}{L{2.2cm}L{4.4cm}L{5.6cm}}{%
    \tableauHeader{Paire} &
    \tableauHeader{Pourquoi elle est interdite} &
    \tableauHeader{Ce qu'elle permet}
}
LIQ + ORD & L'ordonnancement contrôle la liquidation & Arrêter un montant et
émettre seul le mandat \\
FRS + ORD & Le RIB doit être fixé par un autre que celui qui émet le
mandat & Changer le bénéficiaire puis le payer \\
FRS + LIQ & Même raison, en amont & Créer un service fait au profit d'un
compte choisi \\
ADM + tout profil métier & L'administrateur peut s'attribuer n'importe
quel droit & Contourner toutes les séparations \\
\end{tableaucours}
\vspace{-0.8em}
\begin{center}{\small\itshape Tableau~2.6. Les paires de profils
incompatibles dans SIDÉPENSE}\end{center}

\begin{boxscenario}{Dossier MINÉQUIR. Le compte mbarga.j}
\textbf{Constatation.} À la date du 28~décembre~2025, le compte
\texttt{mbarga.j} détient les profils LIQ, ORD et FRS (matrice des
profils extraite le 23~mars~2026).

\medskip
\textbf{Analyse} (confiance élevée). Le compte réunit les trois paires
interdites du Tableau~2.6 qui ne demandent pas l'administration. Un seul
compte suffisait pour modifier le RIB, liquider et ordonnancer~: le
contrôle interne du circuit reposait entièrement sur la confiance en une
personne, et cette personne n'était plus en poste.
\end{boxscenario}

\subsection{Quand on ne peut pas séparer}

Dans un petit service, trois agents ne peuvent pas tenir cinq profils
sans cumul. La pratique admet alors des \termecle{contrôles
compensatoires}\index{contrôle compensatoire}~: une revue, par une
personne indépendante, de toutes les opérations réalisées par un compte
qui cumule des profils. Le contrôle compensatoire doit être écrit,
exécuté et prouvé~; sinon, le cumul est une déficience.

\separateur

% ============================================================================
\section{Les limites du contrôle interne}
\label{sec:c02-limites}
% ============================================================================

Le COSO rappelle que le contrôle interne, même bien conçu, a des
limites~\cite{coso2013icfr}. L'auditeur doit les connaître, pour ne pas
reprocher à un dispositif ce qu'aucun dispositif ne peut empêcher, et pour
savoir où chercher.

\begin{itemize}
    \item \textbf{Le jugement humain} peut être erroné~; les décisions
          sont prises sous la contrainte du temps.
    \item \textbf{Les défaillances} existent~: fatigue, erreur, oubli,
          surtout en fin de gestion.
    \item \textbf{La collusion} de deux personnes ou plus défait la
          séparation des tâches.
    \item \textbf{Le contournement par la direction}~: un responsable peut
          passer outre un contrôle qu'il a lui-même mis en place.
    \item \textbf{Les événements extérieurs} échappent à l'entité.
\end{itemize}

\begin{boxmetier}{Contrôleur financier}
Le contrôle a priori que vous exercez est une activité de contrôle au sens
du COSO (principe~10). S'il passe par SIDÉPENSE, demandez à voir le profil
qui porte votre visa, et la liste des personnes qui le détiennent~: votre
visa ne vaut que si personne d'autre ne peut l'apposer.
\end{boxmetier}

\begin{boxmetier}{Auditeur interne}
Le COSO fait du pilotage (principes~16 et~17) une composante du contrôle
interne~: votre propre travail en fait partie. Gardez une liste des
déficiences que vous avez signalées et de leur suivi~; le vérificateur la
demandera.
\end{boxmetier}

\separateur

\begin{boxresume}
\begin{itemize}
    \item \textbf{Un processus orienté vers des objectifs.} Le contrôle
          interne donne une assurance raisonnable sur les opérations,
          l'information et la conformité.
    \item \textbf{Cinq composantes, dix-sept principes.} Environnement,
          risques, activités de contrôle, information et communication,
          pilotage~; chacune doit être en place, fonctionner, et
          fonctionner avec les autres.
    \item \textbf{Chaque principe a une traduction informatique}, et une
          pièce qui la prouve.
    \item \textbf{Préventif, détectif, correctif~; manuel, automatisé,
          semi-automatisé.} Un dispositif solide combine les modes.
    \item \textbf{Les contrôles généraux d'abord.} Sans maîtrise des accès
          et des changements, aucun contrôle d'application n'est fiable.
    \item \textbf{La séparation se lit dans les profils.} Elle n'existe que
          si l'application l'impose~; sinon, il faut un contrôle
          compensatoire prouvé.
    \item \textbf{Une mission écrite du CONSUPE.} Le décret n°~2013/287
          (art.~28) le charge d'évaluer les structures de contrôle
          interne.
\end{itemize}
\end{boxresume}

\begin{boxexo}[title={\ding{45}\quad Exercice~2.1. Classer des contrôles \niveaubase}]
Classez chaque contrôle selon son moment et son mode.
\begin{enumerate}[(a)]
    \item SIDÉPENSE refuse un engagement supérieur au crédit disponible.
    \item Le chef de la liquidation signe chaque soir la liste des mandats
          du jour, éditée par l'application.
    \item Le poste comptable rejette un mandat dont la facture scannée est
          illisible.
    \item La DAF émet un ordre de reversement après un paiement indu.
\end{enumerate}
\end{boxexo}

\begin{boxexo}[title={\ding{45}\quad Exercice~2.2. Composantes et principes \niveaubase}]
Rattachez chacun des manques suivants, relevés au MINÉQUIR, à une
composante et à un principe du COSO~: (a) aucune cartographie des risques
de SIDÉPENSE~; (b) aucune procédure de désactivation des comptes lors des
mutations~; (c) les anomalies de fin de gestion ne remontent jamais au
directeur~; (d) la cellule informatique n'a aucune fiche de poste.
\end{boxexo}

\begin{boxexo}[title={\ding{45}\quad Exercice~2.3. Une matrice des profils \niveaumoyen}]
La DAF compte six utilisateurs de SIDÉPENSE~: ABENA (ENG), BELLO (LIQ,
ORD), CHOUPO (LIQ), DJOMO (ORD, FRS), ESSOMBA (FRS), FOUDA (ADM).
\begin{enumerate}[(a)]
    \item Identifiez les comptes en conflit au regard du Tableau~2.6.
    \item Proposez une nouvelle répartition sans conflit, sans créer de
          compte.
    \item Si une répartition sans conflit est impossible, décrivez le
          contrôle compensatoire et la preuve qu'il doit laisser.
\end{enumerate}
\end{boxexo}

\begin{boxexo}[title={\ding{45}\quad Exercice~2.4. Les contrôles généraux d'abord \niveauexpert}]
Le chef de la cellule informatique vous montre que SIDÉPENSE bloque
désormais les doublons (même fournisseur, même montant, même facture).
Dites quels contrôles généraux vous devez tester avant de conclure que ce
blocage est efficace pour toute l'année 2026, quelle pièce vous demandez
pour chacun, et ce que vous concluez si l'une de ces pièces manque.
\end{boxexo}

\begin{boxinfo}
\textbf{Transition.} Le contrôle interne est maintenant défini, composante
par composante. Reste la question de l'auditeur~: comment savoir s'il est
efficace~? Le Chapitre~\ref{chap:eval-sci} en donne la démarche, du
questionnaire à la conclusion.
\end{boxinfo}
}{}
        \IfFileExists{partie1/P01_C03_evaluer_controle_interne.tex}{% ============================================================================
% Fichier  : partie1/P01_C03_evaluer_controle_interne.tex
% Titre    : Évaluer le Système de Contrôle Interne
% Auteur   : MINKA MI NGUIDJOÏ Thierry Emmanuel
% Version  : 0.1 -- novembre 2026
% ============================================================================
% Chapitre central du cours (décision de l'auteur, 07/11/2026).
% Sources : COSO 2013 (IFACI), chap. 3, §§ 138-140 (déficience, déficience
% majeure) ; INF3226 C03 (tailles d'échantillon, NON / NEK) ; décret
% n° 2013/287, art. 28.
% Faits MINÉQUIR établis ici (à garder) :
%  - contrôle 5 du P00 : pièces du poste comptable reçues le 2 avril 2026,
%    le contrôle des pièces jointes a fonctionné -> OUI ; lecture finale
%    3 OUI, 5 NON, 0 NEK ;
%  - 31 mutations ou départs en 2025 ; 9 comptes encore actifs au 23/03/2026 ;
%  - déficiences : A (accès et séparation, contrôles 1 et 2) majeure ;
%    B (fichier des fournisseurs et doublons, contrôles 3 et 4) ;
%    C (revue de fin de gestion, contrôle 7).
% ============================================================================

\chapter{Évaluer le Système de Contrôle Interne}
\label{chap:eval-sci}
\chaptermark{Évaluer le contrôle interne}

\epigraph{%
    « Un contrôle mal conçu est inefficace, même exécuté à la lettre~; un
    contrôle bien conçu mais non exécuté ne protège de rien. »%
}{Règle du cours}

\niveauChapitre{Intermédiaire}{Chapitre~\ref{chap:sci}.}

\begin{boxobjectifs}
\begin{itemize}
    \item Conduire l'évaluation d'un système de contrôle interne en cinq
          étapes~: décrire, questionner, cheminer, évaluer la conception,
          tester le fonctionnement.
    \item Construire une matrice des risques et des contrôles pour un
          circuit financier informatisé.
    \item Utiliser un questionnaire de contrôle interne, en distinguant la
          réponse de l'audité de la preuve.
    \item Réaliser un test de cheminement et en tirer ce qu'il établit, et
          ce qu'il n'établit pas.
    \item Dimensionner un test de fonctionnement et coter le contrôle
          \OUI{}, \NON{} ou \NEK{}.
    \item Classer les déficiences (déficience, déficience majeure) selon le
          COSO~2013.
    \item Formuler une conclusion sur l'efficacité du système de contrôle
          interne.
\end{itemize}
\end{boxobjectifs}

\bigskip

% ============================================================================
\section*{Pourquoi ce chapitre}
% ============================================================================

Dans le prologue, la vérificatrice a coté huit contrôles. Elle savait
quels contrôles attendre, comment les tester, et ce que voulait dire un
\NON{}. Ce chapitre donne cette démarche, pas à pas. Il est le plus
important du cours~: chaque chapitre de la Partie~IV l'applique à un
domaine, et le cas intégrateur la reprend en entier.

La question directrice est celle de l'article~28 du décret
n°~2013/287~: évaluer les structures de contrôle interne. Elle se
décompose en deux questions plus précises. \emph{Le dispositif est-il
conçu pour réduire les risques~?} \emph{A-t-il fonctionné pendant toute
la période~?}

\separateur

% ============================================================================
\section{La démarche en cinq étapes}
\label{sec:c03-demarche}
% ============================================================================

\begin{tableaucours}{cL{2.6cm}L{4.8cm}L{4.4cm}}{%
    \tableauHeader{N°} &
    \tableauHeader{Étape} &
    \tableauHeader{Question} &
    \tableauHeader{Production}
}
1 & Décrire & Comment le circuit fonctionne-t-il, et où sont les
risques~? & Narratif, diagramme, matrice des risques et des contrôles \\
2 & Questionner & Quels contrôles l'entité dit-elle exercer~? &
Questionnaire de contrôle interne rempli \\
3 & Cheminer & Le circuit fonctionne-t-il comme on me l'a décrit~? &
Fiche de test de cheminement \\
4 & Évaluer la conception & Ces contrôles, s'ils fonctionnent,
réduisent-ils les risques~? & Cotation de la conception \\
5 & Tester le fonctionnement & Ont-ils fonctionné toute la période~? &
Feuilles de test, cotation finale \\
\end{tableaucours}
\vspace{-0.8em}
\begin{center}{\small\itshape Tableau~3.1. Les cinq étapes de l'évaluation
du contrôle interne}\end{center}

Les étapes s'enchaînent, et chacune peut arrêter la suivante. Un contrôle
dont la conception est défaillante n'a pas à être testé dans son
fonctionnement~: il est déjà une déficience. Un contrôle que l'entité ne
déclare pas, mais que la matrice des risques exige, est un \NON{} dès
l'étape~2. Les étapes~6 et~7, classer les déficiences et conclure, sont
traitées aux sections~\ref{sec:c03-deficiences} et~\ref{sec:c03-conclure}.

\separateur

% ============================================================================
\section{Décrire le système}
\label{sec:c03-decrire}
% ============================================================================

\subsection{Le narratif et le diagramme}

L'auditeur commence par décrire le circuit, avec les mots de l'entité~:
qui fait quoi, dans quel ordre, avec quel écran, et quelle trace reste.
Le narratif s'accompagne d'un diagramme, qui fait apparaître les étapes,
les acteurs et les points de contrôle.

\begin{center}
\begin{tikzpicture}[
    etape/.style = {draw=CouleurBordDef, fill=CouleurDef, rounded corners=3pt,
        text width=2.5cm, minimum height=0.9cm, align=center,
        font=\scriptsize\bfseries\color{CouleurPrimaire}},
    ctrl/.style = {draw=CouleurAccent, fill=CouleurAccent!15, circle,
        inner sep=1pt, font=\scriptsize\bfseries},
    arrow/.style = {-{Stealth[length=4pt]}, thick, color=CouleurSecondaire}
]
\node[etape] (e) at (0,0) {Engagement\\(ENG)};
\node[etape] (l) at (3.4,0) {Liquidation\\(LIQ)};
\node[etape] (o) at (6.8,0) {Ordonnancement\\(ORD)};
\node[etape] (p) at (10.2,0) {Paiement\\(poste comptable)};
\node[etape, text width=3.4cm] (f) at (5.1,-2.0) {Fichier des fournisseurs\\(FRS)};
\draw[arrow] (e) -- (l);
\draw[arrow] (l) -- (o);
\draw[arrow] (o) -- (p);
\draw[arrow, dashed] (f) -- (l);
\draw[arrow, dashed] (f) -- (o);
\node[ctrl] at (1.7,0.75) {C1};
\node[ctrl] at (5.1,0.75) {C2};
\node[ctrl] at (8.5,0.75) {C3};
\node[ctrl] at (7.4,-2.0) {C4};
\end{tikzpicture}
\end{center}
\vspace{-0.4em}
\begin{center}{\small\itshape Figure~3.1. Le circuit de la dépense dans
SIDÉPENSE et ses points de contrôle (C1 contrôle du crédit~; C2 contrôle
de la liquidation~; C3 contrôle des doublons et des pièces~; C4 contrôle
des modifications de RIB)}\end{center}

\subsection{La matrice des risques et des contrôles}

La description aboutit à une \termecle{matrice des risques et des
contrôles}\index{matrice des risques et des contrôles}~: pour chaque
risque du circuit, le ou les contrôles qui doivent le couvrir. C'est la
liste de ce que l'auditeur s'attend à trouver, établie \emph{avant}
d'interroger l'entité, à partir des textes, des référentiels et de son
expérience.

\begin{tableaucours}{L{4.0cm}L{4.6cm}L{1.6cm}L{1.8cm}}{%
    \tableauHeader{Risque} &
    \tableauHeader{Contrôle attendu} &
    \tableauHeader{Type} &
    \tableauHeader{Principe}
}
R1. Un agent parti agit encore & Désactivation du compte à la date de la
mutation & Préventif & 11 \\
R2. Un agent émet seul un mandat & Profils LIQ et ORD jamais réunis &
Préventif & 10, 11 \\
R3. Le bénéficiaire est changé & Seconde validation de tout RIB modifié &
Préventif & 10 \\
R4. Une facture est payée deux fois & Blocage des doublons dans
l'application & Préventif & 10 \\
R5. Le mandat n'est pas justifié & Vérification des pièces jointes par le
poste comptable & Préventif & 10 \\
R6. Une opération ne laisse pas de trace & Journalisation des opérations
sensibles & Détectif & 13 \\
R7. Une anomalie de fin de gestion passe & Revue des mandats émis hors
heures et des RIB modifiés & Détectif & 16 \\
R8. Les traces disparaissent & Conservation protégée des journaux &
Détectif & 13 \\
\end{tableaucours}
\vspace{-0.8em}
\begin{center}{\small\itshape Tableau~3.2. Matrice des risques et des
contrôles du circuit de la dépense (MINÉQUIR)}\end{center}

\begin{boxalert}
\textbf{Violation typique.} Construire la liste des contrôles à partir de
ce que l'entité déclare faire. On ne trouve alors que des contrôles qui
existent~; ceux qui manquent n'apparaissent jamais.

\smallskip
\textbf{Règle~:} la matrice part des risques, pas des contrôles déclarés.
Un risque sans contrôle en face est déjà une constatation.
\end{boxalert}

\separateur

% ============================================================================
\section{Questionner : le questionnaire de contrôle interne}
\label{sec:c03-qci}
% ============================================================================

\begin{boxdef}
\termecle{Questionnaire de contrôle interne}\index{questionnaire de contrôle
interne} (\sigle{QCI})~: liste de questions fermées, tirée de la matrice des
risques et des contrôles, qui demande à l'entité si chaque contrôle attendu
existe, qui l'exerce, à quelle fréquence, et quelle trace il laisse.
\end{boxdef}

Une question bien formée appelle une réponse vérifiable. «~Les droits
sont-ils bien gérés~?~» n'appelle qu'un «~oui~». «~Qui désactive le compte
d'un agent muté, dans quel délai, et sur quelle pièce~?~» appelle une
réponse que l'on peut tester.

\begin{tableaucours}{cL{5.4cm}L{3.0cm}L{3.2cm}}{%
    \tableauHeader{N°} &
    \tableauHeader{Question} &
    \tableauHeader{Réponse déclarée} &
    \tableauHeader{Pièce demandée}
}
1 & Qui informe la cellule informatique d'une mutation, et sous quel
délai~? & «~Le service du personnel, en principe~» & Procédure, exemples
de notifications \\
2 & Qui désactive le compte, et quelle trace reste~? & «~La cellule, dès
qu'elle est informée~» & Journal des désactivations 2025 \\
3 & Les profils LIQ et ORD peuvent-ils être réunis sur un compte~? &
«~Non~» & Matrice des profils \\
4 & Qui peut modifier un RIB, et avec quelle validation~? & «~Le service
de la liquidation~» & Liste des détenteurs du profil FRS, journal des
modifications \\
5 & L'application bloque-t-elle les doublons~? & «~Oui, le numéro de
facture est unique~» & Paramétrage, essai \\
6 & Une revue des mandats de fin de gestion est-elle faite~? & «~Le DAF
suit les gros montants~» & Comptes rendus de revue \\
\end{tableaucours}
\vspace{-0.8em}
\begin{center}{\small\itshape Tableau~3.3. Extrait du QCI sur le circuit
de la dépense, rempli par la DAF et la cellule informatique}\end{center}

\begin{boxinfo}
\textbf{Une réponse au QCI est une déclaration, pas une preuve.} Elle dit
ce que l'entité croit faire. Le cheminement et les tests disent ce qu'elle
fait. La réponse à la question~5 en est l'exemple~: le numéro de facture
est bien unique, mais le suffixe «~bis~» suffit à en faire un autre
numéro. Le contrôle déclaré existe~; il ne couvre pas le risque.
\end{boxinfo}

\separateur

% ============================================================================
\section{Cheminer : le test de cheminement}
\label{sec:c03-cheminement}
% ============================================================================

\begin{boxdef}
\termecle{Test de cheminement}\index{test de cheminement}
(\textit{walkthrough})~: suivi d'une opération réelle, du début à la fin
du circuit, en examinant à chaque étape l'écran utilisé, la trace produite
et le contrôle exercé. Il établit que l'auditeur a compris le circuit et
que les contrôles existent tels qu'ils sont décrits~; il n'établit pas
qu'ils ont fonctionné pendant toute la période.
\end{boxdef}

Mme~NGO~BIYONG choisit le mandat 2025-08817, le premier paiement du
décompte n°~3, et le suit dans SIDÉPENSE, assise à côté de l'agent qui
l'a traité.

\begin{tableaucours}{L{2.4cm}L{5.4cm}L{4.4cm}}{%
    \tableauHeader{Étape} &
    \tableauHeader{Ce que la vérificatrice observe} &
    \tableauHeader{Trace obtenue}
}
Engagement & Le crédit est contrôlé~: un essai à un montant supérieur est
refusé & Capture de l'écran de refus \\
Liquidation & L'agent joint la facture et le PV de réception, puis valide
& Historique du mandat \\
Ordonnancement & Un autre agent valide~; l'application ne vérifie pas
qu'il est différent du liquidateur & Historique (deux noms différents,
par choix et non par contrainte) \\
Transmission & Le mandat part au poste comptable avec ses pièces
scannées & Bordereau électronique \\
Paiement & Le poste comptable vérifie les pièces, puis paie & Visa
électronique, relevé \\
\end{tableaucours}
\vspace{-0.8em}
\begin{center}{\small\itshape Tableau~3.4. Test de cheminement du mandat
2025-08817}\end{center}

\begin{boxscenario}{Dossier MINÉQUIR. Ce que montre le cheminement}
\textbf{Constatation.} Sur le mandat 2025-08817, liquidation et
ordonnancement ont été faits par deux agents différents. L'application
n'empêche pas qu'un même agent fasse les deux.

\medskip
\textbf{Analyse} (confiance élevée). La séparation a été respectée sur ce
mandat, mais par habitude, pas par contrainte. Un cheminement réussi ne
prouve donc pas que le contrôle existe~: il montre ici que la séparation
dépend de la bonne volonté des agents. C'est un défaut de
\emph{conception}, que le cheminement suffit à établir.
\end{boxscenario}

\separateur

% ============================================================================
\section{Évaluer la conception, tester le fonctionnement}
\label{sec:c03-tester}
% ============================================================================

\subsection{Deux questions distinctes}

\begin{tableaucours}{L{3.0cm}L{4.8cm}L{4.4cm}}{%
    \tableauHeader{Évaluation} &
    \tableauHeader{Question} &
    \tableauHeader{Comment}
}
Conception & Le contrôle, s'il fonctionne, réduit-il le risque~? &
Lecture de la procédure, du paramétrage, cheminement \\
Fonctionnement & Le contrôle a-t-il fonctionné comme prévu pendant toute
la période~? & Test d'un échantillon d'exécutions et de leurs traces,
ou test du paramétrage \\
\end{tableaucours}
\vspace{-0.8em}
\begin{center}{\small\itshape Tableau~3.5. Conception et fonctionnement
d'un contrôle}\end{center}

\subsection{Combien d'exécutions tester~?}

Le fonctionnement d'un contrôle manuel se teste sur un échantillon dont la
taille dépend de la fréquence du contrôle. Le cours retient les ordres de
grandeur du Tableau~3.6~; ce sont des conventions de travail, que la
méthodologie de chaque institution peut fixer autrement.

\begin{tableaucours}{L{3.6cm}L{3.2cm}L{5.4cm}}{%
    \tableauHeader{Fréquence du contrôle} &
    \tableauHeader{Exécutions par an} &
    \tableauHeader{Échantillon indicatif}
}
Annuelle & 1 & 1 \\
Trimestrielle & 4 & 2 \\
Mensuelle & 12 & 2 à 5 \\
Hebdomadaire & 52 & 5 à 15 \\
Quotidienne & environ 250 & 20 à 40 \\
Plusieurs fois par jour & plus de 250 & 25 à 60 \\
Automatisée & & Un essai qui doit passer et un qui doit être bloqué, si
les contrôles généraux sont efficaces \\
\end{tableaucours}
\vspace{-0.8em}
\begin{center}{\small\itshape Tableau~3.6. Tailles d'échantillon
indicatives (convention du cours)}\end{center}

\subsection{Coter}

\begin{itemize}
    \item \textbf{\OUI{}}~: le contrôle est bien conçu et a fonctionné sur
          tout l'échantillon (ou l'essai automatisé a donné le résultat
          attendu).
    \item \textbf{\NON{}}~: la conception est défaillante, ou une
          exécution au moins a manqué, ou la trace que le contrôle devait
          produire n'existe pas.
    \item \textbf{\NEK{}}~: malgré des diligences normales, l'auditeur ne
          peut pas conclure~: pièce non remise, sources contradictoires
          qu'il ne peut pas départager. Il écrit les diligences faites et
          relance.
\end{itemize}

\begin{boxobjet}{C-ACC-01 · Comptes des agents mutés}
\textbf{Objectif.} Aucun agent muté ou parti ne peut agir dans
SIDÉPENSE.

\textbf{Risque couvert.} R1~: un compte resté actif est utilisé par son
ancien titulaire ou par un tiers.

\textbf{Critère.} Le compte est désactivé au plus tard à la date d'effet
de la mutation (procédure à écrire~; à défaut, principe~11 du COSO).

\textbf{Test de conception.} Existe-t-il une notification du service du
personnel à la cellule informatique~? Réponse~: non, aucune procédure
écrite.

\textbf{Test de fonctionnement.} Rapprocher la liste des 31 mutations et
départs de 2025 de la liste des comptes actifs au 23~mars~2026.

\textbf{Résultat.} 9 comptes sur 31 sont encore actifs, dont
\texttt{mbarga.j}.

\textbf{Cote.} \NON{} (conception et fonctionnement défaillants).
\end{boxobjet}

\begin{boxalert}
\textbf{Violation typique.} Coter \NEK{} un contrôle dont la procédure
exige un procès-verbal, alors que l'entité a confirmé qu'il n'en existe
aucun.

\smallskip
\textbf{Règle~:} l'absence d'une trace exigée par le critère est un \NON{}
établi. \NEK{} est réservé au cas où l'auditeur ne peut pas savoir.
\end{boxalert}

\begin{boxscenario}{Dossier MINÉQUIR. Le contrôle 5 est levé}
\textbf{Constatation.} Le 2~avril~2026, le poste comptable remet les
visas électroniques des mandats 2025-08817 et 2025-09342. Les deux
portent la vérification des pièces jointes~; les pièces étaient
complètes et lisibles.

\medskip
\textbf{Analyse} (confiance élevée). Le contrôle des pièces a fonctionné
tel qu'il est conçu~: il vérifie la présence et la lisibilité des pièces,
pas leur unicité. Le contrôle~5 passe de \NEK{} à \OUI{}. Le double
paiement relève d'un autre contrôle, le blocage des doublons (contrôle~4),
coté \NON{}.
\end{boxscenario}

\separateur

% ============================================================================
\section{Classer les déficiences}
\label{sec:c03-deficiences}
% ============================================================================

\subsection{Deux niveaux}

Le COSO~2013 distingue deux niveaux de déficience~\cite{coso2013icfr}.

\begin{boxdef}
\termecle{Déficience du contrôle interne}\index{déficience}~: insuffisance
d'une ou plusieurs composantes ou d'un ou plusieurs principes, qui réduit
la probabilité que l'entité atteigne ses objectifs.

\medskip
\termecle{Déficience majeure}\index{déficience majeure}~: déficience, ou
combinaison de déficiences, qui réduit \emph{fortement} cette probabilité.
\end{boxdef}
{\small\itshape Définitions reformulées d'après le COSO~2013, §~138.}

\medskip

Deux règles du COSO (§§~139 et~140) commandent la conclusion. Lorsqu'une
déficience majeure existe, l'entité ne peut pas conclure que son système
de contrôle interne est efficace. Et une déficience majeure d'une
composante, ou d'un principe, ne peut pas être compensée par le bon
fonctionnement d'une autre composante, ou d'autres principes.

\subsection{Une grille de classement}

Pour décider si une déficience est majeure, le cours retient quatre
questions. Elles ne remplacent pas le jugement~; elles l'obligent à
s'expliquer.

\begin{tableaucours}{cL{6.4cm}L{5.4cm}}{%
    \tableauHeader{N°} &
    \tableauHeader{Question} &
    \tableauHeader{Penche vers «~majeure~» si…}
}
1 & Quelle perte, ou quelle irrégularité, la déficience rend-elle
possible~? & Elle porte sur des montants importants ou sur tout le
circuit \\
2 & Avec quelle probabilité~? & Rien n'empêche qu'elle se produise, ou
elle s'est déjà produite \\
3 & Un contrôle compensatoire la couvre-t-il, et est-il prouvé~? & Aucun,
ou il n'est pas prouvé \\
4 & Ouvre-t-elle une possibilité de fraude~? & Oui, en particulier par
contournement de la séparation des tâches \\
\end{tableaucours}
\vspace{-0.8em}
\begin{center}{\small\itshape Tableau~3.7. Grille de classement des
déficiences (convention du cours)}\end{center}

\subsection{Les déficiences du MINÉQUIR}

Les cinq \NON{} de la matrice ne font pas cinq déficiences~: plusieurs
contrôles manquants peuvent relever d'une même cause. La vérificatrice
les regroupe en trois déficiences.

\begin{tableaucours}{cL{4.6cm}L{1.8cm}L{1.8cm}L{3.0cm}}{%
    \tableauHeader{} &
    \tableauHeader{Déficience} &
    \tableauHeader{Contrôles} &
    \tableauHeader{Principes} &
    \tableauHeader{Classement}
}
A & Les accès ne sont pas maîtrisés~: comptes des partants actifs,
séparation LIQ / ORD non imposée & 1, 2 & 10, 11 & Majeure (questions 1
à 4) \\
B & Le fichier des fournisseurs n'est pas protégé et les doublons ne sont
pas bloqués & 3, 4 & 10 & Déficience \\
C & Aucune revue des opérations de fin de gestion & 7 & 16 & Déficience \\
\end{tableaucours}
\vspace{-0.8em}
\begin{center}{\small\itshape Tableau~3.8. Classement des déficiences du
circuit de la dépense}\end{center}

\begin{boxinfo}
\textbf{Pourquoi B n'est pas majeure à elle seule.} Si les accès étaient
maîtrisés, une modification de RIB et un mandat en double exigeraient
deux personnes différentes, et la collusion resterait nécessaire. C'est
la déficience~A qui rend B dangereuse. Le classement dit cette relation~:
la combinaison A et B est majeure, et c'est elle qui a permis le double
paiement.
\end{boxinfo}

\separateur

% ============================================================================
\section{Conclure sur l'efficacité du système}
\label{sec:c03-conclure}
% ============================================================================

\subsection{La lecture de conformité}

\begin{boxlecture}{Circuit de la dépense, après les tests}
\textbf{Constatation.} Sur huit contrôles attendus, trois existent et
fonctionnent (contrôles~5, 6 et~8), cinq manquent ou n'ont pas fonctionné.
Aucun contrôle ne reste à établir.

\medskip
\textbf{Analyse} (confiance élevée). Les contrôles qui fonctionnent sont
la vérification des pièces et deux contrôles de détection. La déficience
majeure~A porte sur les accès, dont dépendent tous les autres contrôles.

\medskip
\hfill\conformite{3}{8}{0}
\end{boxlecture}

\subsection{Une échelle simple}

Pour situer un domaine et suivre son évolution d'une mission à l'autre,
le cours propose une échelle à cinq niveaux. C'est une convention
pédagogique, pas une norme~: elle aide à dire où l'on en est, elle ne
remplace pas la conclusion.

\begin{tableaucours}{cL{2.6cm}L{9.2cm}}{%
    \tableauHeader{Niveau} &
    \tableauHeader{Nom} &
    \tableauHeader{Description}
}
0 & Inexistant & Le risque n'est pas identifié~; aucun contrôle \\
1 & Informel & Des contrôles existent, par habitude, sans procédure ni
trace \\
2 & Défini & Les contrôles sont écrits, mais leur exécution n'est pas
prouvée \\
3 & Appliqué & Les contrôles sont écrits, exécutés et tracés~; les tests
de fonctionnement sont concluants \\
4 & Piloté & Les contrôles sont surveillés, les déficiences suivies et
corrigées \\
\end{tableaucours}
\vspace{-0.8em}
\begin{center}{\small\itshape Tableau~3.9. Échelle de maturité du
contrôle interne (convention du cours)}\end{center}

Le circuit de la dépense du MINÉQUIR est au niveau~1~: la séparation des
tâches existe par habitude (cheminement), sans être imposée ni écrite.

\subsection{Formuler la conclusion}

La conclusion porte sur un objectif, une période et un périmètre, et dit
sur quoi elle repose. Celle de Mme~NGO~BIYONG s'écrit ainsi~:

\begin{quote}\itshape
Sur la base des travaux décrits au rapport (matrice des risques et des
contrôles, questionnaire, test de cheminement, tests de fonctionnement
sur l'exercice 2025), le système de contrôle interne du circuit de la
dépense de la DAF du MINÉQUIR n'est pas efficace~: il présente une
déficience majeure dans la maîtrise des accès à SIDÉPENSE, qui a permis
un double paiement de 187\,450\,000~FCFA.
\end{quote}

\begin{boxprocedure}
\textbf{Le dossier d'évaluation d'un circuit} contient, dans cet ordre~:
\begin{enumerate}
    \item le narratif et le diagramme, validés par l'entité~;
    \item la matrice des risques et des contrôles, datée~;
    \item le QCI rempli, avec le nom et la fonction de chaque répondant~;
    \item la fiche du test de cheminement et ses traces~;
    \item une feuille de test par contrôle~: conception, échantillon,
          résultat, cote, pièces~;
    \item le tableau de classement des déficiences, avec la réponse aux
          quatre questions pour chaque déficience majeure~;
    \item la conclusion, et la référence des pièces sur lesquelles elle
          repose.
\end{enumerate}
\end{boxprocedure}

\begin{boxmetier}{Vérificateur}
L'article~28 du décret n°~2013/287 vous charge d'évaluer les structures
de contrôle interne. Joignez à votre rapport la matrice des risques et des
contrôles~: elle montre que votre conclusion repose sur une liste établie
à l'avance, et non sur ce que l'entité a bien voulu montrer.
\end{boxmetier}

\begin{boxmetier}{Inspecteur}
Une inspection déclenchée par une alerte peut sauter l'étape du QCI, mais
pas celle du cheminement~: suivez au moins une opération réelle dans
l'application avant de conclure sur un circuit.
\end{boxmetier}

\begin{boxmetier}{Auditeur interne}
Votre évaluation annuelle du contrôle interne sera relue par le CONSUPE.
Tenez une matrice des risques et des contrôles par circuit, et mettez-la
à jour à chaque changement d'application~: c'est votre meilleure défense
devant un vérificateur.
\end{boxmetier}

\separateur

\begin{boxresume}
\begin{itemize}
    \item \textbf{Cinq étapes.} Décrire, questionner, cheminer, évaluer la
          conception, tester le fonctionnement~; puis classer et conclure.
    \item \textbf{La matrice part des risques}, pas des contrôles déclarés.
    \item \textbf{Le QCI recueille des déclarations}~; le cheminement et les
          tests produisent des preuves.
    \item \textbf{Conception d'abord.} Un contrôle mal conçu est une
          déficience, sans test de fonctionnement.
    \item \textbf{Trois cotes.} \OUI{}, \NON{}, \NEK{}~; l'absence d'une
          trace exigée est un \NON{}.
    \item \textbf{Deux niveaux de déficience.} Une déficience majeure
          interdit de conclure à l'efficacité, et ne se compense pas.
    \item \textbf{Une conclusion située.} Un objectif, une période, un
          périmètre, et les travaux sur lesquels elle repose.
\end{itemize}
\end{boxresume}

\begin{boxexo}[title={\ding{45}\quad Exercice~3.1. Formuler des questions \niveaubase}]
Réécrivez les questions suivantes du QCI pour qu'elles appellent une
réponse vérifiable~: (a) «~Les sauvegardes sont-elles faites~?~»~; (b)
«~Les mots de passe sont-ils sûrs~?~»~; (c) «~Le fichier de paie est-il
contrôlé~?~».
\end{boxexo}

\begin{boxexo}[title={\ding{45}\quad Exercice~3.2. Dimensionner et coter \niveaubase}]
Le chef de la liquidation signe chaque soir la liste des mandats du jour.
\begin{enumerate}[(a)]
    \item Combien de listes testez-vous pour l'exercice 2025~?
    \item Sur l'échantillon, deux listes ne sont pas signées. Quelle
          cote~?
    \item Les listes de novembre n'ont pas été retrouvées, et le service
          affirme qu'elles ont été signées. Quelle cote pour novembre, et
          que faites-vous~?
\end{enumerate}
\end{boxexo}

\begin{boxexo}[title={\ding{45}\quad Exercice~3.3. Une matrice pour la paie \niveaumoyen}]
Construisez la matrice des risques et des contrôles du circuit de la paie
de la CAPNO (entrée d'un agent, modification d'un élément de solde,
départ, calcul, virement), avec au moins six risques, le contrôle attendu
pour chacun, son type et le principe du COSO correspondant.
\end{boxexo}

\begin{boxexo}[title={\ding{45}\quad Exercice~3.4. Classer et conclure \niveauexpert}]
À la CAPNO, l'auditeur interne relève~: (1) 37 agents payés ne figurent
pas au fichier du personnel~; (2) le profil qui crée un agent peut aussi
valider la paie~; (3) aucun rapprochement entre la paie et le fichier du
personnel n'est fait~; (4) les sauvegardes ne sont pas testées.
\begin{enumerate}[(a)]
    \item Regroupez ces constats en déficiences, et classez chacune à
          l'aide du Tableau~3.7.
    \item Rédigez la conclusion sur l'efficacité du contrôle interne de la
          paie.
    \item Dites ce qui changerait votre classement.
\end{enumerate}
\end{boxexo}

\begin{boxinfo}
\textbf{Transition.} La Partie~I a donné l'objet et la méthode~:
évaluer le contrôle interne d'un système d'information. Reste à savoir au
nom de quoi l'on évalue. Le Chapitre~\ref{chap:institutions} présente les
institutions qui contrôlent les systèmes d'information publics au
Cameroun, et leurs compétences.
\end{boxinfo}
}{}

    % ======================================================================
    % PARTIE 2 — CADRE INSTITUTIONNEL ET NORMATIF (J03)
    % « Le mandat précède la méthode. »
    % Standards de référence : décret n° 2013/287 ; loi n° 2010/012 ;
    %   ISSAI 100, 400, 4000 ; GUID 5100 ; normes IIA 2017 ; COBIT 2019 ;
    %   ISO/IEC 27001 et 27002 ; loi n° 2024/017.
    % Niveau : fondamental
    % ======================================================================
    \part{Cadre Institutionnel et Normatif}
    \label{part:cadre}

        \IfFileExists{partie2/P02_C04_qui_controle_si_publics.tex}{% ============================================================================
% Fichier  : partie2/P02_C04_qui_controle_si_publics.tex
% Titre    : Qui Contrôle les SI Publics au Cameroun
% Auteur   : MINKA MI NGUIDJOÏ Thierry Emmanuel
% Version  : 0.1 -- novembre 2026
% ============================================================================
% Textes lus (dossier Normes) : décret n° 2013/287 (art. 1, 2, 3, 5, 27, 28) ;
% loi n° 2010/012 (art. 4 pt 9, art. 7 al. 2, art. 13 al. 1, art. 32).
% Chambre des Comptes, lois de 2018 : textes absents du dossier Normes ;
% aucun article cité, rôle décrit en termes généraux (décision de l'auteur :
% la Chambre des Comptes s'y retrouve « par extension »).
% Fait CAPNO établi ici : rapport d'audit interne de HAMADOU transmis au
% CONSUPE (art. 5) le 15 mai 2026.
% ============================================================================

\chapter{Qui Contrôle les SI Publics au Cameroun}
\label{chap:institutions}
\chaptermark{Qui contrôle les SI publics}

\epigraph{%
    « Le mandat précède la méthode. »%
}{Règle du cours}

\niveauChapitre{Fondamental}{Chapitre~\ref{chap:audit-public}.}

\begin{boxobjectifs}
\begin{itemize}
    \item Citer les dispositions du décret n°~2013/287 qui fondent le
          contrôle des systèmes d'information par le CONSUPE.
    \item Expliquer le rôle de l'ANTIC dans l'audit de sécurité des
          systèmes d'information publics, au titre de la loi n°~2010/012.
    \item Situer la Chambre des Comptes et expliquer pourquoi elle rencontre
          les systèmes d'information par extension.
    \item Situer les autres acteurs~: inspections générales, contrôle
          financier, audit interne, CDBF.
    \item Déterminer, pour une situation donnée, qui est compétent et
          comment les interventions s'articulent.
\end{itemize}
\end{boxobjectifs}

\bigskip

% ============================================================================
\section*{Pourquoi ce chapitre}
% ============================================================================

Dans le prologue, la vérificatrice a exigé une extraction de SIDÉPENSE
sans avoir à demander la permission. Elle le pouvait parce qu'un texte
lui en donne le droit. Ce chapitre présente ces textes. Il répond à trois
questions~: \emph{qui a compétence pour contrôler un système
d'information public~? Sur quoi porte chaque compétence~? Que devient un
constat d'un organe à l'autre~?}

Un principe de méthode gouverne le chapitre~: seuls sont cités entre
guillemets les textes lus dans leur version disponible. Les autres sont
décrits dans leurs grandes lignes, et l'auditeur doit les lire avant de
les invoquer.

\separateur

% ============================================================================
\section{Le Contrôle Supérieur de l'État}
\label{sec:c04-consupe}
% ============================================================================

\subsection{Une institution supérieure de contrôle}

Le Contrôle Supérieur de l'État est organisé par le décret n°~2013/287 du
4~septembre~2013~\cite{decret2013287consupe}. Ses services relèvent de
l'autorité directe du Président de la République, à qui ils rendent
compte, et sont dirigés par un ministre délégué à la Présidence (art.~1).

\begin{boxjuris}
\textbf{\loicam{Décret n°~2013/287}, article~2, alinéa~1}~: \textit{« Les
services du Contrôle supérieur de l'Etat constituent l'institution
supérieure de contrôle des finances publiques (ISC) du Cameroun. Ils sont
chargés de l'audit externe. »}

\medskip
\textbf{Article~2, alinéa~3}~: \textit{« Dans l'exercice de leurs
attributions, les services du Contrôle supérieur de l'Etat effectuent~:
un contrôle de conformité et de régularité~; un contrôle financier~; un
contrôle de performance~; un contrôle de l'environnement~; un contrôle des
systèmes d'information. »}

\medskip
\textbf{Article~2, alinéa~4}~: \textit{« Les contrôles sont effectués par
les vérificateurs du Contrôle supérieur de l'Etat sur pièces et sur place.
A cet effet, et nonobstant les contrôles effectués par les autres organes
de contrôle, les responsables des structures auditées sont tenus de
présenter les pièces relatives à la période auditée. »}

\medskip
\textbf{Portée pour l'audit.} Trois conséquences pratiques. Le contrôle
des systèmes d'information est une compétence \emph{explicite}, écrite à
côté du contrôle de régularité~: il n'a pas à être déduit d'une autre
mission. Il s'exerce sur place, donc dans l'application, et pas seulement
sur les documents qui en sortent. Enfin, l'obligation de présenter les
pièces vaut «~nonobstant~» les autres contrôles~: une entité ne peut pas
refuser une extraction au motif que l'ANTIC ou une inspection l'a déjà
auditée.
\end{boxjuris}

\subsection{Le champ du contrôle}

L'article~2, alinéa~1, fixe un champ large. Il couvre notamment les
services publics, les établissements publics, les collectivités
territoriales décentralisées, les entreprises publiques et parapubliques,
les projets à financements extérieurs, et les organismes qui bénéficient
de concours financiers, d'avals ou de garanties de l'État. Partout où
l'une de ces entités exploite un système d'information, le contrôle des
SI s'applique.

L'organisation des services en tire les conséquences. L'article~27,
alinéa~2, crée au sein de la Division des inspections et de contrôle des
administrations publiques une \emph{section des administrations
financières, économiques, sociales et des systèmes d'informations}.
L'article~28 charge chaque section du suivi et de l'évaluation des
activités des structures de contrôle interne
(Chapitre~\ref{chap:sci}).

\subsection{Les suites : le CDBF}

\begin{boxjuris}
\textbf{\loicam{Décret n°~2013/287}, article~3, alinéa~1}~: \textit{« Les
services du Contrôle supérieur de l'Etat concourent à la sanction des
ordonnateurs et des gestionnaires des deniers publics, dans les
conditions prévues par les lois et règlements en vigueur. A ce titre, le
ministre délégué en charge du Contrôle supérieur de l'Etat préside le
conseil de discipline budgétaire et financière, dont l'organisation et le
fonctionnement sont fixés par un texte particulier. »}

\medskip
\textbf{Portée pour l'audit.} Un constat du vérificateur peut fonder une
poursuite devant le Conseil de discipline budgétaire et financière
(\sigle{CDBF}). Il doit donc être établi avec la rigueur d'une pièce de
procédure~: critère cité, preuve conservée, contradictoire respecté
(Chapitre~\ref{chap:constater}).
\end{boxjuris}

\subsection{Le lien avec l'audit interne}

\begin{boxjuris}
\textbf{\loicam{Décret n°~2013/287}, article~5}~: \textit{« (1) Les
structures d'audit interne transmettent aux services du Contrôle supérieur
de l'Etat, à la diligence des chefs de départements ministériels
concernés, un exemplaire de chacun des rapports rédigés à la suite de
toutes les interventions menées au cours de l'exercice, ainsi qu'un
exemplaire du rapport annuel d'activités. (2) Après examen des rapports
visés à l'alinéa 1er ci-dessus, les services du Contrôle supérieur de
l'Etat peuvent, le cas échéant, saisir le conseil de discipline
budgétaire et financière. »}
\end{boxjuris}

\begin{boxscenario}{Dossier CAPNO. Un rapport qui voyage}
\textbf{Constatation.} Le 15~mai~2026, M.~HAMADOU, auditeur interne de la
CAPNO, remet son rapport sur la paie~: 37 agents payés ne figurent pas au
fichier du personnel. Conformément à l'article~5, le rapport est transmis
au CONSUPE.

\medskip
\textbf{Analyse} (confiance élevée). Le rapport de l'auditeur interne
devient une pièce que le CONSUPE peut examiner, et sur laquelle il peut
saisir le CDBF. Ses constats doivent donc être reproductibles par un
vérificateur~: extraction datée, empreinte, critère cité. Un rapport qui
écrit seulement «~des anomalies ont été relevées dans la paie~» ne pourra
pas être repris.
\end{boxscenario}

\separateur

% ============================================================================
\section{L'ANTIC : l'audit de sécurité}
\label{sec:c04-antic}
% ============================================================================

L'Agence nationale des technologies de l'information et de la
communication (\sigle{ANTIC}) est le régulateur de la sécurité
électronique. La loi n°~2010/012 relative à la cybersécurité et à la
cybercriminalité définit l'audit de sécurité et en confie la régulation à
l'Agence~\cite{loi2010012}.

\begin{boxjuris}
\textbf{\loicam{Loi n°~2010/012}, article~4, point~9}~: \textit{« Audit de
sécurité~: examen méthodique des composantes et des acteurs de la
sécurité, de la politique, des mesures, des solutions, des procédures et
des moyens mis en oeuvre par une organisation, pour sécuriser son
environnement, effectuer des contrôles de conformité, des contrôles
d'évaluation de l'adéquation des moyens (organisationnels, techniques,
humains, financiers) investis au regard des risques encourus,
d'optimisation, de rationalité et de performance. »}

\medskip
\textbf{Article~7, alinéa~2} (extrait)~: l'Agence est notamment chargée
\textit{« de s'assurer de la régularité, de l'effectivité des audits de
sécurité des systèmes d'information suivant les normes en la matière, des
organismes publics et des autorités de certification »}.

\medskip
\textbf{Article~32}~: \textit{« (1) Les réseaux de communications
électroniques et les systèmes d'information sont soumis à un régime
d'audit de sécurité obligatoire et périodique de leurs systèmes de
sécurité par l'Agence. (2) L'audit de sécurité et les mesures d'impact de
gravité sont effectués chaque année ou lorsque les circonstances
l'exigent. (3) Les rapports d'audit sont confidentiels et adressés au
Ministre chargé des Télécommunications. »}

\medskip
\textbf{Portée pour l'audit.} L'ANTIC porte sur la \emph{sécurité} des
systèmes~: protection, disponibilité, intégrité. Elle ne juge pas la
régularité des opérations financières qui y transitent. Pour le
vérificateur, l'existence et la date du dernier audit de sécurité sont
une pièce à demander~; son absence est un constat. Son contenu est
confidentiel~: l'auditeur demande au moins la date, le périmètre et le
plan d'action qui en est issu.
\end{boxjuris}

\separateur

% ============================================================================
\section{La Chambre des Comptes : par extension}
\label{sec:c04-chambre}
% ============================================================================

La Chambre des Comptes de la Cour suprême juge les comptes des comptables
publics et contrôle la gestion des entités publiques. Son texte
d'organisation ne figure pas dans le dossier des normes du cours~; il
n'est donc pas cité ici, et l'auditeur doit le lire avant de l'invoquer.

Son rapport aux systèmes d'information est indirect, mais réel. Les
comptes qu'elle juge sont produits par des applications~: un compte de
gestion édité par une application de comptabilité ne vaut que ce que
valent les contrôles de cette application. La Chambre rencontre donc les
SI \textbf{par extension}~: lorsqu'elle s'interroge sur la fiabilité d'un
compte, elle s'interroge sur le système qui l'a produit.

\begin{boxinfo}
\textbf{Ce que cela change pour le juge des comptes.} Un écart entre deux
états produits par la même application, ou entre l'application et les
relevés bancaires, n'est pas seulement une irrégularité comptable~: c'est
le symptôme d'un contrôle d'application ou d'un contrôle d'interface
défaillant (Chapitre~\ref{chap:applications}). Le juge peut demander
quel contrôle aurait dû l'empêcher, et qui en avait la charge.
\end{boxinfo}

\separateur

% ============================================================================
\section{Les autres acteurs}
\label{sec:c04-autres}
% ============================================================================

\begin{tableaucours}{L{3.2cm}L{3.8cm}L{5.2cm}}{%
    \tableauHeader{Acteur} &
    \tableauHeader{Nature du contrôle} &
    \tableauHeader{Rapport au système d'information}
}
Inspections générales (ministère des Finances, ministères techniques) &
Contrôle interne de l'administration, hiérarchique &
Contrôle sur pièces et sur place, souvent déclenché par un événement \\
Contrôle financier & Contrôle a priori de la dépense & Son visa passe de
plus en plus par l'application~: il est un profil \\
Audit interne des ministères et des entités & Évaluation du contrôle
interne pour la direction & Rapports transmis au CONSUPE (art.~5) \\
CDBF & Sanction des fautes de gestion & Juge des suites, sur la base des
constats \\
Chambre des Comptes & Juridictionnel, sur les comptes & Par extension,
via les comptes produits par les SI \\
ANTIC & Régulation de la sécurité électronique & Audit de sécurité des
SI des organismes publics \\
\end{tableaucours}
\vspace{-0.8em}
\begin{center}{\small\itshape Tableau~4.1. Les acteurs du contrôle des SI
publics (rôles généraux~; textes à lire avant citation, sauf le décret
n°~2013/287 et la loi n°~2010/012)}\end{center}

\separateur

% ============================================================================
\section{Qui fait quoi, et comment s'articuler}
\label{sec:c04-articuler}
% ============================================================================

\subsection{Trois angles sur un même système}

Pour un même système, trois questions différentes sont posées par trois
institutions.

\begin{tableaucours}{L{2.6cm}L{4.8cm}L{4.8cm}}{%
    \tableauHeader{Institution} &
    \tableauHeader{Question} &
    \tableauHeader{Exemple sur SIDÉPENSE}
}
CONSUPE & Le système et son contrôle interne garantissent-ils la
régularité, la fiabilité et la performance~? & La séparation entre
liquidation et ordonnancement est-elle imposée~? \\
ANTIC & Le système est-il sécurisé selon les normes~? & Les mots de
passe, les sauvegardes, la protection du réseau sont-ils conformes~? \\
Chambre des Comptes & Les comptes produits sont-ils réguliers et
sincères~? & Le compte de gestion reflète-t-il les mandats payés~? \\
\end{tableaucours}
\vspace{-0.8em}
\begin{center}{\small\itshape Tableau~4.2. Trois angles sur un même
système}\end{center}

\subsection{Règles d'articulation}

\begin{enumerate}
    \item \textbf{S'informer des autres interventions.} Le vérificateur
          demande, dès la prise de connaissance, les rapports d'audit
          interne (art.~5), la date du dernier audit de sécurité de
          l'ANTIC et les éventuelles procédures en cours devant la Chambre
          des Comptes ou le CDBF.
    \item \textbf{S'appuyer sans déléguer.} Un rapport d'un autre organe
          oriente les tests~; il ne remplace pas la preuve que le
          vérificateur obtient lui-même.
    \item \textbf{Rester dans son angle.} Le vérificateur qui constate une
          faille de sécurité la consigne et la signale~; il ne conduit pas
          un audit de sécurité à la place de l'ANTIC.
    \item \textbf{Préserver la suite.} Dès qu'un constat peut fonder une
          sanction ou une poursuite, la preuve doit être conservée selon
          les règles du Chapitre~\ref{chap:preuve}.
\end{enumerate}

\begin{boxalert}
\textbf{Violation typique.} Reprendre dans un rapport la conclusion d'un
audit de sécurité («~le SI est sécurisé~») comme preuve que les contrôles
d'accès de l'application sont efficaces.

\smallskip
\textbf{Règle~:} un audit de sécurité et une évaluation du contrôle
interne ne répondent pas à la même question. Chacun cite l'autre, aucun
ne remplace l'autre.
\end{boxalert}

\begin{boxmetier}{Vérificateur}
Citez toujours dans l'ordre de mission la disposition qui fonde votre
contrôle~: «~article~2, alinéa~3, du décret n°~2013/287~». Face à un
refus d'extraction, l'alinéa~4 du même article suffit.
\end{boxmetier}

\begin{boxmetier}{Juge des comptes}
Quand un compte vous paraît douteux, demandez la liste des applications
qui l'ont produit et le dernier rapport du CONSUPE ou de l'audit interne
qui les a examinées. Le défaut d'un contrôle d'application explique
souvent une irrégularité comptable.
\end{boxmetier}

\begin{boxmetier}{Contrôleur financier}
Votre visa électronique est une activité de contrôle que le CONSUPE
évaluera. Gardez la liste des personnes qui détiennent votre profil, et
signalez par écrit toute attribution que vous n'avez pas demandée.
\end{boxmetier}

\separateur

\begin{boxresume}
\begin{itemize}
    \item \textbf{Une compétence explicite.} Le décret n°~2013/287 (art.~2,
          al.~3) range le contrôle des SI parmi les cinq contrôles du
          CONSUPE, institution supérieure de contrôle chargée de l'audit
          externe.
    \item \textbf{Sur pièces et sur place.} Les pièces électroniques sont
          des pièces~; leur présentation est obligatoire, nonobstant les
          autres contrôles (al.~4).
    \item \textbf{Des suites.} Le CONSUPE préside le CDBF (art.~3) et reçoit
          les rapports des audits internes (art.~5).
    \item \textbf{L'ANTIC audite la sécurité.} Loi n°~2010/012~: audit de
          sécurité obligatoire, périodique, confidentiel.
    \item \textbf{La Chambre des Comptes, par extension.} Elle rencontre les
          SI à travers les comptes qu'ils produisent.
    \item \textbf{Trois angles, aucun ne remplace l'autre.} S'informer,
          s'appuyer sans déléguer, rester dans son angle, préserver la
          suite.
\end{itemize}
\end{boxresume}

\begin{boxexo}[title={\ding{45}\quad Exercice~4.1. Qui est compétent~? \niveaubase}]
Pour chaque situation, dites quelle institution est principalement
concernée, et sur quel texte vous vous fondez quand il est cité dans ce
chapitre.
\begin{enumerate}[(a)]
    \item Le réseau d'une recette des finances n'a jamais fait l'objet
          d'un audit de sécurité.
    \item Le compte de gestion d'un poste comptable diffère des paiements
          enregistrés dans l'application.
    \item Une DAF refuse de remettre l'extraction des mandats à un
          vérificateur, au motif qu'une inspection est en cours.
    \item L'auditeur interne d'un ministère achève un rapport sur la paie.
\end{enumerate}
\end{boxexo}

\begin{boxexo}[title={\ding{45}\quad Exercice~4.2. Lire le décret \niveaubase}]
Citez les deux dispositions du décret n°~2013/287 qui permettent à un
vérificateur (a) de conduire une mission dont l'objet est un système
d'information~; (b) d'évaluer le dispositif de contrôle interne d'un
ministère.
\end{boxexo}

\begin{boxexo}[title={\ding{45}\quad Exercice~4.3. Articuler \niveaumoyen}]
Au moment où le CONSUPE ouvre sa mission au MINÉQUIR, l'ANTIC annonce un
audit de sécurité de SIDÉPENSE. Rédigez la note de la vérificatrice à son
chef de division~: ce qu'elle demande à l'ANTIC, ce qu'elle ne fera pas
elle-même, et ce qu'elle fera pour préserver la preuve.
\end{boxexo}

\begin{boxexo}[title={\ding{45}\quad Exercice~4.4. Un rapport reproductible \niveauexpert}]
Réécrivez le constat suivant du rapport de M.~HAMADOU pour qu'un
vérificateur puisse le reprendre devant le CDBF~: «~Des anomalies
importantes ont été relevées dans la paie, notamment des agents qui ne
semblent pas exister.~» Indiquez le critère, la preuve, la méthode et ce
qui reste à établir.
\end{boxexo}

\begin{boxinfo}
\textbf{Transition.} Le mandat est posé. Reste à savoir selon quelles
règles professionnelles le conduire, et à quels critères comparer ce que
l'on trouve. Le Chapitre~\ref{chap:normes} présente les normes et les
référentiels de l'auditeur public.
\end{boxinfo}
}{}
        \IfFileExists{partie2/P02_C05_normes_referentiels.tex}{% ============================================================================
% Fichier  : partie2/P02_C05_normes_referentiels.tex
% Titre    : Normes et Référentiels de l'Auditeur Public
% Auteur   : MINKA MI NGUIDJOÏ Thierry Emmanuel
% Version  : 0.1 -- novembre 2026
% ============================================================================
% Sources : INF3226 C13 et C14 (resserrés) ; GUID 5100 et manuel IDI
% (dossier Normes) ; normes IIA 2017 (DOC-20230213, DOC-20330213) ; COBIT
% 2019 ; ISO/IEC 27001 et 27002 (2022) ; lois n° 2010/012 et 2024/017.
% Aucune norme ISO ou COBIT n'est reproduite (droits) : résumés seulement.
% ============================================================================

\chapter{Normes et Référentiels de l'Auditeur Public}
\label{chap:normes}
\chaptermark{Normes et référentiels}

\epigraph{%
    « La norme dit comment auditer~; le référentiel dit à quoi comparer. »%
}{Règle du cours}

\niveauChapitre{Fondamental}{Chapitre~\ref{chap:institutions}.}

\begin{boxobjectifs}
\begin{itemize}
    \item Distinguer une norme professionnelle (comment auditer) d'un
          référentiel de critères (à quoi comparer).
    \item Situer les normes de l'INTOSAI (ISSAI, GUID~5100) et le manuel de
          l'IDI dans le travail du vérificateur.
    \item Situer les normes de l'IIA dans le travail de l'auditeur interne.
    \item Utiliser COBIT~2019 et ISO/IEC~27001 et~27002 comme sources de
          critères, sans les confondre avec la loi.
    \item Choisir, pour un test donné, un critère explicite et le citer
          correctement.
    \item Énoncer les principes de déontologie communs aux auditeurs.
\end{itemize}
\end{boxobjectifs}

\bigskip

% ============================================================================
\section*{Pourquoi ce chapitre}
% ============================================================================

Le prologue l'a montré~: une cote \NON{} n'a de sens que si l'on sait à
quoi l'on compare. Ce chapitre sépare deux choses que l'on confond
souvent. Les \emph{normes professionnelles} disent comment conduire la
mission~; les \emph{référentiels} et les \emph{textes} fournissent les
critères. L'auditeur public a besoin des deux.

\separateur

% ============================================================================
\section{Norme, référentiel, critère}
\label{sec:c05-notions}
% ============================================================================

\begin{boxdef}
\termecle{Norme professionnelle}\index{norme professionnelle}~: règle qui
encadre la conduite d'une mission~: indépendance, planification, preuve,
documentation, rapport. Elle s'impose à l'auditeur.

\medskip
\termecle{Référentiel}\index{référentiel}~: ensemble structuré de bonnes
pratiques ou d'exigences (COBIT, ISO/IEC~27002) dans lequel l'auditeur
puise des critères. Il ne s'impose à l'entité que si un texte ou un
engagement le rend obligatoire.

\medskip
\termecle{Critère d'audit}\index{critère d'audit}~: exigence de référence,
choisie avant le test, à laquelle on compare la situation constatée.
\end{boxdef}

\begin{tableaucours}{L{3.0cm}L{4.4cm}L{4.8cm}}{%
    \tableauHeader{Source} &
    \tableauHeader{Force} &
    \tableauHeader{Exemple de critère}
}
Loi, décret & Obligatoire & Les SI des organismes publics font l'objet
d'un audit de sécurité (loi n°~2010/012, art.~32) \\
Règle interne & Obligatoire pour l'entité & Procédure de la DAF sur la
désactivation des comptes \\
Référentiel adopté & Obligatoire si l'entité l'a adopté & Politique de
sécurité fondée sur ISO/IEC~27002 \\
Bonne pratique & Recommandation & Un objectif de gestion COBIT~2019 \\
\end{tableaucours}
\vspace{-0.8em}
\begin{center}{\small\itshape Tableau~5.1. Hiérarchie des sources de
critères}\end{center}

\begin{boxalert}
\textbf{Violation typique.} Écrire «~non conforme à ISO/IEC~27002~» pour
une administration qui n'a jamais adopté cette norme.

\smallskip
\textbf{Règle~:} on n'est «~non conforme~» qu'à une règle qui s'impose.
Face à une simple bonne pratique, on écrit «~l'entité ne dispose pas de…,
alors que les bonnes pratiques (référence) recommandent…~», et l'on
classe l'écart en conséquence.
\end{boxalert}

\separateur

% ============================================================================
\section{Les normes de l'INTOSAI}
\label{sec:c05-intosai}
% ============================================================================

\subsection{Le cadre}

L'INTOSAI réunit les institutions supérieures de contrôle des finances
publiques, dont le CONSUPE fait partie au titre de l'article~2 du décret
n°~2013/287. Son cadre de prises de position professionnelles comprend
des principes, des normes (\sigle{ISSAI}) et des lignes directrices
(\sigle{GUID}).

\begin{tableaucours}{L{2.6cm}L{9.6cm}}{%
    \tableauHeader{Texte} &
    \tableauHeader{Objet}
}
ISSAI~100 & Principes fondamentaux de l'audit du secteur public~: les
trois types d'audit, les éléments de l'audit, les principes généraux \\
ISSAI~200, 300, 400 & Principes de l'audit financier, de performance et de
conformité \\
ISSAI~4000 & Norme de l'audit de conformité \\
GUID~5100 & Lignes directrices sur la vérification des systèmes
d'information \\
Manuel de l'IDI & Mise en œuvre pratique de l'audit de conformité selon les
ISSAI \\
\end{tableaucours}
\vspace{-0.8em}
\begin{center}{\small\itshape Tableau~5.2. Les textes de l'INTOSAI utiles à
l'auditeur des SI publics}\end{center}

\subsection{Le GUID 5100}

Le GUID~5100 complète les principes de l'ISSAI~100 pour les missions qui
portent sur les systèmes d'information~\cite{intosai2019guid5100}. Il
précise qu'un audit des SI peut être autonome ou s'insérer dans un audit
financier, de conformité ou de performance, et il énumère treize domaines
que la mission peut couvrir, parmi lesquels la politique et la
gouvernance des TI, le développement et l'acquisition des systèmes, la
gestion des opérations, la sécurité de l'information, la conformité
réglementaire, la continuité et la reprise après sinistre, et les
contrôles d'applications. Le cours suit la même logique~:
la Partie~IV reprend ces domaines.

\subsection{L'audit de conformité selon l'IDI}

Le manuel de l'IDI décrit l'audit de conformité comme une relation entre
trois parties~: l'auditeur, la partie responsable (l'entité auditée) et
les utilisateurs du rapport~\cite{idi2021conformite}. Il insiste sur deux
exigences que l'auditeur des SI applique directement~: identifier des
\emph{critères} appropriés avant le travail de terrain, et documenter
suffisamment la mission pour qu'un auditeur expérimenté, sans lien avec
elle, comprenne les travaux et les conclusions.

\separateur

% ============================================================================
\section{Les normes de l'audit interne}
\label{sec:c05-iia}
% ============================================================================

L'auditeur interne d'un ministère ou d'un établissement public suit les
normes de l'Institut des auditeurs internes (\sigle{IIA}). Le dossier des
normes contient l'édition 2017 du Cadre de référence international des
pratiques professionnelles~\cite{iia2017cripp}~; l'IIA l'a remplacée par
les normes mondiales d'audit interne publiées en
2024~\cite{iia2024gias}. Les exigences qui comptent ici sont stables
d'une édition à l'autre~: indépendance et objectivité, compétence,
planification fondée sur les risques, preuve suffisante, communication
des résultats et suivi des recommandations.

\begin{boxinfo}
\textbf{Vérificateur et auditeur interne parlent la même langue.} Les
normes de l'INTOSAI et celles de l'IIA convergent sur l'essentiel~:
critère, preuve, documentation, contradictoire. C'est ce qui permet au
CONSUPE de s'appuyer sur les rapports d'audit interne qu'il reçoit au
titre de l'article~5 du décret n°~2013/287.
\end{boxinfo}

\separateur

% ============================================================================
\section{Les référentiels de critères}
\label{sec:c05-referentiels}
% ============================================================================

\subsection{COBIT 2019}

COBIT~2019, de l'ISACA, est un référentiel de gouvernance et de gestion
des informations et des technologies~\cite{isaca2018cobitintro}. Il
organise quarante objectifs en cinq domaines~: un domaine de gouvernance
(évaluer, diriger, surveiller) et quatre domaines de
gestion~\cite{isaca2018cobitobjectifs}. L'auditeur public y puise des
critères pour la gouvernance de la fonction informatique, les projets, et
la gestion des changements (Chapitres~\ref{chap:exploitation}
et~\ref{chap:gouvernance}).

\subsection{ISO/IEC 27001 et 27002}

ISO/IEC~27001 fixe les exigences d'un système de management de la
sécurité de l'information~; ISO/IEC~27002 détaille les mesures de
sécurité, regroupées dans l'édition 2022 en quatre thèmes~: mesures
organisationnelles, liées aux personnes, physiques et
technologiques~\cite{iso27001_2022,iso27002_2022}. Ces normes sont sous
droits~: le cours les résume, il ne les reproduit pas.

\subsection{Les textes camerounais comme critères}

\begin{tableaucours}{L{3.4cm}L{8.8cm}}{%
    \tableauHeader{Texte} &
    \tableauHeader{Ce qu'il fournit comme critère}
}
Décret n°~2013/287 & Le droit du vérificateur à obtenir les pièces,
électroniques comprises \\
Loi n°~2010/012 & L'audit de sécurité des SI des organismes publics, les
obligations de sécurité des exploitants \\
Loi n°~2018/012 (régime financier) & Les principes de l'exécution de la
dépense et la séparation ordonnateur / comptable (à lire dans le
texte)~\cite{minfi2023regimefinancier} \\
Loi n°~2024/017 & La protection des données personnelles traitées par
les SI publics (paie, contribuables, usagers) \\
\end{tableaucours}
\vspace{-0.8em}
\begin{center}{\small\itshape Tableau~5.3. Textes camerounais utiles comme
critères}\end{center}

\subsection{Choisir et citer un critère}

\begin{boxprocedure}
\textbf{Fiche de critère}, une par test~:
\begin{itemize}
    \item l'énoncé du critère, en une phrase vérifiable~;
    \item sa source exacte (texte, article, alinéa~; ou référentiel et
          objectif), et sa force (Tableau~5.1)~;
    \item la date à laquelle le texte a été lu, et dans quelle version~;
    \item le test qui permettra de dire \OUI{} ou \NON{}.
\end{itemize}
\end{boxprocedure}

\begin{boxscenario}{Dossier MINÉQUIR. Le critère du contrôle 1}
\textbf{Constatation.} La DAF n'a aucune procédure écrite sur la
désactivation des comptes.

\medskip
\textbf{Analyse} (confiance élevée). Faute de règle interne, la
vérificatrice fonde son critère sur une bonne pratique
(principe~11 du COSO~; gestion des accès dans ISO/IEC~27002) et le
signale comme tel. Son constat portera donc sur deux points~: l'absence
de règle interne, qui est elle-même une déficience de conception, et les
9 comptes encore actifs, qui en sont la conséquence.
\end{boxscenario}

\separateur

% ============================================================================
\section{Déontologie}
\label{sec:c05-deontologie}
% ============================================================================

Les codes de l'INTOSAI, de l'IIA et de l'ISACA convergent sur quatre
exigences.

\begin{tableaucours}{L{3.0cm}L{9.2cm}}{%
    \tableauHeader{Exigence} &
    \tableauHeader{Ce qu'elle impose dans une mission SI}
}
Intégrité & Ne jamais modifier une donnée de l'entité~; travailler sur
une copie \\
Indépendance et objectivité & Ne pas auditer un système que l'on a
conçu, paramétré ou exploité \\
Compétence & Se faire assister par un spécialiste pour les tests
techniques que l'on ne maîtrise pas, et le dire au rapport \\
Confidentialité & Protéger les extractions (paie, contribuables) et les
détruire selon les règles à la fin de la mission \\
\end{tableaucours}
\vspace{-0.8em}
\begin{center}{\small\itshape Tableau~5.4. Quatre exigences
déontologiques dans une mission SI}\end{center}

\begin{boxmetier}{Vérificateur}
Écrivez dans le programme de travail, pour chaque test, la fiche de
critère. Lors du contradictoire, la première contestation porte presque
toujours sur le critère~: «~aucun texte ne nous l'impose~».
\end{boxmetier}

\begin{boxmetier}{Auditeur interne}
Si votre ministère a adopté une politique de sécurité fondée sur
ISO/IEC~27002, ses mesures deviennent des règles internes~: citez la
politique, pas la norme.
\end{boxmetier}

\separateur

\begin{boxresume}
\begin{itemize}
    \item \textbf{Norme et référentiel.} La norme dit comment auditer~; le
          référentiel fournit des critères.
    \item \textbf{L'INTOSAI pour le vérificateur.} ISSAI~100, 400 et 4000,
          GUID~5100, manuel de l'IDI.
    \item \textbf{L'IIA pour l'auditeur interne.} Mêmes exigences
          essentielles~: critère, preuve, documentation.
    \item \textbf{COBIT et ISO/IEC~27002 comme sources de critères}, qui ne
          s'imposent que si l'entité les a adoptés.
    \item \textbf{Une hiérarchie des critères.} Loi, règle interne,
          référentiel adopté, bonne pratique~; l'écart se qualifie selon la
          force du critère.
    \item \textbf{Une fiche de critère par test}, avec la source exacte et la
          date de lecture.
    \item \textbf{Quatre exigences déontologiques.} Intégrité,
          indépendance, compétence, confidentialité.
\end{itemize}
\end{boxresume}

\begin{boxexo}[title={\ding{45}\quad Exercice~5.1. Norme ou référentiel~? \niveaubase}]
Classez~: (a) GUID~5100~; (b) ISO/IEC~27002~; (c) ISSAI~4000~; (d) COBIT~2019~;
(e) normes de l'IIA~; (f) loi n°~2024/017.
\end{boxexo}

\begin{boxexo}[title={\ding{45}\quad Exercice~5.2. Qualifier l'écart \niveaumoyen}]
La CAPNO n'a pas testé ses sauvegardes depuis deux ans. Rédigez le
constat dans deux cas~: (a) la politique de sécurité de la CAPNO exige un
test semestriel~; (b) aucune règle interne n'existe.
\end{boxexo}

\begin{boxexo}[title={\ding{45}\quad Exercice~5.3. Fiches de critère \niveaumoyen}]
Établissez la fiche de critère des contrôles~2 (séparation LIQ / ORD) et~4
(blocage des doublons) du circuit de la dépense du MINÉQUIR
(Tableau~3.2).
\end{boxexo}

\begin{boxexo}[title={\ding{45}\quad Exercice~5.4. Un conflit d'intérêts \niveauexpert}]
Un vérificateur a participé, deux ans plus tôt, au comité de pilotage du
projet SIDÉPENSE. Analysez la situation au regard des exigences
déontologiques, et proposez ce que son chef de division doit décider.
\end{boxexo}

\begin{boxinfo}
\textbf{Transition.} Le mandat et les normes sont posés. La Partie~III
conduit la mission elle-même, en commençant par sa préparation
(Chapitre~\ref{chap:preparer}).
\end{boxinfo}
}{}

    % ======================================================================
    % PARTIE 3 — DÉMARCHE DE MISSION (J04, J05)
    % « Pas de constat sans critère, pas de conclusion sans preuve. »
    % Standards de référence : ISO 19011, ISSAI 4000, GUID 5100, manuel IDI ;
    %   articles A0 (hal-05701007) et A1 (hal-05701270).
    % Niveau : intermédiaire
    % ======================================================================
    \part{Démarche de Mission}
    \label{part:demarche}

        \IfFileExists{partie3/P03_C06_preparer_mission.tex}{% ============================================================================
% Fichier  : partie3/P03_C06_preparer_mission.tex
% Titre    : La Mission : Principes, Préparation et Décomposition
% Auteur   : MINKA MI NGUIDJOÏ Thierry Emmanuel
% Version  : 0.2 -- novembre 2026 (réécriture étoffée, alignée sur INF3226
%            C07 à C09, A0 et A1)
% ============================================================================
% Sources : ISO 19011:2018 (principes) ; ISSAI 4000 via le manuel IDI ;
% A0 « Vers l'auditologie » (hal-05701007) : axiomes 1, 2, 5, 6, proposition
% de décomposition hiérarchique, gabarit canonique ; A1 (hal-05701270).
% Miller (1956) ; Simon (1962).
% Faits MINÉQUIR (LAB-01) : ordre de mission du 16/03/2026 ; équipe
% NGO BIYONG (chef de mission) + ESSAMA ; 4 domaines D1 Accès, D2 Fichier
% des fournisseurs, D3 Circuit du mandat, D4 Exploitation ; cotation des
% risques R1-R8 inchangée (R1 et R3 à 9) ; contrôle D1.1.2 = désactivation
% à la date d'effet (2 critères).
% ============================================================================

\chapter[La mission : principes, préparation, décomposition]{La Mission : Principes, Préparation et Décomposition}
\label{chap:preparer}
\chaptermark{Préparer la mission}

\epigraph{%
    « Une mission se gagne avant le terrain. »%
}{Adage des vérificateurs}

\niveauChapitre{Intermédiaire}{Chapitres~\ref{chap:eval-sci}
et~\ref{chap:normes}.}

\begin{boxobjectifs}
\begin{itemize}
    \item Énoncer les principes de l'audit et les relier aux axiomes de
          l'auditologie.
    \item Décrire le cycle d'une mission en cinq phases, avec les entrées et
          les sorties de chacune.
    \item Rédiger un ordre de mission et des termes de référence qui fixent
          l'objectif, le périmètre, la période et les critères.
    \item Organiser la demande de pièces en trois strates, et décider si une
          pièce absente vaut \NON{} ou \NEK{}.
    \item Conduire une prise de connaissance bornée, par couches du système
          d'information.
    \item Hiérarchiser les zones d'audit en croisant le risque et la
          connaissance que l'on en a.
    \item Décomposer le domaine audité selon le gabarit canonique à huit
          niveaux, et rédiger une fiche d'assertion par contrôle.
    \item Établir le programme de travail.
\end{itemize}
\end{boxobjectifs}

\bigskip

% ============================================================================
\section*{Pourquoi ce chapitre}
% ============================================================================

Une mission sur un système d'information dispose de peu de temps et d'un
objet immense~: des milliers d'opérations, des dizaines de profils, des
journaux de plusieurs gigaoctets. Tout tester est impossible, et ce que
l'on ne regarde pas reste inconnu. La préparation décide donc de ce que la
mission saura, et de ce qu'elle ignorera. Elle se juge à une question~:
\emph{le programme de travail porte-t-il sur les zones où le risque est
élevé et la connaissance faible~?}

Ce chapitre suit la méthode du cours d'audit des systèmes d'information de
l'auteur, qui est l'instanciation, pour le praticien, du programme de
recherche de l'auditologie~\cite{minka2026auditologie}. Il en reprend les
trois outils de préparation~: le cycle de la mission, l'évaluation croisée
du risque et de la connaissance, et le gabarit canonique de décomposition.

\medskip
\textbf{Plan du chapitre.}
\begin{enumerate}
    \item Les principes de l'audit (section~\ref{sec:c06-principes}).
    \item Le cycle de la mission (section~\ref{sec:c06-cycle}).
    \item L'ordre de mission et les termes de référence
          (section~\ref{sec:c06-ordre}).
    \item Les pièces à demander (section~\ref{sec:c06-pieces}).
    \item La prise de connaissance (section~\ref{sec:c06-connaissance}).
    \item Hiérarchiser~: risque et connaissance
          (section~\ref{sec:c06-risques}).
    \item Décomposer~: le gabarit canonique et la fiche d'assertion
          (section~\ref{sec:c06-decomposer}).
    \item Le programme de travail (section~\ref{sec:c06-programme}).
\end{enumerate}

\separateur

% ============================================================================
\section{Les principes de l'audit}
\label{sec:c06-principes}
% ============================================================================

La norme ISO~19011 énonce sept principes de l'audit~\cite{iso19011_2018}~;
les normes de l'INTOSAI et de l'IIA en portent l'équivalent. Ils ne sont
pas des vœux~: chacun commande une pratique de la mission.

\begin{tableaucours}{L{3.4cm}L{4.6cm}L{4.2cm}}{%
    \tableauHeader{Principe} &
    \tableauHeader{Ce qu'il impose} &
    \tableauHeader{Dans une mission SI}
}
Intégrité & Agir honnêtement, avec diligence & Ne jamais modifier une
donnée de l'entité \\
Restitution impartiale & Rapporter fidèlement, y compris les difficultés
& Écrire les \NEK{} restants au rapport \\
Conscience professionnelle & Exercer le soin et le jugement requis &
Vérifier l'intégrité de chaque extraction \\
Confidentialité & Protéger l'information obtenue & Protéger les fichiers
de paie et de fournisseurs \\
Indépendance & Fonder l'objectivité des conclusions & Ne pas auditer un
système que l'on a conçu \\
Approche fondée sur la preuve & Conclure sur des preuves vérifiables &
Partir du système, pas des classeurs \\
Approche par les risques & Porter l'effort là où il compte & Croiser
risque et connaissance (section~\ref{sec:c06-risques}) \\
\end{tableaucours}
\vspace{-0.8em}
\begin{center}{\small\itshape Tableau~6.1. Les sept principes de l'audit
(d'après ISO~19011:2018)}\end{center}

\begin{boxinfo}
\textbf{Des principes aux axiomes.} L'article \emph{Vers l'auditologie}
montre que plusieurs de ces principes se déduisent d'axiomes plus
fondamentaux~\cite{minka2026auditologie}. L'approche fondée sur la preuve
découle de l'\emph{observabilité partielle} (axiome~1)~: l'auditeur ne voit
qu'une partie du système, il doit donc dire sur quoi il s'appuie.
L'indépendance n'est pas un axiome, mais une conséquence de
l'\emph{interaction stratégique} (axiome~4) et de l'\emph{opposabilité}
(axiome~6)~: une conclusion qu'un tiers peut contester doit venir de
quelqu'un qui n'a pas intérêt au résultat.
\end{boxinfo}

\subsection{Les principes en conflit}

Les principes se heurtent parfois. La confidentialité demande de limiter la
diffusion d'une extraction~; la restitution impartiale demande d'en
montrer les résultats. La règle du cours est de ne jamais sacrifier
l'opposabilité~: on peut anonymiser une liste d'agents dans le rapport, à
condition de conserver au dossier la liste nominative et son empreinte.

\separateur

% ============================================================================
\section{Le cycle de la mission}
\label{sec:c06-cycle}
% ============================================================================

\subsection{Cinq phases}

\begin{center}
\begin{tikzpicture}[
    ph/.style = {draw=CouleurBordDef, fill=CouleurDef, rounded corners=3pt,
        text width=2.1cm, minimum height=0.9cm, align=center,
        font=\scriptsize\bfseries\color{CouleurPrimaire}},
    arrow/.style = {-{Stealth[length=4pt]}, thick, color=CouleurSecondaire}
]
\foreach \t [count=\i] in {Préparer, Planifier, Exécuter, Conclure, Suivre}
    \node[ph] (p\i) at (2.55*\i, 0) {\t};
\foreach \i [evaluate=\i as \j using int(\i+1)] in {1,2,3,4}
    \draw[arrow] (p\i) -- (p\j);
\draw[arrow, dashed] (p5.south) -- ++(0,-0.55) -| (p1.south)
    node[pos=0.25, below, font=\scriptsize\itshape, color=CouleurSecondaire]
    {le cycle suivant part de l'état laissé par le précédent};
\end{tikzpicture}
\end{center}
\vspace{-0.2em}
\begin{center}{\small\itshape Figure~6.1. Le cycle de la mission}\end{center}

La flèche de retour n'est pas un ornement. L'auditologie pose
l'\emph{individuation temporelle} (axiome~5)~: le système audité évolue,
et son état à la mission suivante dépend de la mission
précédente~\cite{minka2026auditologie}. Le suivi n'est donc pas une
formalité~: il est la première phase du cycle suivant
(Chapitre~\ref{chap:constater}).

\subsection{Entrées et sorties de chaque phase}

\begin{tableaucours}{L{2.2cm}L{4.6cm}L{5.4cm}}{%
    \tableauHeader{Phase} &
    \tableauHeader{Entrées} &
    \tableauHeader{Sorties}
}
Préparer & Mandat, rapports antérieurs, textes & Ordre de mission, termes
de référence, demande de pièces \\
Planifier & Prise de connaissance, cotation des risques & Décomposition,
fiches d'assertion, programme de travail \\
Exécuter & Programme, pièces, accès au système & Feuilles de test, cotes
\OUI{}, \NON{}, \NEK{}, preuves et empreintes \\
Conclure & Cotes et preuves & Mesures par domaine, observations, rapport
provisoire, contradictoire, rapport définitif \\
Suivre & Plan d'action & Tableau de suivi, comparaison des cycles \\
\end{tableaucours}
\vspace{-0.8em}
\begin{center}{\small\itshape Tableau~6.2. Entrées et sorties des cinq
phases}\end{center}

\subsection{L'équipe et la supervision}

Une mission du CONSUPE compte au moins un chef de mission et un
vérificateur. Le chef de mission répond du programme, de la cohérence des
cotes et du rapport~; il vise chaque feuille de test. La supervision n'est
pas une relecture finale~: elle a lieu pendant l'exécution, quand une cote
peut encore être reprise. Au MINÉQUIR, Mme~NGO~BIYONG revoit chaque soir
les feuilles de M.~ESSAMA.

\separateur

% ============================================================================
\section{L'ordre de mission et les termes de référence}
\label{sec:c06-ordre}
% ============================================================================

\subsection{L'ordre de mission}

L'ordre de mission est la première pièce du dossier. Il fonde le droit
d'accès du vérificateur (décret n°~2013/287, art.~2, al.~4) et fixe ce que
la mission doit démontrer.

\begin{tableaucours}{L{3.0cm}L{9.2cm}}{%
    \tableauHeader{Rubrique} &
    \tableauHeader{Ordre de mission du 16 mars 2026 (MINÉQUIR)}
}
Fondement & Décret n°~2013/287, art.~2, al.~3 (contrôle des SI) et
al.~4 \\
Objectif & Évaluer le contrôle interne du circuit de la dépense de la DAF
dans SIDÉPENSE, et la régularité des paiements de l'exercice \\
Périmètre & DAF du MINÉQUIR~; application SIDÉPENSE~; poste comptable
pour la seule étape du paiement \\
Période & Exercice 2025 \\
Durée & Quatre semaines, du 16~mars au 10~avril~2026 \\
Équipe & Mme~NGO~BIYONG (chef de mission), M.~ESSAMA (vérificateur) \\
Livrables & Rapport provisoire, contradictoire, rapport définitif \\
\end{tableaucours}
\vspace{-0.8em}
\begin{center}{\small\itshape Tableau~6.3. Les rubriques d'un ordre de
mission}\end{center}

\subsection{Les termes de référence : quatre questions}

Les termes de référence précisent l'ordre de mission. Ils répondent à
quatre questions, et à elles seules.

\begin{enumerate}
    \item \textbf{Que doit-on démontrer~?} L'objectif, écrit comme une
          question fermée~: «~Le contrôle interne du circuit de la dépense
          dans SIDÉPENSE est-il efficace pour l'exercice 2025~?~»
    \item \textbf{Sur quoi~?} Le périmètre~: entités, applications,
          processus inclus, et ce qui est exclu (ici, la paie et la recette).
    \item \textbf{Par rapport à quoi~?} Les critères et leur source
          (Chapitre~\ref{chap:normes}).
    \item \textbf{Avec quelle assurance~?} Raisonnable ou limitée
          (Chapitre~\ref{chap:audit-public}), et les seuils de décision que
          la mission retiendra (Chapitre~\ref{chap:constater}).
\end{enumerate}

\begin{boxalert}
\textbf{Violation typique.} Un ordre de mission qui dit «~contrôler le
système informatique de la DAF~», sans objectif, sans période et sans
périmètre.

\smallskip
\textbf{Règle~:} un objectif s'écrit comme une question à laquelle le
rapport répondra par oui ou par non~; le périmètre dit ce qui est exclu
autant que ce qui est inclus.
\end{boxalert}

\separateur

% ============================================================================
\section{Les pièces à demander}
\label{sec:c06-pieces}
% ============================================================================

\subsection{Trois strates}

La demande de pièces part dès le premier jour, par écrit, avec un délai.
Elle s'organise en trois strates, de la plus générale à la plus probante.

\begin{tableaucours}{L{2.8cm}L{4.8cm}L{4.6cm}}{%
    \tableauHeader{Strate} &
    \tableauHeader{Contenu} &
    \tableauHeader{Au MINÉQUIR}
}
1. Gouvernance & Organigramme, décisions, schéma directeur, politiques &
Organigramme de la DAF, décision de création de la cellule \\
2. Fonctionnement & Procédures, manuels, contrats, comptes rendus &
Procédures du circuit, contrat de maintenance de SIDÉPENSE \\
3. Traces & Extractions, journaux, paramètres, faites devant l'auditeur &
Comptes et profils, mandats 2025, journal des opérations \\
\end{tableaucours}
\vspace{-0.8em}
\begin{center}{\small\itshape Tableau~6.4. Les trois strates de
pièces}\end{center}

\begin{tableaucours}{cL{6.2cm}L{3.0cm}L{2.2cm}}{%
    \tableauHeader{N°} &
    \tableauHeader{Pièce} &
    \tableauHeader{Destinataire} &
    \tableauHeader{Délai}
}
1 & Liste des utilisateurs de SIDÉPENSE et de leurs profils, extraite
devant la vérificatrice & Cellule informatique & J+2 \\
2 & Liste des mutations et départs de 2025 & Ressources humaines & J+3 \\
3 & Extraction des mandats 2025 et du journal des opérations & Cellule
informatique & J+5 \\
4 & Journal des modifications du fichier des fournisseurs & Cellule
informatique & J+5 \\
5 & Procédures écrites du circuit et de la gestion des comptes & DAF &
J+3 \\
6 & Dernier rapport d'audit de sécurité de l'ANTIC (date, périmètre, plan
d'action) & Cellule informatique & J+5 \\
\end{tableaucours}
\vspace{-0.8em}
\begin{center}{\small\itshape Tableau~6.5. Demande de pièces de la mission
MINÉQUIR}\end{center}

\subsection{Une pièce absente : NON ou NEK~?}

Une pièce qui ne vient pas n'a pas toujours le même sens. Le cours
applique une règle simple, cohérente avec l'instrument de l'article
A1~\cite{minka2026nek}.

\begin{itemize}
    \item \textbf{\NON{}}~: la pièce devait exister (une règle l'impose), et
          l'entité confirme par écrit qu'elle n'existe pas. L'absence est
          une constatation.
    \item \textbf{\NEK{}}~: la pièce peut exister ailleurs, ou n'a pas
          encore été produite malgré la relance. On ne sait pas~; on le dit,
          avec les diligences faites.
\end{itemize}
Un \NEK{} qui reste à la fin de la mission devient une \emph{limitation
d'étendue}, écrite au rapport (Chapitre~\ref{chap:constater}).

\separateur

% ============================================================================
\section{La prise de connaissance}
\label{sec:c06-connaissance}
% ============================================================================

\subsection{Une prise de connaissance bornée}

La prise de connaissance n'a pas pour but de tout comprendre, mais de
savoir assez pour formuler des hypothèses de risque. L'article A0 rappelle
que l'auditeur est un agent cognitivement limité
(axiome~2)~\cite{minka2026auditologie}~: une prise de connaissance sans
borne produit un dossier que personne ne relit. Le cours la limite à deux
jours sur quatre semaines, et la conduit par couches.

\subsection{Lire le système d'information en couches}

\begin{tableaucours}{L{2.8cm}L{4.8cm}L{4.6cm}}{%
    \tableauHeader{Couche} &
    \tableauHeader{Question} &
    \tableauHeader{Au MINÉQUIR}
}
Organisation & Qui décide, qui exécute, qui contrôle~? & DAF propriétaire,
cellule exécutante, contrôle financier \\
Processus & Quelles étapes, dans quel ordre~? & Engagement, liquidation,
ordonnancement, paiement \\
Applications & Quelles applications portent les étapes~? & SIDÉPENSE,
logiciel du poste comptable \\
Données & Quelles données de référence conditionnent tout~? & Fichier des
fournisseurs, table des profils \\
Infrastructure & Où tourne le système, qui l'administre~? & Serveur de la
cellule, sauvegardes locales \\
\end{tableaucours}
\vspace{-0.8em}
\begin{center}{\small\itshape Tableau~6.6. Les couches du système
d'information}\end{center}

\subsection{La cartographie des applications et des flux}

\begin{tableaucours}{L{2.8cm}L{3.2cm}L{3.0cm}L{3.2cm}}{%
    \tableauHeader{Application} &
    \tableauHeader{Fonction} &
    \tableauHeader{Propriétaire} &
    \tableauHeader{Flux sortant}
}
SIDÉPENSE & Engagement, liquidation, ordonnancement, fournisseurs & DAF
& Mandats vers le poste comptable (fichier quotidien) \\
Logiciel du poste comptable & Paiement, comptabilité & Poste comptable &
Ordres de virement vers les banques \\
Fichier du personnel & Mutations, départs & Direction des ressources
humaines & Aucun vers SIDÉPENSE (constat) \\
\end{tableaucours}
\vspace{-0.8em}
\begin{center}{\small\itshape Tableau~6.7. Cartographie des applications
du périmètre}\end{center}

La dernière ligne est déjà un résultat~: aucun flux ne relie le fichier du
personnel à SIDÉPENSE. Les mutations ne sont connues de la cellule
informatique que si quelqu'un pense à les lui signaler. La cartographie a
désigné un risque avant tout test.

\subsection{Conduire un entretien de prise de connaissance}

\begin{boxprocedure}
\textbf{Entretien de prise de connaissance}~:
\begin{itemize}
    \item préparer une liste de questions ouvertes, par couche~;
    \item être deux~: l'un interroge, l'autre note~;
    \item demander, pour chaque affirmation, la pièce qui la montre~;
    \item relire le compte rendu à l'interlocuteur et le faire valider~;
    \item noter à part ce qui est déclaré sans pièce~: ce sont les futurs
          \NEK{} à lever.
\end{itemize}
\end{boxprocedure}

\separateur

% ============================================================================
\section{Hiérarchiser : risque et connaissance}
\label{sec:c06-risques}
% ============================================================================

\subsection{La cotation du risque}

Pour chaque risque de la matrice (Chapitre~\ref{chap:eval-sci}), on cote
de 1 (faible) à 3 (élevé) le \termecle{risque inhérent}\index{risque
inhérent} (gravité et probabilité de l'anomalie si aucun contrôle
n'existait) et le \termecle{risque lié au contrôle}\index{risque lié au
contrôle} (probabilité que le contrôle interne, tel qu'on le connaît, ne
l'empêche ni ne la détecte). Leur produit, de 1 à 9, est la priorité.

\begin{tableaucours}{L{5.4cm}ccc}{%
    \tableauHeader{Risque} &
    \tableauHeader{Inhérent} &
    \tableauHeader{Contrôle} &
    \tableauHeader{Priorité}
}
R1. Un agent parti agit encore & 3 & 3 & 9 \\
R2. Un agent émet seul un mandat & 3 & 2 & 6 \\
R3. Le bénéficiaire est changé & 3 & 3 & 9 \\
R4. Une facture est payée deux fois & 3 & 2 & 6 \\
R5. Le mandat n'est pas justifié & 2 & 1 & 2 \\
R6. Une opération ne laisse pas de trace & 2 & 1 & 2 \\
R7. Une anomalie de fin de gestion passe & 2 & 3 & 6 \\
R8. Les traces disparaissent & 2 & 2 & 4 \\
\end{tableaucours}
\vspace{-0.8em}
\begin{center}{\small\itshape Tableau~6.8. Cotation préliminaire des
risques du circuit de la dépense (convention du cours)}\end{center}

\subsection{Le risque et la connaissance}

La cotation du risque ne suffit pas. L'auditologie ajoute une seconde
dimension~: ce que l'on \emph{sait} de chaque zone. Une zone à risque élevé
que l'on connaît bien (un audit récent l'a testée) appelle moins d'effort
qu'une zone à risque moyen dont on ignore tout. C'est la distinction d'A0
entre non-conformité \emph{connue} et non-conformité
\emph{inconnue}~\cite{minka2026auditologie}~: la seconde est la plus
dangereuse, parce qu'elle ne se voit pas.

\begin{tableaucours}{L{3.0cm}L{4.4cm}L{4.8cm}}{%
    \tableauHeader{} &
    \tableauHeader{Connaissance élevée} &
    \tableauHeader{Connaissance faible}
}
\textbf{Risque élevé} & Tester ce qui a changé depuis le dernier audit &
\textbf{Priorité absolue}~: tests étendus, preuves de strate~3 \\
\textbf{Risque faible} & Cheminement ou revue analytique & Prise de
connaissance complémentaire, puis décider \\
\end{tableaucours}
\vspace{-0.8em}
\begin{center}{\small\itshape Tableau~6.9. Croiser le risque et la
connaissance}\end{center}

\begin{boxscenario}{Dossier MINÉQUIR. Où regarder d'abord}
\textbf{Constatation.} R1 (agents partis) et R3 (RIB) sont cotés 9. Aucun
audit n'a jamais porté sur SIDÉPENSE~; l'inspection de janvier n'a examiné
que des dossiers papier.

\medskip
\textbf{Analyse} (confiance moyenne). Les deux risques sont à la fois
élevés et inconnus~: ce sont les zones de priorité absolue. La
vérificatrice leur consacre la première semaine, puis les trois risques
cotés~6. Elle ne teste R5 et R6 que par un cheminement~: leur priorité ne
justifie pas un échantillon.
\end{boxscenario}

\separateur

% ============================================================================
\section{Décomposer : le gabarit canonique et la fiche d'assertion}
\label{sec:c06-decomposer}
% ============================================================================

\subsection{Pourquoi décomposer}

Un auditeur ne peut pas examiner «~les accès~» d'un bloc. George Miller a
montré que la mémoire de travail ne retient qu'environ sept éléments à la
fois, plus ou moins deux~\cite{miller1956}~; Herbert Simon, que les
systèmes complexes sont hiérarchiques et presque
décomposables~\cite{simon1962complexity}. L'article A0 en tire une
proposition~: tout instrument qui couvre un espace d'audit non trivial doit
être organisé en une décomposition hiérarchique dont chaque niveau respecte
cette borne~\cite{minka2026auditologie}.

\begin{boxdef}
\termecle{Espace d'audit}\index{espace d'audit}~: ensemble des objets sur
lesquels porte une mission (processus, contrôles, données, traces), tel que
le délimite l'ordre de mission.

\medskip
\termecle{Décomposition}\index{décomposition}~: découpage hiérarchique de
l'espace d'audit en niveaux emboîtés, chaque élément d'un niveau se divisant
en éléments du niveau inférieur.
\end{boxdef}

La décomposition apporte quatre choses~: la \textbf{couverture} (on sait ce
qui a été audité et ce qui ne l'a pas été), la \textbf{traçabilité} (chaque
conclusion remonte à une trace, ce qui la rend opposable, axiome~6), la
\textbf{localisation de l'incertitude} (les \NEK{} se situent à un endroit
précis) et la \textbf{répartition du travail}.

\subsection{Le gabarit canonique à huit niveaux}

Le programme de l'auditologie propose un \termecle{gabarit
canonique}\index{gabarit canonique} de décomposition à huit
niveaux~\cite{minka2026auditologie}. Les quatre premiers disent \emph{ce
qui est audité}~; les quatre derniers, \emph{comment on le vérifie}. Le
contrôle fait la charnière~: c'est l'objet sur lequel l'auditeur se
prononce.

\begin{center}
\begin{tikzpicture}[
    niv/.style = {draw=CouleurBordDef, fill=CouleurDef, rounded corners=3pt,
        minimum width=2.9cm, minimum height=0.58cm, align=center,
        font=\footnotesize\bfseries\color{CouleurPrimaire}},
    note/.style = {font=\scriptsize\itshape\color{CouleurSecondaire}, anchor=west, text width=7.2cm},
    fleche/.style = {-{Stealth[length=4pt]}, thick, color=CouleurSecondaire}
]
\foreach \t/\q [count=\i] in {
    1. Domaine/D1 Accès à SIDÉPENSE,
    2. Sous-domaine/D1.1 Cycle de vie des comptes,
    3. Processus/Mutation ou départ d'un agent,
    4. Activité/Retrait des accès,
    5. Contrôle/D1.1.2 Désactivation à la date d'effet,
    6. Test/Rapprocher mutations et comptes actifs,
    7. Preuve/Extraction des comptes et son empreinte,
    8. Trace/Champ « date de désactivation » et journal} {
    \node[niv] (n\i) at (0, -0.78*\i) {\t};
    \node[note] at (1.6, -0.78*\i) {\q};
}
\foreach \i [evaluate=\i as \j using int(\i+1)] in {1,...,7}
    \draw[fleche] (n\i) -- (n\j);
\end{tikzpicture}
\end{center}
\vspace{-0.2em}
\begin{center}{\small\itshape Figure~6.2. Le gabarit canonique appliqué au
contrôle des agents partis (d'après~\cite{minka2026auditologie})}\end{center}

\begin{boxalert}
\textbf{Violation typique.} Mettre dans la même liste «~une politique
d'accès existe~» (un document) et «~le journal est activé~» (une trace),
au même niveau que «~les accès~» (un domaine).

\smallskip
\textbf{Règle~:} chaque élément se place à un seul niveau, celui de sa
nature. Un mélange de niveaux rend la couverture invérifiable.
\end{boxalert}

\subsection{Preuve et trace : deux choses différentes}

La \emph{trace} existe dans le système, que l'auditeur passe ou non~: une
entrée de journal, un champ de date. La \emph{preuve} est ce que l'auditeur
en tire et conserve~: l'extraction horodatée, son empreinte, la requête
qui l'a produite. Une trace non extraite n'est pas une preuve~; une preuve
qui ne renvoie à aucune trace (une capture d'écran sans contexte) est
fragile.

\subsection{La décomposition de la mission MINÉQUIR}

\begin{center}
\begin{tikzpicture}[
    niv/.style = {draw=CouleurBordDef, fill=CouleurDef, rounded corners=3pt,
        text width=2.9cm, minimum height=0.8cm, align=center,
        font=\scriptsize\bfseries\color{CouleurPrimaire}},
    arrow/.style = {-{Stealth[length=4pt]}, thick, color=CouleurSecondaire}
]
\node[niv, text width=6cm] (obj) at (5.1,0) {Objectif~: le contrôle interne
du circuit de la dépense est-il efficace~?};
\node[niv] (d1) at (0,-1.6) {D1 Accès};
\node[niv] (d2) at (3.4,-1.6) {D2 Fichier des fournisseurs};
\node[niv] (d3) at (6.8,-1.6) {D3 Circuit du mandat};
\node[niv] (d4) at (10.2,-1.6) {D4 Exploitation};
\foreach \d in {d1,d2,d3,d4} \draw[arrow] (obj.south) -- (\d.north);
\node[niv, font=\tiny] at (0,-3.0) {Contrôles 1, 2};
\node[niv, font=\tiny] at (3.4,-3.0) {Contrôle 3};
\node[niv, font=\tiny] at (6.8,-3.0) {Contrôles 4, 5, 7};
\node[niv, font=\tiny] at (10.2,-3.0) {Contrôles 6, 8};
\end{tikzpicture}
\end{center}
\vspace{-0.4em}
\begin{center}{\small\itshape Figure~6.3. Les quatre domaines de la
mission MINÉQUIR et les huit contrôles du prologue}\end{center}

\subsection{La fiche d'assertion}

Le gabarit dit \emph{où} se trouve un contrôle. La \termecle{fiche
d'assertion}\index{fiche d'assertion} dit \emph{ce que l'auditeur affirme}
à son sujet, et ce qui arrive si l'affirmation est fausse.

\begin{tableaucours}{L{2.4cm}L{5.0cm}L{4.8cm}}{%
    \tableauHeader{Attribut} &
    \tableauHeader{Rôle} &
    \tableauHeader{Forme}
}
Objectif & Ce que le contrôle doit garantir & À l'infinitif~: « S'assurer
que… » \\
Critère & La condition vérifiable qui définit la conformité & Au présent~:
une condition unique, observable \\
Question & La question posée pour coter le contrôle & Interrogative~;
réponse \OUI{}, \NON{} ou \NEK{} \\
Risque & L'état de non-conformité si la réponse est \NON{} & L'antithèse
du critère \\
Conséquence & L'effet sur l'entité si le risque se réalise & Financier,
opérationnel, juridique \\
\end{tableaucours}
\vspace{-0.8em}
\begin{center}{\small\itshape Tableau~6.10. Les attributs de la fiche
d'assertion}\end{center}

\begin{boxobjet}{D1.1.2 · Désactivation à la date d'effet}
\textbf{Position~:} D1 Accès $\rightarrow$ D1.1 Cycle de vie des comptes
$\rightarrow$ mutation ou départ $\rightarrow$ retrait des accès
$\rightarrow$ contrôle D1.1.2.

\textbf{Objectif~:} s'assurer qu'aucun agent muté ou parti ne peut agir
dans SIDÉPENSE.

\textbf{Critère~1~:} le compte est désactivé au plus tard à la date d'effet
de la mutation. \textbf{Question~1~:} chaque compte d'agent muté en 2025
a-t-il été désactivé au plus tard à la date d'effet~?

\textbf{Critère~2~:} la DRH notifie par écrit chaque mutation à la cellule
informatique. \textbf{Question~2~:} chaque mutation de 2025 a-t-elle fait
l'objet d'une notification écrite~?

\textbf{Test~:} rapprocher la liste des 31 mutations et départs de 2025 de
la liste des comptes actifs au 23~mars~2026~; demander les notifications.

\textbf{Preuve et trace~:} extractions et empreintes~; champ de date de
désactivation~; notifications.

\textbf{Risque~:} un compte resté actif est utilisé par son ancien
titulaire ou par un tiers.

\textbf{Conséquence~:} paiement indu, impossibilité d'imputer une
opération à une personne, suites devant le CDBF.
\end{boxobjet}

\subsection{De la fiche à la cotation}

Chaque critère reçoit sa propre question et sa propre cote. La cote du
contrôle s'en déduit par une règle de composition~: un seul critère
\NON{} suffit pour que le contrôle soit \NON{}~; à défaut, un seul critère
\NEK{} donne \NEK{}~; le contrôle n'est \OUI{} que si tous ses critères le
sont.

\begin{tableaucours}{L{3.4cm}L{3.4cm}cL{3.8cm}}{%
    \tableauHeader{Critère 1 (désactivé)} &
    \tableauHeader{Critère 2 (notifié)} &
    \tableauHeader{Cote D1.1.2} &
    \tableauHeader{Lecture}
}
\OUI{} & \OUI{} & \OUI{} & Contrôle conforme \\
\OUI{} & \NON{} & \NON{} & Désactivé par chance, sans procédure \\
\NEK{} & \OUI{} & \NEK{} & Notifié, mais désactivation non établie \\
\NON{} & \NEK{} & \NON{} & Un défaut établi suffit \\
\end{tableaucours}
\vspace{-0.8em}
\begin{center}{\small\itshape Tableau~6.11. Composer les cotes de deux
critères}\end{center}

\subsection{Les règles de qualité}

\begin{itemize}
    \item \textbf{Atomicité}~: l'objectif, le critère et la question ne
          contiennent ni «~et~» ni «~ou~»~; toute assertion composée est
          scindée.
    \item \textbf{Exhaustivité}~: les éléments d'un niveau couvrent
          entièrement l'élément du niveau supérieur.
    \item \textbf{Non-redondance}~: un contrôle n'appartient qu'à un
          domaine, sinon il est compté deux fois.
    \item \textbf{Borne cognitive}~: au plus sept à neuf éléments sous un
          même parent~; au-delà, on ajoute un niveau.
\end{itemize}

\separateur

% ============================================================================
\section{Le programme de travail}
\label{sec:c06-programme}
% ============================================================================

\begin{boxprocedure}
\textbf{Le programme de travail}, une ligne par contrôle~:
\begin{itemize}
    \item le code et l'intitulé du contrôle (position dans le gabarit)~;
    \item le risque couvert, sa priorité et la connaissance que l'on en a~;
    \item la fiche d'assertion (critères, questions)~;
    \item le test prévu (cheminement, échantillon, essai, analyse de toute
          la population) et sa taille (Chapitre~\ref{chap:preuve})~;
    \item les pièces à demander, par strate~;
    \item les seuils de décision retenus (Chapitre~\ref{chap:constater})~;
    \item le membre de l'équipe et la date prévue.
\end{itemize}
\end{boxprocedure}

\begin{boxmetier}{Vérificateur}
Gardez la trace de chaque relance. Une pièce non remise à l'échéance,
malgré relance écrite, devient soit un \NON{} (si elle devait exister),
soit une limitation d'étendue à mentionner au rapport.
\end{boxmetier}

\begin{boxmetier}{Inspecteur}
Même pour une inspection courte, écrivez l'objectif et le périmètre en
tête de votre note de mission~: ils déterminent la portée de votre
conclusion (Chapitre~\ref{chap:prologue}).
\end{boxmetier}

\begin{boxmetier}{Auditeur interne}
Votre plan d'audit annuel est une cotation des risques et des
connaissances à l'échelle de l'entité~: les applications financières, et
leurs jointures, doivent y figurer, avec la date du dernier audit qui les
a testées.
\end{boxmetier}

\separateur

\begin{boxresume}
\begin{itemize}
    \item \textbf{Sept principes, des axiomes derrière.} La preuve vient de
          l'observabilité partielle~; l'indépendance, de l'interaction
          stratégique et de l'opposabilité.
    \item \textbf{Un cycle en cinq phases}~: préparer, planifier, exécuter,
          conclure, suivre~; le suivi ouvre le cycle suivant.
    \item \textbf{Quatre questions de référence}~: quoi démontrer, sur quoi,
          par rapport à quoi, avec quelle assurance.
    \item \textbf{Trois strates de pièces}~; une pièce absente est \NON{} si
          elle devait exister, \NEK{} sinon.
    \item \textbf{Risque et connaissance.} La priorité absolue va aux zones
          à risque élevé que l'on connaît mal.
    \item \textbf{Le gabarit canonique}~: huit niveaux, du domaine à la
          trace~; une fiche d'assertion par contrôle, une cote par critère.
    \item \textbf{Un programme ligne par ligne}, écrit avant le terrain.
\end{itemize}
\end{boxresume}

\begin{boxexo}[title={\ding{45}\quad Exercice~6.1. Un ordre de mission \niveaubase}]
Rédigez l'ordre de mission et les termes de référence d'une vérification
de la paie de la CAPNO (Garoua) pour l'exercice 2025, en répondant aux
quatre questions de la section~\ref{sec:c06-ordre}.
\end{boxexo}

\begin{boxexo}[title={\ding{45}\quad Exercice~6.2. Risque et connaissance \niveaumoyen}]
Pour la paie de la CAPNO, cotez au moins six risques, indiquez pour chacun
la connaissance dont vous disposez (élevée ou faible) et placez-les dans le
Tableau~6.9. Dites lesquels vous testerez par analyse de toute la
population, et lesquels par cheminement.
\end{boxexo}

\begin{boxexo}[title={\ding{45}\quad Exercice~6.3. Une fiche d'assertion \niveaumoyen}]
Rédigez la fiche d'assertion du contrôle «~seconde validation de toute
modification de RIB~» (domaine D2), avec sa position dans le gabarit, deux
critères atomiques, leurs questions, le test, la preuve, la trace, le
risque et la conséquence. Donnez la cote du contrôle si le premier critère
est \OUI{} et le second \NEK{}.
\end{boxexo}

\begin{boxexo}[title={\ding{45}\quad Exercice~6.4. Décomposer e-RECETTES \niveauexpert}]
Décomposez l'objectif «~le système e-RECETTES garantit-il l'exhaustivité
des recettes enregistrées~?~» selon le gabarit canonique jusqu'au niveau
du contrôle, pour au moins trois sous-domaines. Vérifiez les quatre règles
de qualité et signalez une violation que vous avez dû corriger.
\end{boxexo}

\begin{boxinfo}
\textbf{Transition.} Le programme dit quoi tester, et où se trouve chaque
contrôle. Le Chapitre~\ref{chap:preuve} dit comment répondre à chaque
question~: avec quelle preuve, et que faire quand on ne sait pas.
\end{boxinfo}
}{}
        \IfFileExists{partie3/P03_C07_collecter_preuve.tex}{% ============================================================================
% Fichier  : partie3/P03_C07_collecter_preuve.tex
% Titre    : Collecter la Preuve : OUI, NON, NEK
% Auteur   : MINKA MI NGUIDJOÏ Thierry Emmanuel
% Version  : 0.2 -- novembre 2026 (réécriture étoffée, alignée sur INF3226
%            C10, A0 et A1)
% ============================================================================
% Sources : A0 (hal-05701007), proposition de modalité d'ignorance explicite ;
% A1 (hal-05701270), instrument OUI/NON/NEK, mesures, indice rho ; GUID 5100 ;
% manuel IDI (documentation).
% Faits MINÉQUIR établis ici : domaine D1 Accès, 5 questions x 4 répondants
% (DAF EBOUMBOU, chef de la cellule ATANGANA, chef de la liquidation
% NDZIE, contrôleur financier MOUKOURI) ; grille binaire de l'audit interne
% 2025 : 17 OUI / 20 ; instrument trichotomique en mars 2026 : 11 OUI, 3 NON,
% 6 NEK -> mu 0,55, U 0,30, H 1,41 ; rho = 0,30.
% Faits CAPNO (LAB-05) : paie de décembre 2025, 4 312 lignes, net total
% 1 293 418 500 FCFA ; première extraction 4 309 lignes (3 manquantes,
% 1 905 000 FCFA) ; extraction refaite le 6 mai 2026.
% ============================================================================

\chapter[Collecter la preuve : OUI, NON, NEK]{Collecter la Preuve : OUI, NON, NEK}
\label{chap:preuve}
\chaptermark{Collecter la preuve}

\epigraph{%
    « Une grille qui interdit de répondre \emph{je ne sais pas}\\
    ne supprime pas l'ignorance~: elle la déguise en certitude. »%
}{Principe de l'auditologie}

\niveauChapitre{Intermédiaire}{Chapitre~\ref{chap:preparer}.}

\begin{boxobjectifs}
\begin{itemize}
    \item Utiliser l'instrument \OUI{}~/ \NON{}~/ \NEK{} auprès des
          répondants et pour coter les contrôles, et justifier la modalité
          \NEK{} à partir des axiomes de l'auditologie.
    \item Mesurer le biais d'une grille binaire par l'indice $\rho$.
    \item Classer les preuves selon leur fiabilité et pratiquer la
          corroboration.
    \item Conduire un entretien d'audit qui produit des réponses cotables.
    \item Calculer la taille d'un échantillon par attributs et en tirer une
          conclusion.
    \item Obtenir une extraction dont on garantit l'origine, l'exhaustivité,
          l'exactitude et l'intégrité.
    \item Tenir des papiers de travail qu'un tiers peut relire et refaire.
\end{itemize}
\end{boxobjectifs}

\bigskip

% ============================================================================
\section*{Pourquoi ce chapitre}
% ============================================================================

Dans le prologue, tout reposait sur deux fichiers et sur huit cotes. Si la
DAF avait pu dire «~ce fichier n'est pas le nôtre~», ou si la vérificatrice
avait coté \OUI{} ce qu'elle ne savait pas, le constat tombait. Ce chapitre
apprend à répondre à chaque question de la fiche d'assertion avec une
preuve qui tient, et à dire honnêtement ce que l'on ne sait pas. Il repose
sur l'instrument de réponse construit par l'article A1~\cite{minka2026nek},
à partir des axiomes de l'article A0~\cite{minka2026auditologie}.

\medskip
\textbf{Plan du chapitre.}
\begin{enumerate}
    \item L'instrument \OUI{}~/ \NON{}~/ \NEK{} (section~\ref{sec:c07-instrument}).
    \item Le biais du choix forcé (section~\ref{sec:c07-biais}).
    \item Les preuves et leur fiabilité (section~\ref{sec:c07-preuves}).
    \item Les entretiens (section~\ref{sec:c07-entretiens}).
    \item Tester et échantillonner (sections~\ref{sec:c07-techniques}
          et~\ref{sec:c07-echantillon}).
    \item Obtenir une extraction (section~\ref{sec:c07-extraction}).
    \item Les papiers de travail (section~\ref{sec:c07-papiers}).
\end{enumerate}

\separateur

% ============================================================================
\section{L'instrument OUI / NON / NEK}
\label{sec:c07-instrument}
% ============================================================================

\subsection{Pourquoi une troisième réponse}

L'article \emph{Vers l'auditologie} établit, à partir de quatre de ses
axiomes, que tout instrument de réponse d'audit valide doit comporter une
modalité explicite de connaissance insuffisante~\cite{minka2026auditologie}.
Le raisonnement est simple.
\begin{itemize}
    \item \textbf{Observabilité partielle} (axiome~1)~: l'auditeur ne voit
          qu'une partie des états du système~; certaines questions n'ont
          pas de réponse établie au moment où il les pose.
    \item \textbf{Cognition limitée} (axiome~2)~: les répondants, comme
          l'auditeur, ne savent pas tout de ce qu'ils gèrent.
    \item \textbf{Mesurabilité} (axiome~3)~: l'incertitude peut se mesurer,
          à condition d'être enregistrée.
    \item \textbf{Interaction stratégique} (axiome~4)~: un répondant a
          parfois intérêt à répondre \OUI{} quand il ne sait pas.
\end{itemize}
Une grille qui n'admet que \OUI{} et \NON{} oblige donc à transformer
l'ignorance en une réponse fausse. L'article A1 construit l'instrument qui
l'évite~\cite{minka2026nek}.

\begin{boxdef}
\termecle{Espace de réponse trichotomique}\index{espace de réponse
trichotomique}~: ensemble $\Omega = \{\OUI, \NON, \NEK\}$ des réponses
admises à une question d'audit.
\begin{itemize}
    \item \OUI{}~: le critère est satisfait, et une preuve l'établit.
    \item \NON{}~: le critère n'est pas satisfait, et une preuve l'établit.
    \item \NEK{} (\textit{Not Enough Knowledge})~: l'information disponible
          ne permet pas de conclure, malgré des diligences normales.
\end{itemize}
\end{boxdef}

\subsection{Deux usages}

\begin{enumerate}
    \item \textbf{Auprès des répondants}~: une même question est posée à
          plusieurs interlocuteurs. Pour chaque question, ou chaque domaine,
          on obtient une distribution de réponses, dont l'article A1 tire
          trois mesures~: $\mu$, la part des \OUI{}~; $U$, la part des
          \NEK{}~; $H$, l'entropie~\cite{minka2026nek}. Le
          Chapitre~\ref{chap:constater} apprend à les lire.
    \item \textbf{Par l'auditeur}~: après ses tests, l'auditeur cote chaque
          contrôle \OUI{}, \NON{} ou \NEK{}, preuve à l'appui. C'est la cote
          qui entre dans le rapport.
\end{enumerate}
Les deux usages se répondent~: une question où les répondants divergent
appelle des tests~; les tests permettent ensuite de coter avec preuve.

\subsection{NEK n'est pas « non applicable »}

\begin{boxalert}
\textbf{Violation typique.} Coter \NEK{} un contrôle qui ne s'applique pas au
périmètre (par exemple, la validation des annulations de recettes dans un
service qui n'encaisse rien).

\smallskip
\textbf{Règle~:} un contrôle sans objet est \emph{non applicable}~: il sort
de la décomposition, avec la raison écrite. \NEK{} est réservé à un contrôle
applicable dont on ne peut pas établir l'état. Les confondre fausse le taux
d'incertitude~$U$.
\end{boxalert}

\separateur

% ============================================================================
\section{Le biais du choix forcé}
\label{sec:c07-biais}
% ============================================================================

Une grille binaire force le répondant qui ne sait pas à choisir. Il choisit
le plus souvent la réponse attendue, ou la plus confortable~: le \OUI{}.
La grille binaire \emph{surestime} donc la conformité. L'article A1 mesure
cette surestimation par un indice de réduction du biais~\cite{minka2026nek}~:
\[
    \rho = \hat{\mu}^{B} - \mu^{T},
\]
où $\hat{\mu}^{B}$ est la part des \OUI{} obtenue avec l'instrument binaire, et
$\mu^{T}$ celle obtenue avec l'instrument trichotomique, sur les mêmes
questions et les mêmes répondants.

\begin{boxlecture}{Domaine D1, binaire puis trichotomique}
En 2025, l'audit interne du MINÉQUIR avait posé cinq questions sur les accès
à SIDÉPENSE à quatre interlocuteurs (le DAF, le chef de la cellule
informatique, le chef de la liquidation, le contrôleur financier), avec une
grille \OUI{}~/ \NON{}~: 17 \OUI{} sur 20 réponses. En mars~2026, la
vérificatrice pose les mêmes questions aux mêmes personnes avec
l'instrument trichotomique~: 11 \OUI{}, 3 \NON{}, 6 \NEK{}.

\medskip
\textbf{Constatation.} $\hat{\mu}^{B} = 17/20 = 0{,}85$~; $\mu^{T} = 11/20 =
0{,}55$~; $\rho = 0{,}30$.

\medskip
\textbf{Analyse} (confiance moyenne). La grille binaire surestimait la
conformité de trente points~: six réponses \OUI{} de 2025 étaient en
réalité des «~je ne sais pas~». C'est ainsi que l'audit interne avait pu
écrire que «~la gestion des accès est satisfaisante~», quelques mois avant
le double paiement.

\medskip
\hfill\lecture{0{,}55}{0{,}30}{1{,}41}
\end{boxlecture}

\begin{boxinfo}
\textbf{Ce que $\rho$ ne dit pas.} L'indice mesure l'écart entre deux
instruments, pas la conformité réelle. Seuls les tests de l'auditeur disent
si les contrôles fonctionnent~: au Chapitre~\ref{chap:acces}, un seul des
six contrôles d'accès sera coté \OUI{}.
\end{boxinfo}

\separateur

% ============================================================================
\section{Les preuves et leur fiabilité}
\label{sec:c07-preuves}
% ============================================================================

\subsection{Suffisante et appropriée}

\begin{boxdef}
\termecle{Preuve d'audit}\index{preuve d'audit}~: information sur laquelle
l'auditeur fonde sa cote et ses conclusions. Elle doit être
\emph{suffisante} (en quantité, au regard du risque) et \emph{appropriée}
(pertinente pour le critère, et fiable).
\end{boxdef}

\subsection{Une hiérarchie de fiabilité}

\begin{tableaucours}{cL{4.2cm}L{4.4cm}L{2.8cm}}{%
    \tableauHeader{Rang} &
    \tableauHeader{Source} &
    \tableauHeader{Exemple} &
    \tableauHeader{Fiabilité}
}
1 & Trace du système, extraite par l'auditeur & Matrice des profils
extraite devant la vérificatrice & La plus élevée \\
2 & Document externe & Relevé bancaire, visa du poste comptable & Élevée \\
3 & Document interne approuvé & Note de service signée & Moyenne \\
4 & Observation directe & La vérificatrice voit un agent liquider &
Ponctuelle \\
5 & Déclaration orale & « Les comptes sont désactivés dès qu'on nous
informe » & La plus faible \\
\end{tableaucours}
\vspace{-0.8em}
\begin{center}{\small\itshape Tableau~7.1. Fiabilité des preuves selon leur
source}\end{center}

Une déclaration orale ne suffit jamais à coter \OUI{}~: elle oriente la
recherche de la preuve. Quand l'auditeur ne dispose que d'une déclaration,
la cote est \NEK{}, et la question devient «~où est la preuve~?~».

\subsection{La corroboration}

Deux preuves indépendantes qui concordent valent plus que la plus fiable
des deux prise seule. Pour les contrôles importants, l'auditeur cherche une
seconde source~: la liste des mutations de la DRH (rang~3) recoupée avec
les dates de désactivation dans SIDÉPENSE (rang~1). Si deux sources se
contredisent sans qu'on puisse les départager, la cote est \NEK{}, et
l'écart devient lui-même un objet d'investigation.

\separateur

% ============================================================================
\section{Les entretiens}
\label{sec:c07-entretiens}
% ============================================================================

L'entretien d'audit, à la différence de l'entretien de prise de
connaissance (Chapitre~\ref{chap:preparer}), recueille des réponses
cotables. Ses questions sont celles des fiches d'assertion, classées par
interlocuteur.

\begin{boxprocedure}
\textbf{Règles de l'entretien d'audit}~:
\begin{enumerate}
    \item un répondant à la fois, pour éviter qu'un chef ne réponde pour son
          équipe~;
    \item lire la question de la fiche telle quelle, sans la reformuler en
          suggérant la réponse~;
    \item annoncer que \NEK{} est une réponse admise et utile~;
    \item noter la réponse, son auteur, la date, et la pièce invoquée~;
    \item demander la pièce séance tenante, ou fixer un délai écrit~;
    \item faire relire et signer le compte rendu.
\end{enumerate}
\end{boxprocedure}

\begin{boxalert}
\textbf{Violation typique.} Demander au DAF «~vous désactivez bien les
comptes des agents mutés, n'est-ce pas~?~».

\smallskip
\textbf{Règle~:} une question qui suggère la réponse produit un \OUI{} sans
valeur. On demande «~qui désactive, quand, sur quelle pièce~?~», et l'on
admet «~je ne sais pas~».
\end{boxalert}

\separateur

% ============================================================================
\section{Les techniques de test}
\label{sec:c07-techniques}
% ============================================================================

\begin{tableaucours}{L{2.6cm}L{4.4cm}L{5.2cm}}{%
    \tableauHeader{Technique} &
    \tableauHeader{Ce qu'elle établit} &
    \tableauHeader{Exemple}
}
Entretien & Ce que l'entité déclare & Le chef de la cellule décrit la
gestion des comptes \\
Observation & Ce qui se passe, au moment où l'on regarde & Assister à une
désactivation de compte \\
Inspection & Ce que contient un document ou un enregistrement & Lire la
matrice des profils \\
Ré-exécution & Que le contrôle produit le résultat attendu & Tenter de
saisir un mandat en double dans un environnement de test \\
Analyse de données & Ce que dit toute la population & Rechercher les
doublons parmi 9\,611 mandats \\
\end{tableaucours}
\vspace{-0.8em}
\begin{center}{\small\itshape Tableau~7.2. Les techniques de
test}\end{center}

\begin{boxalert}
\textbf{Violation typique.} Faire une ré-exécution dans l'application de
production («~pour voir si le doublon passe~»).

\smallskip
\textbf{Règle~:} l'auditeur ne crée, ne modifie ni ne supprime jamais une
donnée de production. Les essais se font dans un environnement de test, ou
par observation d'une opération réelle.
\end{boxalert}

\separateur

% ============================================================================
\section{Échantillonner dans le système}
\label{sec:c07-echantillon}
% ============================================================================

\subsection{Population et échantillon}

\begin{boxdef}
\termecle{Population}\index{population}~: ensemble des opérations sur
lesquelles porte la conclusion. Dans un SI, elle se prend dans le système,
après rapprochement de ses totaux (Chapitre~\ref{chap:prologue}).

\medskip
\termecle{Échantillon}\index{échantillon}~: partie de la population
examinée en détail. Il est \emph{aléatoire} quand chaque opération a la même
chance d'être tirée, \emph{raisonné} quand l'auditeur choisit les opérations
selon un critère (montant élevé, opération hors heures).
\end{boxdef}

L'échantillon aléatoire permet d'étendre la conclusion à la population~;
l'échantillon raisonné va droit aux anomalies probables, mais ne permet pas
cette extension. L'auditeur dit toujours lequel il a utilisé.

\subsection{L'échantillonnage par attributs, démontré}

On teste un attribut (le contrôle a-t-il été exécuté~?) sur $n$ opérations
tirées au hasard. On veut pouvoir dire, avec une confiance $c$, que le taux
réel d'écart ne dépasse pas un taux toléré $t$, si l'on ne trouve aucun
écart.

Si le taux réel d'écart était exactement $t$, la probabilité de ne trouver
aucun écart sur $n$ tirages indépendants serait $(1 - t)^n$. Pour que ce
résultat soit peu probable, au plus $1 - c$, il faut~:
\[
    (1 - t)^n \leq 1 - c
    \quad\Longleftrightarrow\quad
    n \geq \frac{\ln(1 - c)}{\ln(1 - t)}.
\]

\begin{tableaucours}{L{3.0cm}L{3.0cm}L{3.0cm}L{3.0cm}}{%
    \tableauHeader{Confiance $c$} &
    \tableauHeader{Taux toléré $t$} &
    \tableauHeader{Calcul} &
    \tableauHeader{Taille $n$}
}
95~\% & 5~\% & $\ln 0{,}05 / \ln 0{,}95 \approx 58{,}4$ & 59 \\
90~\% & 5~\% & $\ln 0{,}10 / \ln 0{,}95 \approx 44{,}9$ & 45 \\
95~\% & 10~\% & $\ln 0{,}05 / \ln 0{,}90 \approx 28{,}4$ & 29 \\
\end{tableaucours}
\vspace{-0.8em}
\begin{center}{\small\itshape Tableau~7.3. Tailles d'échantillon par
attributs, sans écart attendu}\end{center}

\subsection{Conclure sur l'échantillon}

\begin{itemize}
    \item \textbf{Aucun écart}~: avec une confiance $c$, le taux d'écart ne
          dépasse pas $t$~; le contrôle est coté \OUI{} pour la période.
    \item \textbf{Un écart ou plus}~: la conclusion par attributs tombe~; le
          contrôle est \NON{}, et l'on examine chaque écart (cause, montant).
    \item \textbf{Des éléments introuvables}~: s'ils sont nombreux, la cote
          est \NEK{}, et l'on écrit les diligences faites.
\end{itemize}

\begin{boxoutils}{Tirer un échantillon aléatoire dans un tableur}
Dans la feuille des mandats, ajouter une colonne~:
\begin{lstlisting}
H2 : =ALEA()
\end{lstlisting}
Copier la colonne et la coller en valeurs (pour figer le tirage), trier sur
H, retenir les $n$ premières lignes. Conserver la feuille triée~: elle
prouve le caractère aléatoire du tirage.
\end{boxoutils}

\subsection{Tester toute la population}

Quand la donnée est fiable, il est souvent plus probant de tester
\emph{toute} la population par une analyse~: doublons, trous de
numérotation, opérations hors heures, montants juste sous un seuil
(Chapitre~\ref{chap:donnees-fraude}). L'analyse désigne des opérations~;
chacune est ensuite examinée pièce par pièce.

\separateur

% ============================================================================
\section{Obtenir une extraction}
\label{sec:c07-extraction}
% ============================================================================

\subsection{Quatre garanties}

\begin{tableaucours}{L{2.8cm}L{4.4cm}L{5.0cm}}{%
    \tableauHeader{Garantie} &
    \tableauHeader{Question} &
    \tableauHeader{Comment l'obtenir}
}
Origine & Le fichier vient-il bien de l'application~? & Assister à
l'extraction, ou la faire faire sous ses yeux \\
Exhaustivité & Toutes les opérations y sont-elles~? & Requête écrite~;
nombre de lignes et totaux rapprochés d'une source indépendante \\
Exactitude & Les valeurs sont-elles celles de l'application~? & Comparer
quelques lignes à l'écran \\
Intégrité & Le fichier n'a-t-il pas changé depuis~? & Empreinte SHA-256
calculée à la remise, contresignée \\
\end{tableaucours}
\vspace{-0.8em}
\begin{center}{\small\itshape Tableau~7.4. Les quatre garanties d'une
extraction}\end{center}

\subsection{L'empreinte et les totaux de contrôle}

\begin{boxdef}
\termecle{Empreinte}\index{empreinte} (\textit{hash})~: suite de caractères
calculée à partir du contenu d'un fichier par une fonction de hachage
(SHA-256). Le moindre changement du fichier change l'empreinte.

\medskip
\termecle{Totaux de contrôle}\index{totaux de contrôle}~: nombre de lignes
et sommes, rapprochés d'une source que l'auditeur n'a pas obtenue de la même
personne (relevé bancaire, compte administratif, état du poste comptable).
\end{boxdef}

L'empreinte prouve que le fichier n'a pas changé~; elle ne prouve pas qu'il
était complet. Il faut les deux.

\begin{boxoutils}{Calculer et vérifier une empreinte}
Sous Windows, dans l'invite de commandes~:
\begin{lstlisting}
certutil -hashfile paie_2025-12.csv SHA256
\end{lstlisting}
Sous Linux (bureau d'analyse de NDAP SUKULU)~:
\begin{lstlisting}
sha256sum paie_2025-12.csv
sha256sum -c MANIFESTE_SHA256.txt
\end{lstlisting}
\end{boxoutils}

\begin{boxscenario}{Dossier CAPNO. Trois lignes de moins}
\textbf{Constatation.} La première extraction de la paie de décembre~2025,
remise sur une clé, compte 4\,309 lignes. L'ordre de virement transmis à la
banque porte sur 1\,293\,418\,500~FCFA pour 4\,312 bénéficiaires. Il manque
trois lignes, pour un total de 1\,905\,000~FCFA.

\medskip
\textbf{Analyse} (confiance élevée). L'extraction n'était pas exhaustive.
M.~HAMADOU la fait refaire devant lui le 6~mai~2026~: 4\,312 lignes, totaux
égaux à l'ordre de virement. Les trois lignes absentes deviennent les
premières à examiner~: pourquoi manquaient-elles~?
\end{boxscenario}

\separateur

% ============================================================================
\section{Les papiers de travail}
\label{sec:c07-papiers}
% ============================================================================

Le manuel de l'IDI le rappelle~: la documentation doit permettre à un
auditeur expérimenté, sans lien avec la mission, de comprendre les travaux,
les preuves et les conclusions~\cite{idi2021conformite}. C'est aussi la
traduction pratique de l'axiome d'\emph{opposabilité}~: une conclusion qui
ne peut pas être refaite ne peut pas être défendue~\cite{minka2026auditologie}.

\begin{boxprocedure}
\textbf{Feuille de test}, une par contrôle testé~:
\begin{enumerate}
    \item référence (position dans le gabarit, ligne du programme)~;
    \item fiche d'assertion (critères et questions)~;
    \item population (source, nombre, totaux, empreinte du fichier)~;
    \item échantillon (méthode, confiance, taux toléré, taille, liste)~;
    \item résultat élément par élément, avec la pièce examinée~;
    \item cote de chaque critère et cote du contrôle, avec la
          constatation~; pour un \NEK{}, les diligences et la relance~;
    \item auteur, date, et visa du chef de mission.
\end{enumerate}
\end{boxprocedure}

\begin{boxinfo}
\textbf{Conserver.} Les extractions sont des données personnelles (paie) ou
sensibles (fournisseurs). Elles sont conservées sur un support protégé, avec
leur empreinte, aussi longtemps que les suites possibles l'exigent, puis
détruites selon les règles de l'institution.
\end{boxinfo}

\begin{boxmetier}{Vérificateur}
Annoncez \NEK{} dès le premier entretien~: «~si vous ne savez pas,
dites-le~; c'est une réponse utile~». Vous obtiendrez moins de \OUI{}, et
davantage de vérité.
\end{boxmetier}

\begin{boxmetier}{Juge des comptes}
Une pièce électronique produite à l'instance vaut si l'on peut en établir
l'origine et l'intégrité. Demandez le procès-verbal de remise et
l'empreinte~; leur absence affaiblit la pièce.
\end{boxmetier}

\begin{boxmetier}{Auditeur interne}
Remplacez vos grilles \OUI{}~/ \NON{} par l'instrument à trois réponses, et
mesurez $\rho$ sur vos propres questionnaires~: vous saurez de combien vos
évaluations passées surestimaient la conformité.
\end{boxmetier}

\separateur

\begin{boxresume}
\begin{itemize}
    \item \textbf{Trois réponses, fondées en théorie.} \OUI{}, \NON{}, \NEK{}~;
          la modalité \NEK{} découle des axiomes~1 à~4 d'A0.
    \item \textbf{NEK n'est pas « non applicable ».} Un contrôle sans objet
          sort de la décomposition.
    \item \textbf{La grille binaire surestime}~; $\rho$ mesure de combien
          (0,30 sur les accès du MINÉQUIR).
    \item \textbf{Une hiérarchie de preuves}~; une déclaration seule donne
          \NEK{}~; deux sources indépendantes valent mieux qu'une.
    \item \textbf{L'entretien ne suggère pas}~; il admet «~je ne sais pas~».
    \item \textbf{$n \ge \ln(1-c) / \ln(1-t)$}~: 59 tirages pour 95~\% de
          confiance et 5~\% d'écart toléré.
    \item \textbf{Empreinte et totaux}~: intégrité et exhaustivité d'une
          extraction.
    \item \textbf{Un tiers doit pouvoir refaire}~: c'est l'opposabilité.
\end{itemize}
\end{boxresume}

\begin{boxexo}[title={\ding{45}\quad Exercice~7.1. Coter avec la bonne preuve \niveaubase}]
Pour chaque situation, donnez la cote et le rang de la preuve
(Tableau~7.1)~: (a) le DAF affirme que les profils sont revus chaque
semestre~; (b) la matrice extraite devant vous montre cinq comptes LIQ +
ORD~; (c) une note signée prévoit la désactivation sous 48~heures, mais le
journal n'a pas encore été remis~; (d) la DRH et la cellule donnent deux
dates de désactivation différentes pour le même agent.
\end{boxexo}

\begin{boxexo}[title={\ding{45}\quad Exercice~7.2. L'indice $\rho$ \niveaubase}]
Sur le domaine D4 (exploitation), la grille binaire de 2025 donnait 14
\OUI{} sur 20~; l'instrument trichotomique donne 8 \OUI{}, 4 \NON{}, 8
\NEK{}. Calculez $\hat{\mu}^{B}$, $\mu^{T}$ et $\rho$, et dites ce que cela
signifie.
\end{boxexo}

\begin{boxexo}[title={\ding{45}\quad Exercice~7.3. Dimensionner \niveaumoyen}]
Vous voulez conclure, avec 95~\% de confiance, que le visa du contrôle
financier manque sur moins de 3~\% des mandats de 2025. (a) Combien de
mandats tirez-vous~? (b) Vous trouvez un mandat sans visa~: que
concluez-vous~? (c) Que proposez-vous plutôt, puisque l'extraction des
9\,611 mandats est disponible~?
\end{boxexo}

\begin{boxexo}[title={\ding{45}\quad Exercice~7.4. Trois lignes \niveauexpert}]
Les trois lignes absentes de la première extraction de la CAPNO
correspondent à trois agents créés le 28~novembre~2025 par le même compte.
Dites ce que cela suggère, ce que vous ne pouvez pas encore affirmer, les
trois pièces que vous demandez, et la cote provisoire du contrôle
«~l'extraction remise est exhaustive~».
\end{boxexo}

\begin{boxinfo}
\textbf{Transition.} Les réponses sont recueillies, les contrôles cotés,
les preuves conservées. Le Chapitre~\ref{chap:constater} apprend à lire
ces réponses par trois mesures, à conclure, à rapporter et à suivre.
\end{boxinfo}
}{}
        \IfFileExists{partie3/P03_C08_constater_rapporter_suivre.tex}{% ============================================================================
% Fichier  : partie3/P03_C08_constater_rapporter_suivre.tex
% Titre    : Mesurer, Conclure, Rapporter, Suivre
% Auteur   : MINKA MI NGUIDJOÏ Thierry Emmanuel
% Version  : 0.2 -- novembre 2026 (réécriture étoffée, alignée sur INF3226
%            C11 et C12, A0 et A1)
% ============================================================================
% Sources : A1 (hal-05701270) : mu, U, H, cas limites, quadrant, divergence ;
% A0 (hal-05701007) : non-conformité connue / inconnue, axiomes 5 et 6 ;
% Shannon (1948), Kullback et Leibler (1951), Cover et Thomas (2006).
% Faits MINÉQUIR (répondants, mars 2026, 20 réponses par domaine) :
%   D1 Accès 11/3/6 (0,55 ; 0,30 ; 1,41) conformité opaque ;
%   D2 Fournisseurs 6/12/2 (0,30 ; 0,10 ; 1,30) non-conformité claire ;
%   D3 Circuit du mandat 15/3/2 (0,75 ; 0,10 ; 1,05) conformité claire ;
%   D4 Exploitation 8/4/8 (0,40 ; 0,40 ; 1,52) non-conformité opaque.
% Audit de suivi de novembre 2026, D1 : 17/2/1 (0,85 ; 0,05 ; 0,75) ;
%   D_KL = 0,346 bit (0,247 avec lissage de Laplace).
% Rapport provisoire 17/04/2026 ; réponse DAF 30/04/2026 ; définitif
% 15/05/2026 ; recommandations R-01 à R-06 (LAB-10).
% ============================================================================

\chapter[Mesurer, conclure, rapporter, suivre]{Mesurer, Conclure, Rapporter, Suivre}
\label{chap:constater}
\chaptermark{Conclure, rapporter, suivre}

\epigraph{%
    « On ne lit jamais l'entropie sans l'incertitude,\\
    ni l'incertitude sans la conformité. »%
}{Règle du cours, d'après A1}

\niveauChapitre{Intermédiaire}{Chapitre~\ref{chap:preuve}.}

\begin{boxobjectifs}
\begin{itemize}
    \item Calculer et interpréter les trois mesures $\mu$, $U$ et $H$ d'un
          domaine, et en connaître les cas limites.
    \item Décomposer l'entropie en part d'ignorance et part de désaccord.
    \item Situer un domaine dans le quadrant conformité~$\times$~incertitude
          et en déduire l'action.
    \item Appliquer la règle de décision du cours~: $U$ d'abord, $\mu$
          ensuite, $H$ pour contrôler.
    \item Rédiger une observation en sept éléments, et la relier au constat
          en cinq éléments.
    \item Conduire la procédure contradictoire et rédiger le rapport.
    \item Suivre les recommandations et comparer deux cycles par la
          divergence de Kullback-Leibler.
    \item Préserver la preuve quand le constat peut fonder une suite.
\end{itemize}
\end{boxobjectifs}

\bigskip

% ============================================================================
\section*{Pourquoi ce chapitre}
% ============================================================================

Les cotes \OUI{}, \NON{} et \NEK{} sont des résultats de travail. Pour
conclure, il faut les lire ensemble, domaine par domaine, sans confondre ce
que l'on sait et ce que l'on ignore. L'article A1 donne pour cela trois
mesures et un quadrant~\cite{minka2026nek}~; l'article A0, la distinction
entre non-conformité connue et non-conformité
inconnue~\cite{minka2026auditologie}. Ce chapitre les applique au dossier
MINÉQUIR, puis suit le rapport jusqu'au plan d'action et au cycle suivant.

\medskip
\textbf{Plan du chapitre.}
\begin{enumerate}
    \item Trois mesures (section~\ref{sec:c08-mesures}) et la décomposition de
          l'entropie (section~\ref{sec:c08-decomp}).
    \item Le quadrant et la règle de décision (sections~\ref{sec:c08-quadrant}
          et~\ref{sec:c08-regle}).
    \item Le constat et l'observation (section~\ref{sec:c08-constat}).
    \item Le rapport et la procédure contradictoire
          (sections~\ref{sec:c08-rapport} et~\ref{sec:c08-contradictoire}).
    \item Suivre les recommandations et comparer les cycles
          (sections~\ref{sec:c08-suivi} et~\ref{sec:c08-cycles}).
    \item Quand le constat devient une affaire (section~\ref{sec:c08-suites}).
\end{enumerate}

\separateur

% ============================================================================
\section{Trois mesures}
\label{sec:c08-mesures}
% ============================================================================

\subsection{Les définitions}

Pour un ensemble de réponses (une question posée à plusieurs répondants, ou
toutes les questions d'un domaine), on note $p(\OUI)$, $p(\NON)$ et
$p(\NEK)$ les parts de chaque réponse.

\begin{boxdef}
\termecle{Taux de conformité}\index{taux de conformité}~: $\mu = p(\OUI)$.

\medskip
\termecle{Taux d'incertitude}\index{taux d'incertitude}~: $U = p(\NEK)$.

\medskip
\termecle{Entropie}\index{entropie}~: dispersion des réponses, en bits,
\[
    H = -\sum_{m \in \{\OUI, \NON, \NEK\}} p(m) \log_2 p(m),
    \qquad 0 \leq H \leq \log_2 3 \approx 1{,}585,
\]
avec la convention $0 \log_2 0 = 0$.
\end{boxdef}
{\small\itshape Définitions de l'article A1~\cite{minka2026nek}~;
l'entropie est celle de Shannon~\cite{shannon1948}.}

\medskip

L'entropie est nulle quand tous répondent la même chose, maximale quand les
trois réponses sont également fréquentes. Au niveau d'un domaine, on agrège
les réponses de toutes ses questions avant de calculer~: les mesures se
recalculent toujours à partir des réponses, jamais comme moyenne des
mesures du niveau inférieur.

\subsection{Quatre cas limites}

\begin{tableaucours}{L{4.6cm}cccL{3.4cm}}{%
    \tableauHeader{Situation} &
    \tableauHeader{$\mu$} & \tableauHeader{$U$} & \tableauHeader{$H$} &
    \tableauHeader{Lecture}
}
Tous répondent \OUI{} & 1 & 0 & 0 & Conformité certaine \\
Tous répondent \NON{} & 0 & 0 & 0 & Non-conformité certaine \\
Tous répondent \NEK{} & 0 & 1 & 0 & \textbf{Ignorance unanime} \\
Réponses également réparties & $1/3$ & $1/3$ & 1,585 & Dispersion maximale \\
\end{tableaucours}
\vspace{-0.8em}
\begin{center}{\small\itshape Tableau~8.1. Les cas limites des trois
mesures (d'après~\cite{minka2026nek})}\end{center}

\begin{boxalert}
\textbf{Violation typique.} Conclure «~certitude~» parce que l'entropie est
nulle.

\smallskip
\textbf{Règle~:} une entropie nulle signifie seulement que tous ont donné la
même réponse. Si cette réponse est \NEK{}, personne ne sait rien~: $H = 0$ et
$U = 1$. On ne lit jamais $H$ sans $U$ ni $\mu$.
\end{boxalert}

\subsection{Les quatre domaines du MINÉQUIR}

En mars~2026, la vérificatrice a posé cinq questions par domaine aux quatre
interlocuteurs du Chapitre~\ref{chap:preuve}, soit vingt réponses par
domaine.

\begin{tableaucours}{L{3.6cm}cccccc}{%
    \tableauHeader{Domaine} &
    \tableauHeader{\OUI{}} & \tableauHeader{\NON{}} & \tableauHeader{\NEK{}} &
    \tableauHeader{$\mu$} & \tableauHeader{$U$} & \tableauHeader{$H$}
}
D1 Accès & 11 & 3 & 6 & 0,55 & 0,30 & 1,41 \\
D2 Fichier des fournisseurs & 6 & 12 & 2 & 0,30 & 0,10 & 1,30 \\
D3 Circuit du mandat & 15 & 3 & 2 & 0,75 & 0,10 & 1,05 \\
D4 Exploitation & 8 & 4 & 8 & 0,40 & 0,40 & 1,52 \\
\end{tableaucours}
\vspace{-0.8em}
\begin{center}{\small\itshape Tableau~8.2. Les trois mesures des quatre
domaines (réponses des répondants, mars~2026)}\end{center}

\begin{boxoutils}{Calculer les trois mesures}
Dans un tableur, avec les effectifs en B2 (\OUI{}), C2 (\NON{}), D2
(\NEK{})~:
\begin{lstlisting}
E2 (mu) : =B2/SOMME(B2:D2)
F2 (U)  : =D2/SOMME(B2:D2)
G2 (H)  : =-SOMME(SI(B2:D2>0;B2:D2/SOMME(B2:D2)*LOG(B2:D2/SOMME(B2:D2);2);0))
\end{lstlisting}
La dernière formule se valide comme formule matricielle dans les tableurs
qui l'exigent.
\end{boxoutils}

\separateur

% ============================================================================
\section{Décomposer l'entropie}
\label{sec:c08-decomp}
% ============================================================================

Une entropie élevée peut venir de deux choses très différentes~: beaucoup de
répondants ne savent pas (ignorance), ou ceux qui savent ne sont pas
d'accord (désaccord entre \OUI{} et \NON{}). On ne traite pas l'ignorance
comme le désaccord.

On regroupe les réponses en deux classes~: \NEK{} (probabilité $U$) et
«~réponse informée~» (\OUI{} ou \NON{}, probabilité $1 - U$). Dans la seconde
classe, la part des \OUI{} vaut $\mu / (1 - U)$. La règle de chaînage de
l'entropie~\cite{shannon1948,cover2006elements} donne, pour $U < 1$~:
\[
    H = \underbrace{H_b(U)}_{\text{ignorance}} +
        \underbrace{(1 - U)\, H_b\!\left(\frac{\mu}{1 - U}\right)}_{\text{désaccord}},
    \qquad H_b(x) = -x\log_2 x - (1-x)\log_2(1-x).
\]

\begin{tableaucours}{L{3.6cm}ccccL{2.8cm}}{%
    \tableauHeader{Domaine} &
    \tableauHeader{$H$} &
    \tableauHeader{$H_b(U)$} &
    \tableauHeader{Désaccord} &
    \tableauHeader{Part d'ignorance} &
    \tableauHeader{Lecture}
}
D1 Accès & 1,41 & 0,88 & 0,53 & 63~\% & Surtout ignorance \\
D2 Fournisseurs & 1,30 & 0,47 & 0,83 & 36~\% & Surtout désaccord \\
D4 Exploitation & 1,52 & 0,97 & 0,55 & 64~\% & Surtout ignorance \\
\end{tableaucours}
\vspace{-0.8em}
\begin{center}{\small\itshape Tableau~8.3. Décomposition de l'entropie de
trois domaines}\end{center}

\begin{boxlecture}{Deux entropies voisines, deux démarches}
D1 et D2 ont des entropies proches (1,41 et 1,30), pour des raisons
opposées.

\medskip
\textbf{Analyse} (confiance élevée). Dans D1, près des deux tiers de
l'entropie viennent de l'ignorance~: personne ne sait vraiment qui
désactive les comptes~; on cherche la preuve. Dans D2, la dispersion vient
du désaccord~: le DAF croit les RIB contrôlés, le chef de la liquidation
sait qu'ils ne le sont pas~; on confronte les deux versions à la trace. Le
même chiffre appelle deux démarches.

\medskip
\hfill\lecture{0{,}55}{0{,}30}{1{,}41}
\end{boxlecture}

\separateur

% ============================================================================
\section{Le quadrant conformité × incertitude}
\label{sec:c08-quadrant}
% ============================================================================

L'article A1 propose de lire ensemble la conformité et l'incertitude dans un
quadrant à quatre régions, chacune appelant une action
différente~\cite{minka2026nek}. Le cours retient, par convention, un seuil de
0,5 pour $\mu$ et de 0,20 pour $U$~; chaque mission peut fixer les siens, à
condition de les écrire dans le programme de travail.

\noindent\begin{minipage}{\linewidth}
\begin{center}
\begin{tikzpicture}[scale=0.95]
\fill[green!18] (3,0) rectangle (6,2);
\fill[red!15] (0,0) rectangle (3,2);
\fill[orange!25] (3,2) rectangle (6,5);
\fill[orange!15] (0,2) rectangle (3,5);
\draw[thick, color=CouleurPrimaire] (0,0) rectangle (6,5);
\draw[dashed, color=CouleurPrimaire] (3,0) -- (3,5);
\draw[dashed, color=CouleurPrimaire] (0,2) -- (6,2);
\node[align=center, font=\footnotesize\bfseries] at (1.5,0.45) {Non-conformité claire};
\node[align=center, font=\footnotesize\bfseries] at (4.5,0.45) {Conformité claire};
\node[align=center, font=\footnotesize\bfseries] at (1.5,4.6) {Non-conformité opaque};
\node[align=center, font=\footnotesize\bfseries] at (4.5,4.6) {Conformité opaque};
\foreach \x/\y/\t in {3.3/3.0/D1, 1.8/1.0/D2, 4.5/1.0/D3, 2.4/4.0/D4}
    \fill[CouleurAccent] (\x,\y) circle (2.5pt) node[above right, font=\scriptsize\bfseries] {\t};
\node[font=\scriptsize] at (3,-0.3) {0,5};
\node[font=\scriptsize] at (-0.35,2) {0,20};
\node[font=\small] at (3,-0.8) {Taux de conformité $\mu$};
\node[font=\small, rotate=90] at (-0.9,2.5) {Taux d'incertitude $U$};
\end{tikzpicture}
\end{center}
\vspace{-0.4em}
\begin{center}{\small\itshape Figure~8.1. Le quadrant diagnostique et les
quatre domaines du MINÉQUIR (d'après~\cite{minka2026nek}~; seuils du cours~;
échelle verticale $U \times 10$)}\end{center}
\end{minipage}

\medskip

\begin{tableaucours}{L{3.4cm}L{4.0cm}L{4.8cm}}{%
    \tableauHeader{Région} &
    \tableauHeader{Situation} &
    \tableauHeader{Action}
}
Conformité claire & On s'accorde, avec assurance, sur la conformité &
Validation courante \\
Non-conformité claire & On s'accorde, avec assurance, sur la
non-conformité & Correctifs \\
\textbf{Conformité opaque} & Des \OUI{} nombreux, mais une part importante
de \NEK{} & \textbf{Alerte prioritaire}~: rechercher la preuve avant toute
conclusion \\
Non-conformité opaque & Des \NON{} nombreux et une part importante de
\NEK{} & Rechercher la preuve et engager des premiers correctifs \\
\end{tableaucours}
\vspace{-0.8em}
\begin{center}{\small\itshape Tableau~8.4. Les quatre régions et leurs
actions (d'après~\cite{minka2026nek})}\end{center}

La conformité opaque est la région la plus dangereuse~: une grille binaire
la confond avec la conformité claire. C'est le domaine D1 du MINÉQUIR, celui
du double paiement. L'article A0 dit la même chose dans son
vocabulaire~: la \emph{non-conformité connue} appelle des correctifs~; la
\emph{non-conformité inconnue}, que cachent les régions opaques, appelle
d'abord des investigations~\cite{minka2026auditologie}.

\separateur

% ============================================================================
\section{La règle de décision du cours}
\label{sec:c08-regle}
% ============================================================================

\begin{boxprocedure}
\textbf{Conclure sur un domaine.}
\begin{enumerate}
    \item \textbf{$U$ d'abord.} Si $U \geq 0{,}20$, le domaine est opaque~:
          aucune conclusion de conformité n'est possible. L'auditeur
          recherche la preuve, ou conclut \emph{sur l'opacité elle-même}.
    \item \textbf{$\mu$ ensuite.} Si $U < 0{,}20$, le domaine est clair~:
          conforme si $\mu \geq 0{,}5$, non conforme sinon.
    \item \textbf{$H$ pour contrôler.} Une entropie élevée avec $U$ faible
          signale un désaccord~: le domaine n'est pas homogène~; on le
          décompose avant de conclure (Chapitre~\ref{chap:preparer}).
    \item \textbf{Les cotes de l'auditeur prévalent.} Les mesures des
          répondants orientent~; la conclusion repose sur les contrôles
          cotés par l'auditeur, preuve à l'appui (Chapitre~\ref{chap:preuve}).
\end{enumerate}
\end{boxprocedure}

\begin{boxlecture}{Les quatre domaines du MINÉQUIR}
\textbf{Analyse} (confiance moyenne).
\begin{itemize}
    \item \emph{D1}~: conformité opaque~; alerte prioritaire. Les tests
          (Chapitre~\ref{chap:acces}) établiront un seul contrôle sur six
          conforme~: l'opacité cachait une non-conformité.
    \item \emph{D2}~: non-conformité claire~; correctifs immédiats sur les
          RIB.
    \item \emph{D3}~: conformité claire selon les répondants~; mais
          l'auditeur a coté \NON{} le blocage des doublons (règle~4)~: la
          conclusion suit les cotes.
    \item \emph{D4}~: non-conformité opaque~; investigations et premiers
          correctifs.
\end{itemize}
\end{boxlecture}

\begin{boxalert}
\textbf{Violation typique.} Moyenner les quatre domaines pour donner une
note globale («~le circuit est conforme à 50~\%~»).

\smallskip
\textbf{Règle~:} une moyenne efface les régions du quadrant. Le rapport
présente chaque domaine dans sa région.
\end{boxalert}

\separateur

% ============================================================================
\section{Le constat en cinq éléments}
\label{sec:c08-constat}
% ============================================================================

\subsection{Le constat}

\begin{tableaucours}{L{2.6cm}L{4.4cm}L{5.2cm}}{%
    \tableauHeader{Élément} &
    \tableauHeader{Question} &
    \tableauHeader{Règle de rédaction}
}
Condition & Qu'a-t-on constaté~? & Un fait vérifiable, chiffré, daté, avec
sa source et sa cote \\
Critère & Que fallait-il~? & La règle citée exactement, avec sa force
(Chapitre~\ref{chap:normes}) \\
Cause & Pourquoi l'écart~? & Ce qui, dans le dispositif, l'explique~; la
cause appelle la recommandation \\
Effet & Qu'est-ce que cela produit ou permet~? & Le risque réalisé ou
possible, chiffré quand on le peut \\
Recommandation & Que faire~? & Une action, un responsable, une échéance \\
\end{tableaucours}
\vspace{-0.8em}
\begin{center}{\small\itshape Tableau~8.5. Les cinq éléments d'un
constat}\end{center}

\subsection{L'observation en sept éléments}

La méthode du cours d'audit de l'auteur développe le constat en une
\termecle{observation}\index{observation} en sept éléments. Les deux
éléments ajoutés sont indispensables à l'opposabilité (axiome~6
d'A0~\cite{minka2026auditologie})~: la trace du contradictoire, et une
conclusion explicite qui sépare la position de l'auditeur de son analyse.

\begin{tableaucours}{L{5.0cm}L{6.0cm}}{%
    \tableauHeader{Constat en cinq éléments} &
    \tableauHeader{Observation en sept éléments}
}
Critère & 1. Norme \\
Condition & 2. Constatation (avec la cote) \\
Effet & 3. Risques \\
(absent) & 4. Contradictoire \\
Cause & 5. Analyse (avec un niveau de confiance) \\
(absent) & 6. Conclusion \\
Recommandation & 7. Recommandations \\
\end{tableaucours}
\vspace{-0.8em}
\begin{center}{\small\itshape Tableau~8.6. Du constat à
l'observation}\end{center}

\begin{boxobjet}{D1.1.2 · Observation}
\textbf{1. Norme~:} le compte d'un agent est désactivé à la date d'effet de
sa mutation (bonne pratique~: COSO~2013, principe~11~; aucune règle
interne).

\textbf{2. Constatation~:} au 23~mars~2026, 9 des 31 agents mutés ou partis
en 2025 ont un compte actif, dont \texttt{mbarga.j}, qui cumule LIQ, ORD et
FRS (extractions et empreintes au dossier). Cote \NON{}.

\textbf{3. Risques~:} un compte resté actif a servi, les 27 et
28~décembre~2025, à modifier un RIB et à émettre seul un mandat de
187\,450\,000~FCFA.

\textbf{4. Contradictoire~:} le 30~avril~2026, la DAF indique que les
mutations ne lui sont pas notifiées par les ressources humaines.

\textbf{5. Analyse} (confiance élevée)~: la cause est l'absence de flux
entre le fichier du personnel et la cellule informatique, et l'absence de
procédure~; non la négligence d'un agent.

\textbf{6. Conclusion~:} le contrôle D1.1.2 n'est pas efficace~;
déficience majeure du domaine D1 (Chapitre~\ref{chap:eval-sci}).

\textbf{7. Recommandations~:} désactiver les 9 comptes~; organiser la
notification écrite des mutations par la DRH~; imposer dans SIDÉPENSE
l'incompatibilité LIQ~/ ORD.
\end{boxobjet}

\begin{boxalert}
\textbf{Violation typique.} Écrire dans les risques~: «~M.~MBARGA a détourné
187\,450\,000~FCFA~».

\smallskip
\textbf{Règle~:} l'observation dit ce que la preuve établit~: un compte, pas
une personne. L'imputation relève d'une autre procédure.
\end{boxalert}

\subsection{Écrire une recommandation}

\begin{tableaucours}{L{5.6cm}L{6.6cm}}{%
    \tableauHeader{Mal écrite} &
    \tableauHeader{Bien écrite}
}
Renforcer la sécurité des accès & Organiser la notification écrite des
mutations de la DRH vers la cellule, avant le 31~juillet~2026 (responsable~:
DRH) \\
Désactiver les comptes inutiles & Désactiver sous 48~heures les 9 comptes
identifiés (annexe~C), puis revoir chaque semestre les comptes actifs
(responsable~: DAF) \\
Sanctionner M.~MBARGA & (Hors du rôle de l'auditeur~: la sanction relève
de l'autorité compétente, après enquête.) \\
\end{tableaucours}
\vspace{-0.8em}
\begin{center}{\small\itshape Tableau~8.7. Écrire une
recommandation}\end{center}

\separateur

% ============================================================================
\section{Le rapport}
\label{sec:c08-rapport}
% ============================================================================

\begin{enumerate}
    \item \textbf{Synthèse}~: l'objectif, la conclusion sur l'efficacité du
          contrôle interne, les domaines dans leur région du quadrant, les
          déficiences majeures, en une page.
    \item \textbf{Mandat et étendue}~: fondement, périmètre, période,
          travaux, seuils de décision, et \termecle{limitations
          d'étendue}\index{limitation d'étendue} (les \NEK{} restants, avec
          les diligences faites).
    \item \textbf{Mesures par domaine}~: $\mu$, $U$, $H$ et région.
    \item \textbf{Observations}, classées par domaine puis par gravité.
    \item \textbf{Tableau des déficiences} et leur classement.
    \item \textbf{Réponses de l'entité} et position de l'auditeur.
    \item \textbf{Annexes}~: matrice des risques et des contrôles,
          décomposition, liste des pièces et de leurs empreintes.
\end{enumerate}

Un \NEK{} qui reste n'est pas caché~: la restitution impartiale l'exige, et
la conclusion ne peut pas porter sur ce qui n'a pas pu être établi.

\separateur

% ============================================================================
\section{La procédure contradictoire}
\label{sec:c08-contradictoire}
% ============================================================================

\subsection{Pourquoi le contradictoire}

Le rapport provisoire est adressé à l'entité, qui répond par écrit dans un
délai fixé. Le contradictoire n'est pas une politesse~: il réduit
l'incertitude. L'audité connaît des faits que l'auditeur ignore~; sa
réponse peut lever un \NEK{}, corriger une cause, ou révéler une pièce. Il
est aussi la condition de l'opposabilité~: une conclusion que l'audité n'a
pas pu discuter se défend mal devant le CDBF ou la Chambre des Comptes.

\begin{tableaucours}{L{4.4cm}L{4.0cm}L{3.8cm}}{%
    \tableauHeader{Réponse de la DAF (30 avril 2026)} &
    \tableauHeader{Position de la vérificatrice} &
    \tableauHeader{Effet sur le rapport}
}
«~Les mutations ne nous ont pas été notifiées par les ressources
humaines.~» & Exact~; c'est la cause relevée~; ne change pas la condition
& Observation maintenue, cause précisée \\
«~Les 9 comptes ont été désactivés le 20~avril.~» & Pièce fournie
(journal) & Mention au rapport~; R-01 mise en œuvre \\
«~La séparation a toujours été respectée en pratique.~» & Contredit par
le mandat 2025-09342 (journal) & Observation maintenue \\
\end{tableaucours}
\vspace{-0.8em}
\begin{center}{\small\itshape Tableau~8.8. Le contradictoire du dossier
MINÉQUIR}\end{center}

\subsection{Maintenir, modifier, retirer}

Pour chaque réponse, l'auditeur décide, et motive~: il \emph{maintient}
l'observation, la \emph{modifie} au vu d'une pièce nouvelle, ou la
\emph{retire}. Une déclaration sans pièce ne fait jamais changer une cote
\NON{}~; elle peut, au mieux, transformer une cause.

\begin{boxprocedure}
\textbf{Fiche de contradictoire}, une par observation~: date d'envoi du
rapport provisoire~; réponse reçue (date, signataire, texte)~; pièces
nouvelles~; position motivée~; modification apportée au rapport définitif.
\end{boxprocedure}

\separateur

% ============================================================================
\section{Suivre les recommandations}
\label{sec:c08-suivi}
% ============================================================================

\begin{tableaucours}{cL{5.0cm}L{2.6cm}L{2.0cm}L{1.8cm}}{%
    \tableauHeader{N°} &
    \tableauHeader{Recommandation} &
    \tableauHeader{Responsable} &
    \tableauHeader{Échéance} &
    \tableauHeader{Statut}
}
R-01 & Désactiver les 9 comptes & Cellule informatique & 20/04/2026 &
Fait \\
R-02 & Imposer l'incompatibilité LIQ / ORD dans SIDÉPENSE & Cellule,
éditeur & 30/06/2026 & En cours \\
R-03 & Double validation des RIB & DAF, cellule & 30/06/2026 & En cours
\\
R-04 & Bloquer les doublons (fournisseur, montant, facture normalisée) &
Cellule, éditeur & 30/06/2026 & En cours \\
R-05 & Notification des mutations par les ressources humaines & DRH &
31/07/2026 & Non commencé \\
R-06 & Revue des mandats de fin de gestion & DAF, contrôle financier &
31/12/2026 & Non commencé \\
\end{tableaucours}
\vspace{-0.8em}
\begin{center}{\small\itshape Tableau~8.9. Plan d'action du MINÉQUIR
(situation au 15~mai~2026)}\end{center}

Le suivi vérifie la mise en œuvre par une preuve, pas par une
déclaration~: pour R-02, la vérificatrice demandera une tentative, en
environnement de test, d'attribuer LIQ et ORD au même compte.

\begin{boxinfo}
\textbf{L'audit modifie ce qu'il observe.} L'article A0 énonce une
proposition de \emph{réflexivité}~: l'acte d'audit fait partie de la
trajectoire du système audité~\cite{minka2026auditologie}. Le plan d'action
change le système~; le cycle suivant n'audite donc pas le même objet. C'est
pourquoi on compare des distributions, et non des taux isolés.
\end{boxinfo}

\separateur

% ============================================================================
\section{Comparer deux cycles}
\label{sec:c08-cycles}
% ============================================================================

\subsection{Pourquoi le taux de conformité ne suffit pas}

Le taux de conformité seul peut rester stable alors que l'état de
connaissance change~: des \NEK{} deviennent des \NON{}, l'entité sait
désormais ce qu'elle ignorait. Il faut comparer les \emph{distributions} de
réponses. C'est la conséquence de l'\emph{individuation temporelle}
(axiome~5)~: les conclusions d'une mission ne se testent que sur plusieurs
cycles~\cite{minka2026auditologie}.

\begin{boxdef}
\termecle{Divergence de Kullback-Leibler}\index{divergence de
Kullback-Leibler}~: pour deux distributions $p$ (nouveau cycle) et $q$
(cycle précédent) sur $\{\OUI, \NON, \NEK\}$~\cite{kullback1951,minka2026nek},
\[
    D_{\mathrm{KL}}(p \,\|\, q) = \sum_{m} p(m) \log_2 \frac{p(m)}{q(m)}.
\]
Elle est positive ou nulle, nulle si et seulement si les deux distributions
sont identiques~; elle n'est pas symétrique (on mesure toujours le nouveau
cycle par rapport à l'ancien)~; elle s'exprime en bits.
\end{boxdef}

\begin{boxscenario}{Dossier MINÉQUIR. Le domaine D1 en novembre 2026}
\textbf{Faits.} En mars~2026, D1 donnait 11~\OUI{}, 3~\NON{}, 6~\NEK{}. À
l'audit de suivi de novembre~2026, après R-01 à R-03, les mêmes questions
posées aux mêmes fonctions donnent 17~\OUI{}, 2~\NON{}, 1~\NEK{}.

\medskip
\textbf{Calcul.} $\mu$ passe de 0,55 à 0,85, $U$ de 0,30 à 0,05, $H$ de 1,41
à 0,75~bit. La divergence vaut
\[
    0{,}85\log_2\frac{0{,}85}{0{,}55} + 0{,}10\log_2\frac{0{,}10}{0{,}15} +
    0{,}05\log_2\frac{0{,}05}{0{,}30} \approx 0{,}534 - 0{,}058 - 0{,}129
    \approx 0{,}35~\text{bit}.
\]

\textbf{Analyse} (confiance moyenne). D1 est passé de la conformité opaque
à la conformité claire. La vérificatrice confirmera par ses propres tests
(revue des comptes actifs, essai de cumul en environnement de test) avant
de conclure.
\end{boxscenario}

\subsection{Le lissage}

Si une réponse est absente au premier cycle et présente au second, un terme
de la divergence devient infini~: ce n'est pas que le changement soit
infini, c'est que le premier échantillon, petit, n'a rien observé. On
ajoute alors une unité fictive à chaque réponse des deux cycles (lissage de
Laplace). La divergence de D1 devient 0,25~bit. On applique le même lissage
à toutes les comparaisons d'une mission, et on l'écrit dans le programme.

\begin{boxoutils}{Calculer la divergence (Python)}
\begin{lstlisting}[language=Python]
from math import log2
def dkl(nouveau, ancien, lissage=0):
    a = [x + lissage for x in nouveau]
    b = [x + lissage for x in ancien]
    p = [x / sum(a) for x in a]
    q = [x / sum(b) for x in b]
    return sum(pi * log2(pi / qi) for pi, qi in zip(p, q) if pi > 0)

print(round(dkl([17, 2, 1], [11, 3, 6]), 3))     # 0.346
print(round(dkl([17, 2, 1], [11, 3, 6], 1), 3))  # 0.247 (lissée)
\end{lstlisting}
\end{boxoutils}

\separateur

% ============================================================================
\section{Quand le constat devient une affaire}
\label{sec:c08-suites}
% ============================================================================

Une observation peut fonder une suite~: saisine du CDBF, que le CONSUPE peut
engager (décret n°~2013/287, art.~3 et~5)~; transmission à la Chambre des
Comptes~; plainte pénale. L'auditeur n'instruit pas ces procédures, mais il
les rend possibles ou impossibles par la qualité de sa preuve.

\begin{itemize}
    \item \textbf{Geler la preuve}~: conserver les extractions originales,
          leurs empreintes et les procès-verbaux~; travailler sur des copies.
    \item \textbf{Ne pas alerter}~: ne pas interroger la personne soupçonnée
          sur les faits avant que l'autorité ait décidé de la suite.
    \item \textbf{Demander la conservation} des journaux au-delà de leur
          durée habituelle.
    \item \textbf{Séparer}~: le rapport traite du contrôle interne~; une note
          distincte, adressée à la hiérarchie, présente les faits
          susceptibles de poursuite.
\end{itemize}

\begin{boxmetier}{Vérificateur}
Présentez dans la synthèse les domaines par région du quadrant, et non par
une note globale. Une «~conformité opaque~» écrite en toutes lettres
protège votre rapport mieux qu'un taux flatteur.
\end{boxmetier}

\begin{boxmetier}{Juge des comptes}
Une limitation d'étendue écrite au rapport dit ce que l'auditeur n'a pas pu
établir~; elle désigne les pièces que l'instruction peut exiger.
\end{boxmetier}

\begin{boxmetier}{Auditeur interne}
Refaites chaque année les mêmes questions aux mêmes fonctions, et calculez
la divergence entre les cycles~: c'est la mesure la plus honnête de
l'effet de vos recommandations.
\end{boxmetier}

\separateur

\begin{boxresume}
\begin{itemize}
    \item \textbf{Trois mesures}~: $\mu$ (conformité), $U$ (incertitude), $H$
          (entropie, au plus $\log_2 3$)~; jamais l'une sans les autres.
    \item \textbf{L'entropie se décompose} en ignorance et désaccord~; chacun
          appelle sa démarche.
    \item \textbf{Le quadrant} sépare quatre régions~; la conformité opaque
          est la plus dangereuse.
    \item \textbf{La règle de décision}~: $U$ d'abord, $\mu$ ensuite, $H$
          pour contrôler~; les cotes de l'auditeur prévalent.
    \item \textbf{Sept éléments}~: l'observation ajoute au constat le
          contradictoire et la conclusion, pour l'opposabilité.
    \item \textbf{Le contradictoire réduit l'incertitude}~; une déclaration
          sans pièce ne fait pas changer un \NON{}.
    \item \textbf{Comparer les cycles} par la divergence de Kullback-Leibler,
          avec le même lissage partout.
    \item \textbf{Préserver la suite}~: geler la preuve, ne pas alerter,
          séparer rapport et note de faits.
\end{itemize}
\end{boxresume}

\begin{boxexo}[title={\ding{45}\quad Exercice~8.1. Calculer et situer \niveaubase}]
Calculez $\mu$, $U$ et $H$ pour (a) 18~\OUI{}, 1~\NON{}, 1~\NEK{}~; (b)
0~\OUI{}, 0~\NON{}, 20~\NEK{}~; (c) 10~\OUI{}, 8~\NON{}, 2~\NEK{}. Situez
chacun dans le quadrant et dites l'action.
\end{boxexo}

\begin{boxexo}[title={\ding{45}\quad Exercice~8.2. Ignorance ou désaccord~? \niveaumoyen}]
Décomposez l'entropie du domaine D3 (15~\OUI{}, 3~\NON{}, 2~\NEK{}) et
dites quelle démarche elle appelle. Expliquez pourquoi la conclusion du
domaine reste \NON{} sur le blocage des doublons.
\end{boxexo}

\begin{boxexo}[title={\ding{45}\quad Exercice~8.3. L'observation \niveaumoyen}]
Rédigez l'observation en sept éléments sur le contrôle «~seconde validation
des modifications de RIB~» (domaine D2), à partir des faits du prologue et
de la réponse de la DAF~: «~les RIB sont modifiés sur demande écrite du
fournisseur~» (aucune demande n'est produite pour BATIRURAL).
\end{boxexo}

\begin{boxexo}[title={\ding{45}\quad Exercice~8.4. Deux cycles \niveauexpert}]
À l'audit de suivi, D4 donne 8~\OUI{}, 10~\NON{}, 2~\NEK{} (contre 8, 4, 8
en mars). (a) Le taux de conformité a-t-il changé~? (b) Calculez la
divergence. (c) Dites si l'exploitation s'est améliorée, et rédigez la
phrase du rapport de suivi.
\end{boxexo}

\begin{boxinfo}
\textbf{Transition.} La méthode est complète, du mandat au cycle suivant. La
Partie~IV l'applique domaine par domaine, en commençant par celui dont tout
dépend~: les accès (Chapitre~\ref{chap:acces}).
\end{boxinfo}
}{}

    % ======================================================================
    % PARTIE 4 — DOMAINES D'AUDIT (J05 à J07)
    % « Chaque domaine a ses contrôles, chaque contrôle a sa trace. »
    % Standards de référence : GUID 5100, COBIT 2019, ISO/IEC 27002,
    %   guide COSO sur la fraude ; lois n° 2010/012 et 2024/017.
    % Niveau : intermédiaire à expert
    % ======================================================================
    \part{Domaines d'Audit}
    \label{part:domaines}

        \IfFileExists{partie4/P04_C09_acces_habilitations.tex}{% ============================================================================
% Fichier  : partie4/P04_C09_acces_habilitations.tex
% Titre    : Accès et Habilitations
% Auteur   : MINKA MI NGUIDJOÏ Thierry Emmanuel
% Version  : 0.1 -- novembre 2026
% ============================================================================
% Sources : INF3226 C18, C21 (accès) ; COSO principe 11 ; ISO/IEC 27002:2022
% (gestion des accès, résumé) ; ACSI LAB-02 (matrice des droits).
% Faits MINÉQUIR établis ici (LAB-03) : 47 comptes actifs au 23/03/2026 ;
% 9 comptes d'agents partis ; 3 comptes génériques (saisie01, daf, admin) ;
% 5 comptes LIQ + ORD (dont mbarga.j) ; 2 comptes FRS + ORD (mbarga.j,
% nkolo.r) ; profil ADM : atangana.p et admin (générique) ; mot de passe
% sans expiration ; dernière revue des droits : aucune.
% ============================================================================

\chapter{Accès et Habilitations}
\label{chap:acces}

\epigraph{%
    « Un droit d'accès est une signature que l'on peut prêter. »%
}{Principe du cours}

\niveauChapitre{Intermédiaire}{Chapitres~\ref{chap:eval-sci}
et~\ref{chap:preuve}.}

\begin{boxobjectifs}
\begin{itemize}
    \item Décrire le cycle de vie d'un compte et les contrôles attendus à
          chaque étape.
    \item Repérer les comptes à risque~: partis, génériques, à privilèges,
          inactifs.
    \item Croiser la liste des comptes avec le fichier du personnel.
    \item Analyser une matrice de profils et identifier les conflits de
          séparation des tâches.
    \item Évaluer la revue périodique des habilitations.
    \item Coter les contrôles d'accès et classer les déficiences.
\end{itemize}
\end{boxobjectifs}

\bigskip

% ============================================================================
\section*{Pourquoi ce chapitre}
% ============================================================================

Les accès sont le premier des contrôles généraux informatiques, et le
plus important~: s'ils ne sont pas maîtrisés, aucun autre contrôle n'est
fiable (Chapitre~\ref{chap:sci}). Le double paiement du prologue est
d'abord une défaillance des accès. Ce chapitre donne la grille complète
pour les évaluer.

\separateur

% ============================================================================
\section{Le cycle de vie d'un compte}
\label{sec:c09-cycle}
% ============================================================================

\begin{tableaucours}{L{2.4cm}L{4.8cm}L{5.0cm}}{%
    \tableauHeader{Étape} &
    \tableauHeader{Contrôle attendu} &
    \tableauHeader{Preuve à demander}
}
Création & Demande écrite, validée par le chef de service, pour un agent
identifié & Formulaires de demande, liste des créations \\
Attribution des profils & Profils conformes à la fiche de poste~; pas
de combinaison interdite & Matrice des profils, demandes \\
Modification & Nouvelle demande validée lors d'un changement de poste &
Historique des profils \\
Revue périodique & Le propriétaire de l'application valide la liste des
comptes et des profils & Revues signées, avec suites \\
Désactivation & À la date d'effet du départ ou de la mutation & Journal
des désactivations, rapproché du fichier du personnel \\
\end{tableaucours}
\vspace{-0.8em}
\begin{center}{\small\itshape Tableau~9.1. Le cycle de vie d'un
compte}\end{center}

\begin{boxdef}
\termecle{Propriétaire d'application}\index{propriétaire d'application}~:
responsable métier qui décide des droits dans une application (pour
SIDÉPENSE, le directeur des affaires financières). La cellule
informatique \emph{exécute} les décisions du propriétaire~; elle ne les
prend pas.
\end{boxdef}

\separateur

% ============================================================================
\section{Les comptes à risque}
\label{sec:c09-risque}
% ============================================================================

\begin{tableaucours}{L{2.8cm}L{4.4cm}L{5.0cm}}{%
    \tableauHeader{Type} &
    \tableauHeader{Risque} &
    \tableauHeader{Test}
}
Compte d'un agent parti & Utilisé par l'ancien titulaire ou par un
tiers & Croiser avec le fichier du personnel \\
Compte générique (\texttt{saisie01}) & Aucune action n'est imputable à
une personne & Lister les comptes sans titulaire nominatif \\
Compte à privilèges (ADM) & Peut tout faire, y compris effacer ses traces
& Lister les détenteurs, vérifier la journalisation de leurs actions \\
Compte inactif & Oublié, donc non surveillé & Dernière connexion de plus
de 90 jours \\
\end{tableaucours}
\vspace{-0.8em}
\begin{center}{\small\itshape Tableau~9.2. Les comptes à risque}\end{center}

\begin{boxoutils}{Croiser les comptes avec le fichier du personnel}
Deux feuilles dans le même classeur~: \texttt{COMPTES} (colonne A
matricule) et \texttt{PERSONNEL} (colonne A matricule, colonne D
situation). Dans \texttt{COMPTES}~:
\begin{lstlisting}
E2 : =SIERREUR(RECHERCHEV(A2;PERSONNEL!$A:$D;4;FAUX);"ABSENT")
\end{lstlisting}
Un filtre sur \texttt{E} donne les comptes d'agents mutés ou partis
(situation autre que «~en poste~») et les comptes sans titulaire au
fichier (\texttt{ABSENT}), qui sont souvent des comptes génériques.
\end{boxoutils}

\begin{boxscenario}{Dossier MINÉQUIR. Les 47 comptes}
\textbf{Constatation.} Au 23~mars~2026, SIDÉPENSE compte 47 comptes
actifs. Le croisement avec le fichier du personnel donne 9 comptes
d'agents mutés ou partis et 3 comptes sans titulaire~:
\texttt{saisie01}, \texttt{daf} et \texttt{admin}. Le profil ADM est
détenu par \texttt{atangana.p} et par le compte générique
\texttt{admin}. Les mots de passe n'expirent jamais.

\medskip
\textbf{Analyse} (confiance élevée). Le compte \texttt{admin} est le plus
grave~: il peut tout faire, et rien de ce qu'il fait n'est imputable à une
personne. Le journal du prologue ne serait d'aucune aide si les
opérations du 28~décembre avaient été faites sous ce compte.
\end{boxscenario}

\separateur

% ============================================================================
\section{La matrice des profils}
\label{sec:c09-matrice}
% ============================================================================

L'analyse part de la matrice extraite par l'auditeur (un compte par ligne,
un profil par colonne) et des paires interdites du Tableau~2.6.

\begin{boxoutils}{Repérer les conflits dans la matrice}
Colonnes \texttt{B} à \texttt{F}~: ENG, LIQ, ORD, FRS, ADM (1 si le
compte détient le profil, 0 sinon). Une colonne par paire interdite~:
\begin{lstlisting}
G2 (LIQ+ORD) : =SI(ET(C2=1;D2=1);"CONFLIT";"")
H2 (FRS+ORD) : =SI(ET(E2=1;D2=1);"CONFLIT";"")
I2 (FRS+LIQ) : =SI(ET(E2=1;C2=1);"CONFLIT";"")
J2 (ADM+...) : =SI(ET(F2=1;SOMME(B2:E2)>0);"CONFLIT";"")
\end{lstlisting}
\end{boxoutils}

\begin{boxscenario}{Dossier MINÉQUIR. Les conflits}
\textbf{Constatation.} Cinq comptes réunissent LIQ et ORD, dont
\texttt{mbarga.j}~; deux réunissent FRS et ORD (\texttt{mbarga.j} et
\texttt{nkolo.r}). Aucune revue des habilitations n'a été faite depuis la
mise en service de SIDÉPENSE en 2024.

\medskip
\textbf{Analyse} (confiance élevée). Le conflit LIQ + ORD n'est pas un
cas isolé~: il touche cinq comptes, ce qui montre que la séparation n'a
jamais été conçue dans l'attribution des profils. L'absence de revue
explique que personne ne l'ait vu.
\end{boxscenario}

\separateur

% ============================================================================
\section{La revue des habilitations}
\label{sec:c09-revue}
% ============================================================================

\begin{boxdef}
\termecle{Revue des habilitations}\index{revue des habilitations}~:
contrôle détectif par lequel le propriétaire de l'application examine,
à intervalle fixé, la liste des comptes et de leurs profils, confirme ou
retire chaque droit, et signe. Elle n'est efficace que si la liste est
exhaustive et extraite de l'application, et si les retraits décidés sont
exécutés.
\end{boxdef}

\begin{boxobjet}{C-ACC-03 · Revue des habilitations}
\textbf{Objectif.} Les droits dans SIDÉPENSE restent conformes aux
fonctions.

\textbf{Risque couvert.} Accumulation de droits, comptes oubliés.

\textbf{Critère.} Revue semestrielle signée par le DAF (bonne pratique~;
aucune règle interne).

\textbf{Test de conception.} Une procédure de revue existe-t-elle~?
Non.

\textbf{Test de fonctionnement.} Sans objet~: aucune revue n'a eu lieu.

\textbf{Cote.} \NON{}.
\end{boxobjet}

\separateur

% ============================================================================
\section{L'authentification}
\label{sec:c09-authentification}
% ============================================================================

L'authentification relie un compte à une personne. Les contrôles attendus
sont simples~: un compte par personne, un mot de passe personnel, une
longueur minimale, un changement périodique ou une authentification à
deux facteurs pour les profils sensibles, un verrouillage après plusieurs
échecs. Ces règles relèvent aussi de l'audit de sécurité de l'ANTIC
(Chapitre~\ref{chap:institutions})~: le vérificateur les note sans
conduire lui-même un audit de sécurité.

\begin{boxalert}
\textbf{Violation typique.} Constater que les mots de passe n'expirent
jamais, et conclure que la sécurité de SIDÉPENSE est défaillante.

\smallskip
\textbf{Règle~:} rester dans son objectif. Pour le contrôle interne de la
dépense, le constat est que l'imputation des opérations aux personnes
n'est pas garantie~; la sécurité d'ensemble relève de l'audit de
sécurité.
\end{boxalert}

\separateur

% ============================================================================
\section{Lecture du domaine}
\label{sec:c09-lecture}
% ============================================================================

\begin{tableaucours}{L{6.2cm}cL{4.4cm}}{%
    \tableauHeader{Contrôle} &
    \tableauHeader{Cote} &
    \tableauHeader{Pièce}
}
Création sur demande validée & \NON{} & Aucune demande archivée \\
Profils sans combinaison interdite & \NON{} & Matrice (5 et 2 conflits) \\
Désactivation à la date du départ & \NON{} & 9 comptes sur 31 \\
Revue des habilitations & \NON{} & Aucune revue \\
Comptes nominatifs & \NON{} & 3 comptes génériques \\
Journalisation des actions des comptes ADM & \OUI{} & Journal \\
\end{tableaucours}
\vspace{-0.8em}
\begin{center}{\small\itshape Tableau~9.3. Lecture du domaine D1
(accès)}\end{center}

\begin{boxlecture}{Domaine D1, accès à SIDÉPENSE}
\textbf{Constatation.} Un contrôle sur six fonctionne.

\textbf{Analyse} (confiance élevée). Le domaine est au niveau~0 ou~1 de
l'échelle du Chapitre~\ref{chap:eval-sci}. La déficience est majeure~:
elle prive de valeur tous les contrôles d'application de SIDÉPENSE.

\medskip
\hfill\conformite{1}{6}{0}
\end{boxlecture}

\begin{boxmetier}{Vérificateur}
Demandez toujours trois listes extraites devant vous~: les comptes et
leurs profils, la dernière connexion de chaque compte, et les
détenteurs des profils d'administration. Croisez-les avec le fichier du
personnel du même jour.
\end{boxmetier}

\begin{boxmetier}{Contrôleur financier}
Votre visa électronique est un profil. Si un compte générique ou un autre
agent peut viser à votre place, votre contrôle n'existe pas~: signalez-le
par écrit.
\end{boxmetier}

\begin{boxmetier}{Auditeur interne}
Proposez à la direction une revue semestrielle des habilitations, signée
par le propriétaire de l'application~: c'est le contrôle le moins coûteux
du domaine, et celui qui aurait détecté les cinq conflits.
\end{boxmetier}

\separateur

\begin{boxresume}
\begin{itemize}
    \item \textbf{Cinq étapes, cinq contrôles.} Création, attribution,
          modification, revue, désactivation.
    \item \textbf{Le propriétaire décide, la cellule exécute.}
    \item \textbf{Quatre comptes à risque.} Partis, génériques, à
          privilèges, inactifs.
    \item \textbf{Croiser avec le personnel}, le même jour, sur des listes
          extraites devant l'auditeur.
    \item \textbf{La matrice révèle les conflits}~; la revue les fait
          disparaître.
    \item \textbf{Rester dans son objectif.} L'authentification se note~;
          l'audit de sécurité relève de l'ANTIC.
\end{itemize}
\end{boxresume}

\begin{boxexo}[title={\ding{45}\quad Exercice~9.1. Classer les comptes \niveaubase}]
Classez selon le Tableau~9.2~: (a) \texttt{admin}~; (b) \texttt{mbarga.j}
au 23~mars~2026~; (c) \texttt{saisie01}~; (d) un compte dont la dernière
connexion date de juin~2025.
\end{boxexo}

\begin{boxexo}[title={\ding{45}\quad Exercice~9.2. Une formule \niveaubase}]
Écrivez la formule qui signale, dans la matrice, un compte qui détient à
la fois ENG et ORD, puis dites si cette combinaison doit être interdite.
\end{boxexo}

\begin{boxexo}[title={\ding{45}\quad Exercice~9.3. La revue qui ne sert à rien \niveaumoyen}]
À la CAPNO, une revue des habilitations est signée chaque semestre. Elle
porte sur une liste saisie à la main par l'informaticien, et aucun retrait
n'a jamais été décidé. Cotez le contrôle et rédigez le constat.
\end{boxexo}

\begin{boxexo}[title={\ding{45}\quad Exercice~9.4. Un petit service \niveauexpert}]
La recette des finances de Bafoussam n'a que trois agents pour
e-RECETTES. Proposez une attribution des profils, les contrôles
compensatoires nécessaires et la preuve que chacun doit laisser.
\end{boxexo}

\begin{boxinfo}
\textbf{Transition.} Les accès maîtrisés, il faut encore que les
programmes et leurs paramètres ne changent pas sans contrôle, et que
l'exploitation tienne. C'est l'objet du Chapitre~\ref{chap:exploitation}.
\end{boxinfo}
}{}
        \IfFileExists{partie4/P04_C10_changements_exploitation_continuite.tex}{% ============================================================================
% Fichier  : partie4/P04_C10_changements_exploitation_continuite.tex
% Titre    : Changements, Exploitation et Continuité
% Auteur   : MINKA MI NGUIDJOÏ Thierry Emmanuel
% Version  : 0.1 -- novembre 2026
% ============================================================================
% Sources : INF3226 C23 (production et continuité, resserré) ; COBIT 2019
% (BAI06, DSS01, DSS04 : cités par leur intitulé général) ; ISO 22301 ;
% loi n° 2010/012 (art. 32).
% Faits e-RECETTES établis ici (LAB-08) : 14 mises en production en 2025,
% 11 sans demande écrite, 3 avec tests documentés ; 1 modification directe
% du taux de pénalité en base par le prestataire le 9 septembre 2025 ;
% journal des sauvegardes 2025 : 365 jours, 23 échecs, plus longue suite
% d'échecs 6 jours (du 14 au 19 août 2025) ; aucune restauration testée ;
% sauvegardes sur un disque dans la salle du serveur ; pas de PCA ; aucun
% audit de sécurité ANTIC.
% ============================================================================

\chapter{Changements, Exploitation et Continuité}
\label{chap:exploitation}
\chaptermark{Changements, exploitation, continuité}

\epigraph{%
    « Une sauvegarde que l'on n'a jamais restaurée est un espoir, pas un
    contrôle. »%
}{Adage des exploitants}

\niveauChapitre{Intermédiaire}{Chapitre~\ref{chap:acces}.}

\begin{boxobjectifs}
\begin{itemize}
    \item Évaluer la gestion des changements~: demande, test, approbation,
          mise en production, trace.
    \item Repérer une modification directe des données ou des paramètres
          hors de toute procédure.
    \item Évaluer l'exploitation~: traitements, incidents, journaux.
    \item Tester les sauvegardes par leur journal et par une restauration.
    \item Évaluer le plan de continuité et le situer par rapport à l'audit
          de sécurité de l'ANTIC.
\end{itemize}
\end{boxobjectifs}

\bigskip

% ============================================================================
\section*{Pourquoi ce chapitre}
% ============================================================================

Un contrôle automatique n'est fiable que si personne ne peut le modifier
sans contrôle. Une donnée n'est conservée que si l'on peut la restaurer.
Ce chapitre évalue les deux autres familles de contrôles généraux
informatiques, sur le dossier e-RECETTES de Bafoussam\footnote{La recette
des finances, le prestataire NUMÉRIKA~SARL, les dates et les montants de
ce dossier sont fictifs~; les mécanismes sont exacts.}.

\separateur

% ============================================================================
\section{La gestion des changements}
\label{sec:c10-changements}
% ============================================================================

\subsection{Le circuit attendu}

\begin{boxdef}
\termecle{Changement}\index{changement}~: toute modification d'un
programme, d'un paramètre (taux, plafond, barème) ou de la structure d'une
base, en environnement de production.

\medskip
\termecle{Gestion des changements}\index{gestion des changements}~:
ensemble des contrôles qui garantissent qu'un changement est demandé,
analysé, testé, approuvé par le propriétaire de l'application, mis en
production par une autre personne que le développeur, et tracé.
\end{boxdef}

\begin{tableaucours}{L{2.6cm}L{5.0cm}L{4.6cm}}{%
    \tableauHeader{Étape} &
    \tableauHeader{Contrôle} &
    \tableauHeader{Preuve}
}
Demande & Écrite, motivée, émise par le métier & Ticket ou fiche \\
Test & En environnement de test, par le métier & Procès-verbal de
recette \\
Approbation & Par le propriétaire de l'application & Signature, date \\
Mise en production & Par une personne distincte du développeur & Journal
des mises en production \\
Trace & Toute modification de production est journalisée & Journal de la
base \\
\end{tableaucours}
\vspace{-0.8em}
\begin{center}{\small\itshape Tableau~10.1. Le circuit d'un
changement}\end{center}

\subsection{Le test}

On rapproche la liste des mises en production (journal technique) de la
liste des demandes (tickets). Toute mise en production sans demande
approuvée est un écart~; toute modification de la base hors mise en
production est un écart grave.

\begin{boxscenario}{Dossier e-RECETTES. Quatorze mises en production}
\textbf{Constatation.} En 2025, le journal compte 14 mises en production
d'e-RECETTES. Onze n'ont aucune demande écrite~; trois seulement ont un
procès-verbal de tests. Le 9~septembre~2025, le journal de la base montre
une modification directe du taux de pénalité de retard, faite par le
compte du prestataire NUMÉRIKA, sans mise en production.

\medskip
\textbf{Analyse} (confiance élevée). Le paramétrage d'e-RECETTES peut être
modifié par le prestataire, sans demande et sans approbation. Tout
contrôle automatique de l'application, y compris le calcul des
pénalités, perd donc sa valeur probante jusqu'à ce que l'on sache ce qui a
été modifié, et quand.
\end{boxscenario}

\separateur

% ============================================================================
\section{L'exploitation}
\label{sec:c10-exploitation}
% ============================================================================

L'exploitation fait tourner le système au quotidien~: traitements
programmés (calculs de fin de mois, transferts vers la trésorerie),
surveillance, gestion des incidents, journaux.

\begin{tableaucours}{L{3.0cm}L{4.6cm}L{4.6cm}}{%
    \tableauHeader{Contrôle} &
    \tableauHeader{Ce que l'on vérifie} &
    \tableauHeader{Preuve}
}
Traitements & Ils s'exécutent comme prévu~; les échecs sont repris &
Journal des traitements, comptes rendus \\
Incidents & Ils sont enregistrés, analysés, clos & Registre des
incidents \\
Journaux & Ils sont activés, protégés et conservés & Paramètres,
échantillon de journaux \\
Interfaces & Les transferts sont complets & Totaux de contrôle à
l'émission et à la réception \\
\end{tableaucours}
\vspace{-0.8em}
\begin{center}{\small\itshape Tableau~10.2. Les contrôles
d'exploitation}\end{center}

\separateur

% ============================================================================
\section{Les sauvegardes}
\label{sec:c10-sauvegardes}
% ============================================================================

\begin{boxobjet}{C-EXP-01 · Sauvegardes et restauration}
\textbf{Objectif.} Les données d'e-RECETTES peuvent être restaurées après
une perte.

\textbf{Risque couvert.} Perte de données de recettes (panne, erreur,
attaque).

\textbf{Critère.} Sauvegarde quotidienne réussie, copie hors du site,
restauration testée au moins une fois par an (bonne pratique~; aucune
règle interne).

\textbf{Test de conception.} La copie est sur un disque dans la salle du
serveur~: pas de copie hors du site.

\textbf{Test de fonctionnement.} Journal des sauvegardes 2025~: 365 jours,
23 échecs, dont 6 consécutifs du 14 au 19~août~2025. Aucune restauration
n'a jamais été tentée.

\textbf{Cote.} \NON{}.
\end{boxobjet}

\begin{boxoutils}{Analyser un journal de sauvegardes}
Le journal donne une ligne par jour (colonne A date, colonne B statut
\texttt{OK} ou \texttt{ECHEC}). Pour compter les échecs et la plus
longue suite~:
\begin{lstlisting}
Nombre d'echecs  : =NB.SI(B:B;"ECHEC")
C2 (suite)       : =SI(B2="ECHEC";C1+1;0)
Plus longue suite: =MAX(C:C)
\end{lstlisting}
La plus longue suite d'échecs mesure la perte maximale de données que
l'entité a courue sans le savoir.
\end{boxoutils}

\separateur

% ============================================================================
\section{La continuité}
\label{sec:c10-continuite}
% ============================================================================

\begin{boxdef}
\termecle{Plan de continuité d'activité}\index{plan de continuité}
(\sigle{PCA})~: document qui fixe comment l'entité continue ses missions
essentielles pendant une interruption, et comment elle revient à la
normale. Il fixe deux durées~: l'interruption maximale admissible, et la
perte de données maximale admissible.
\end{boxdef}

Le vérificateur ne conduit pas un audit de continuité selon ISO~22301~;
il vérifie trois choses~: qu'un plan existe pour les applications
financières, qu'il a été testé, et que ses durées sont compatibles avec
les sauvegardes réelles. Pour la sécurité d'ensemble, il demande la date
et le plan d'action du dernier audit de sécurité de l'ANTIC (loi
n°~2010/012, art.~32)~; son absence est un constat
(Chapitre~\ref{chap:institutions}).

\begin{boxscenario}{Dossier e-RECETTES. Ni plan, ni audit}
\textbf{Constatation.} La recette des finances n'a pas de plan de
continuité. Aucun audit de sécurité de l'ANTIC n'a été réalisé depuis la
mise en service d'e-RECETTES. Les données sont hébergées par NUMÉRIKA
hors du pays (Chapitre~\ref{chap:gouvernance}).

\medskip
\textbf{Analyse} (confiance élevée). Une panne du serveur d'hébergement
arrêterait l'encaissement des recettes sans solution de repli, et la
dernière sauvegarde utilisable pourrait dater de plusieurs jours.
\end{boxscenario}

\separateur

\begin{boxlecture}{Domaine D4, changements et exploitation d'e-RECETTES}
\textbf{Constatation.} Sur sept contrôles attendus (demande écrite,
test, approbation, absence de modification directe, journalisation,
sauvegardes, plan de continuité), un seul fonctionne~: la
journalisation, qui a permis de voir la modification du 9~septembre.

\textbf{Analyse} (confiance élevée). Déficience majeure~: le
paramétrage n'est pas maîtrisé et la restauration n'est pas garantie.

\medskip
\hfill\conformite{1}{7}{0}
\end{boxlecture}

\begin{boxmetier}{Vérificateur}
Demandez le journal des mises en production et la liste des demandes de
la même période, puis rapprochez-les ligne à ligne. Une mise en
production sans demande se voit en quelques minutes.
\end{boxmetier}

\begin{boxmetier}{Juge des comptes}
Si un taux ou un barème a été modifié en cours d'exercice sans
procédure, les montants calculés par l'application avant et après la
modification ne reposent pas sur la même règle~: demandez la date exacte
et l'acte qui l'autorisait.
\end{boxmetier}

\begin{boxmetier}{Auditeur interne}
Proposez un test de restauration annuel, en présence d'un témoin du
métier, avec procès-verbal. C'est la seule preuve qu'une sauvegarde
sert à quelque chose.
\end{boxmetier}

\separateur

\begin{boxresume}
\begin{itemize}
    \item \textbf{Un changement suit un circuit.} Demande, test,
          approbation, mise en production par un tiers, trace.
    \item \textbf{Rapprocher les mises en production des demandes.} Ce qui
          n'a pas de demande est un écart.
    \item \textbf{La modification directe en base} est l'écart le plus
          grave~: elle prive les contrôles automatiques de valeur.
    \item \textbf{L'exploitation se prouve par ses journaux}~: traitements,
          incidents, interfaces.
    \item \textbf{Une sauvegarde se prouve par une restauration}, et se
          mesure par la plus longue suite d'échecs.
    \item \textbf{Continuité et sécurité.} Plan testé~; audit de sécurité
          de l'ANTIC demandé, pas refait.
\end{itemize}
\end{boxresume}

\begin{boxexo}[title={\ding{45}\quad Exercice~10.1. Le circuit \niveaubase}]
Pour la modification du taux de pénalité du 9~septembre~2025, dites
quelles étapes du Tableau~10.1 ont manqué et quelle preuve aurait dû
exister pour chacune.
\end{boxexo}

\begin{boxexo}[title={\ding{45}\quad Exercice~10.2. Le journal des sauvegardes \niveaubase}]
Un journal de 30 jours contient la suite~: 10 \texttt{OK}, 3
\texttt{ECHEC}, 5 \texttt{OK}, 4 \texttt{ECHEC}, 8 \texttt{OK}. Donnez le
nombre d'échecs, la plus longue suite, et ce qu'elle signifie.
\end{boxexo}

\begin{boxexo}[title={\ding{45}\quad Exercice~10.3. Les effets d'un changement \niveaumoyen}]
Le taux de pénalité est passé de 10~\% à 5~\% le 9~septembre. Dites quels
tests vous faites sur les recettes encaissées après cette date, et ce que
vous écrivez dans l'effet du constat.
\end{boxexo}

\begin{boxexo}[title={\ding{45}\quad Exercice~10.4. Classer \niveauexpert}]
Classez les déficiences du domaine D4 d'e-RECETTES (Tableau~3.7 du
Chapitre~\ref{chap:eval-sci}) et rédigez la conclusion du domaine.
\end{boxexo}

\begin{boxinfo}
\textbf{Transition.} Les contrôles généraux évalués, on peut enfin se
fier, ou non, aux contrôles intégrés aux applications. Le
Chapitre~\ref{chap:applications} les teste sur la chaîne de la dépense et
de la recette.
\end{boxinfo}
}{}
        \IfFileExists{partie4/P04_C11_controles_application.tex}{% ============================================================================
% Fichier  : partie4/P04_C11_controles_application.tex
% Titre    : Contrôles d'Application de la Dépense et de la Recette
% Auteur   : MINKA MI NGUIDJOÏ Thierry Emmanuel
% Version  : 0.1 -- novembre 2026
% ============================================================================
% Sources : INF3226 C18 (applications et données, resserré) ; GUID 5100
% (gestion des contrôles d'applications) ; COSO principe 10.
% Faits établis ici (LAB-09) :
%  MINÉQUIR, 9 611 mandats 2025 : liquidation > engagement : 0 ;
%   212 mandats > 50 000 000 FCFA, dont 4 sans visa du contrôle financier ;
%   doublons même fournisseur + même montant : 3 paires (P00) ; avec le
%   numéro de facture normalisé en plus : 1 paire (BATIRURAL) ; avec le
%   numéro brut : 0 (c'est pourquoi le contrôle de SIDÉPENSE n'a rien vu).
%  e-RECETTES 2025 : quittances Q-000001 à Q-018640 ; 37 numéros absents ;
%   112 annulations, dont 61 par le même agent (FOMEKONG).
% ============================================================================

\chapter{Contrôles d'Application de la Dépense et de la Recette}
\label{chap:applications}
\chaptermark{Contrôles d'application}

\epigraph{%
    « L'application fait ce qu'on lui a dit~; l'auditeur vérifie ce qu'on
    lui a dit. »%
}{Principe du cours}

\niveauChapitre{Intermédiaire}{Chapitres~\ref{chap:eval-sci}
et~\ref{chap:acces}.}

\begin{boxobjectifs}
\begin{itemize}
    \item Classer les contrôles d'application~: saisie, traitement,
          interface, sortie.
    \item Identifier les contrôles d'application attendus sur la chaîne de
          la dépense et sur celle de la recette.
    \item Tester un contrôle automatique, en sachant à quelle condition un
          essai unique suffit.
    \item Tester un contrôle par l'analyse de toute la population quand
          l'essai n'est pas possible.
    \item Coter les contrôles d'application et en tirer une conclusion.
\end{itemize}
\end{boxobjectifs}

\bigskip

% ============================================================================
\section*{Pourquoi ce chapitre}
% ============================================================================

Une fois les contrôles généraux évalués, l'auditeur sait s'il peut se
fier aux contrôles intégrés à l'application. Ce chapitre les passe en
revue sur la dépense (SIDÉPENSE) et sur la recette (e-RECETTES), et
montre comment les tester sans jamais toucher à la production.

\separateur

% ============================================================================
\section{Quatre familles de contrôles d'application}
\label{sec:c11-familles}
% ============================================================================

\begin{tableaucours}{L{2.4cm}L{4.6cm}L{5.2cm}}{%
    \tableauHeader{Famille} &
    \tableauHeader{Objectif} &
    \tableauHeader{Exemples}
}
Saisie & Seules des données valides entrent & Champs obligatoires,
formats, listes de valeurs, doublons, plafonds \\
Traitement & Les calculs et les enchaînements sont justes & Montant
liquidé au plus égal à l'engagement, calcul des pénalités \\
Interface & Les données transmises arrivent toutes, et intactes &
Totaux de contrôle entre SIDÉPENSE et le poste comptable \\
Sortie & Les états sont complets et ne vont qu'aux destinataires prévus
& Numérotation continue des quittances, états signés \\
\end{tableaucours}
\vspace{-0.8em}
\begin{center}{\small\itshape Tableau~11.1. Les quatre familles de contrôles
d'application}\end{center}

\separateur

% ============================================================================
\section{Sur la chaîne de la dépense}
\label{sec:c11-depense}
% ============================================================================

\begin{tableaucours}{L{4.0cm}L{2.0cm}L{6.2cm}}{%
    \tableauHeader{Contrôle attendu} &
    \tableauHeader{Famille} &
    \tableauHeader{Ce qu'il empêche}
}
Engagement au plus égal au crédit disponible & Saisie & Le dépassement
de crédit \\
Liquidation au plus égale à l'engagement & Traitement & Le paiement d'un
montant non engagé \\
Blocage des doublons (fournisseur, montant, facture normalisée) & Saisie
& Le double paiement \\
Visa du contrôle financier au-delà d'un seuil & Traitement & Le paiement
d'un mandat important sans visa \\
Imputation cohérente avec la nature de la dépense & Saisie & Les
imputations de complaisance \\
Totaux de contrôle de la transmission au poste comptable & Interface &
La perte ou l'ajout d'un mandat en cours de route \\
\end{tableaucours}
\vspace{-0.8em}
\begin{center}{\small\itshape Tableau~11.2. Les contrôles d'application de
la dépense}\end{center}

\begin{boxdef}
\termecle{Normalisation}\index{normalisation}~: transformation d'une
donnée avant comparaison (suppression des espaces, des tirets, des
suffixes comme «~bis~», passage en majuscules), pour que deux valeurs
qui désignent la même chose soient reconnues comme égales. Un contrôle de
doublon sans normalisation se contourne par une lettre.
\end{boxdef}

\separateur

% ============================================================================
\section{Sur la chaîne de la recette}
\label{sec:c11-recette}
% ============================================================================

\begin{tableaucours}{L{4.4cm}L{2.0cm}L{5.8cm}}{%
    \tableauHeader{Contrôle attendu} &
    \tableauHeader{Famille} &
    \tableauHeader{Ce qu'il empêche}
}
Numérotation continue et automatique des quittances & Sortie & Une
recette encaissée sans quittance enregistrée \\
Annulation d'une quittance sur motif et validation d'un second agent &
Traitement & L'annulation de complaisance \\
Rapprochement quotidien des encaissements et des versements & Interface &
Une recette non reversée \\
Calcul automatique des pénalités selon le taux en vigueur & Traitement &
Les remises non autorisées \\
\end{tableaucours}
\vspace{-0.8em}
\begin{center}{\small\itshape Tableau~11.3. Les contrôles d'application de
la recette}\end{center}

\begin{boxscenario}{Dossier e-RECETTES. Trous et annulations}
\textbf{Constatation.} En 2025, les quittances vont de Q-000001 à
Q-018640~; 37 numéros n'existent pas dans la base. Sur 112 annulations,
61 ont été faites par le même agent, M.~FOMEKONG, sans validation d'un
second agent~: e-RECETTES ne l'exige pas.

\medskip
\textbf{Analyse} (confiance moyenne). Les 37 numéros absents peuvent
correspondre à des quittances émises puis supprimées, ce que la
numérotation continue devait empêcher. La concentration des annulations
sur un agent n'est pas une preuve de fraude, mais un indice qui justifie
l'examen pièce par pièce des 61 annulations
(Chapitre~\ref{chap:donnees-fraude}).
\end{boxscenario}

\separateur

% ============================================================================
\section{Tester un contrôle automatique}
\label{sec:c11-tester}
% ============================================================================

\subsection{L'essai, si les contrôles généraux tiennent}

Un contrôle automatique fait la même chose à chaque opération. Si les
contrôles généraux sont efficaces (le programme et le paramètre n'ont pas
changé pendant la période), un essai qui doit passer et un essai qui doit
être bloqué suffisent, en environnement de test, avec un paramétrage
identique à la production.

\subsection{L'analyse de la population, sinon}

Si les contrôles généraux ne tiennent pas, ou s'il n'existe pas
d'environnement de test, on teste le contrôle par son \emph{résultat}~:
on cherche, dans toute la population, les opérations que le contrôle
aurait dû bloquer. Zéro opération trouvée sur toute l'année est une preuve
forte de fonctionnement~; une seule suffit à coter \NON{}.

\begin{boxoutils}{Tester deux contrôles de SIDÉPENSE par la population}
Feuille des 9\,611 mandats~: D montant liquidé, E montant engagé, F
montant, G visa (\texttt{OUI} ou vide).
\begin{lstlisting}
Liquidation > engagement :
H2 : =SI(D2>E2;"ECART";"")
Mandats > 50 M sans visa :
I2 : =SI(ET(F2>50000000;G2<>"OUI");"SANS VISA";"")
\end{lstlisting}
\end{boxoutils}

\begin{boxscenario}{Dossier MINÉQUIR. Deux contrôles par leur résultat}
\textbf{Constatation.} Aucun des 9\,611 mandats n'a une liquidation
supérieure à l'engagement. Sur les 212 mandats de plus de
50\,000\,000~FCFA, 4 n'ont pas de visa du contrôle financier.

\medskip
\textbf{Analyse} (confiance moyenne). Le contrôle de la liquidation a
fonctionné toute l'année (\OUI{}), mais la confiance reste moyenne~: les
accès n'étant pas maîtrisés, un compte ADM aurait pu le désactiver puis
le réactiver. Le contrôle du visa n'est pas bloquant (\NON{})~: les 4
mandats sont à examiner un par un.
\end{boxscenario}

\begin{boxobjet}{C-APP-03 · Blocage des doublons}
\textbf{Objectif.} Une même facture n'est payée qu'une fois.

\textbf{Critère.} SIDÉPENSE refuse un mandat de même fournisseur, même
montant et même facture normalisée qu'un mandat existant.

\textbf{Test de conception.} Le contrôle porte sur le numéro de facture
brut, sans normalisation.

\textbf{Test de fonctionnement.} Analyse des 9\,611 mandats~: avec le
numéro brut, aucune paire~; avec le numéro normalisé, une paire, celle du
prologue (les deux autres paires de même fournisseur et même montant ont
des factures différentes, justifiées par des livraisons distinctes).

\textbf{Cote.} \NON{}.
\end{boxobjet}

\begin{boxlecture}{Contrôles d'application de SIDÉPENSE}
\textbf{Constatation.} Sur six contrôles attendus (Tableau~11.2), trois
fonctionnent~: crédit, liquidation, transmission~; trois non~: doublons,
visa, imputation (non bloquante).

\textbf{Analyse} (confiance moyenne). Les contrôles qui fonctionnent
restent fragiles tant que la déficience majeure des accès n'est pas
corrigée.

\medskip
\hfill\conformite{3}{6}{0}
\end{boxlecture}

\begin{boxmetier}{Vérificateur}
Avant tout essai, demandez s'il existe un environnement de test dont le
paramétrage est identique à la production. S'il n'en existe pas, testez
par la population~: c'est plus long, mais c'est sans risque pour
l'entité.
\end{boxmetier}

\begin{boxmetier}{Contrôleur financier}
Si votre visa n'est pas bloquant dans l'application, chaque mandat qui
passe sans lui est un mandat que vous n'avez pas contrôlé. Demandez que le
seuil soit paramétré comme contrôle bloquant.
\end{boxmetier}

\begin{boxmetier}{Juge des comptes}
Une numérotation de quittances avec des trous est un indice
d'exhaustivité défaillante des recettes. Demandez l'explication de chaque
numéro absent.
\end{boxmetier}

\separateur

\begin{boxresume}
\begin{itemize}
    \item \textbf{Quatre familles.} Saisie, traitement, interface, sortie.
    \item \textbf{Dépense~:} crédit, liquidation, doublons, visa,
          imputation, transmission.
    \item \textbf{Recette~:} numérotation continue, annulations validées,
          rapprochement, pénalités.
    \item \textbf{Normaliser avant de comparer}~: un doublon se cache
          derrière une lettre.
    \item \textbf{Un essai suffit} si les contrôles généraux tiennent~;
          sinon, tester par la population.
    \item \textbf{Zéro sur toute la population} est une preuve forte~; un
          seul cas suffit à coter \NON{}.
\end{itemize}
\end{boxresume}

\begin{boxexo}[title={\ding{45}\quad Exercice~11.1. Classer \niveaubase}]
Rattachez à une famille~: (a) champ RIB obligatoire~; (b) total des
mandats transmis égal au total reçu~; (c) calcul de la pénalité de
retard~; (d) quittance imprimée avec un numéro.
\end{boxexo}

\begin{boxexo}[title={\ding{45}\quad Exercice~11.2. Normaliser \niveaubase}]
Écrivez la formule de tableur qui transforme «~F-2025-117 bis~» et
«~F2025117~» en une même valeur.
\end{boxexo}

\begin{boxexo}[title={\ding{45}\quad Exercice~11.3. Les 37 numéros \niveaumoyen}]
Décrivez votre démarche pour les 37 numéros de quittance absents
d'e-RECETTES~: pièces demandées, personnes interrogées, et ce qui ferait
passer l'indice au constat.
\end{boxexo}

\begin{boxexo}[title={\ding{45}\quad Exercice~11.4. Essai ou population \niveauexpert}]
La DAF vous propose de tester le nouveau blocage des doublons par un essai
en production, «~avec un mandat fictif qu'on annulera~». Dites pourquoi
vous refusez, ce que vous proposez, et à quelles conditions vous
conclurez \OUI{} pour l'année 2026.
\end{boxexo}

\begin{boxinfo}
\textbf{Transition.} L'analyse de toute la population a montré sa
puissance. Le Chapitre~\ref{chap:donnees-fraude} en fait un outil à part
entière, pour la paie et les marchés, et pour la détection de la fraude.
\end{boxinfo}
}{}
        \IfFileExists{partie4/P04_C12_donnees_fraude.tex}{% ============================================================================
% Fichier  : partie4/P04_C12_donnees_fraude.tex
% Titre    : Analyse de Données et Détection de la Fraude
% Auteur   : MINKA MI NGUIDJOÏ Thierry Emmanuel
% Version  : 0.1 -- novembre 2026
% ============================================================================
% Sources : INF3226 C06 (fraude) et C27 (audit continu) ; COSO 2013
% principe 8 ; guide COSO-ACFE (dossier Normes) ; Cressey (1953) ; Benford
% (1938) ; Nigrini (2012).
% Faits CAPNO (LAB-06) : 4 312 agents payés en décembre 2025 ; 37 absents
% du fichier du personnel, net mensuel cumulé 9 842 500 FCFA ; 12 comptes
% bancaires partagés par 27 agents ; 9 agents de plus de 60 ans au
% 31/12/2025 encore payés.
% Faits MINÉQUIR (LAB-07) : 23 bons de commande entre 4 800 000 et
% 4 999 999 FCFA en 2025, adressés à 3 fournisseurs, sous le seuil de
% 5 000 000 FCFA (seuil fictif de mise en concurrence de la DAF) ; chiffre 4
% en premier chiffre : 15,8 % des montants des bons (attendu 9,7 %).
% ============================================================================

\chapter{Analyse de Données et Détection de la Fraude}
\label{chap:donnees-fraude}
\chaptermark{Données et fraude}

\epigraph{%
    « Les données ne mentent pas~; elles attendent qu'on leur pose la
    bonne question. »%
}{Adage des auditeurs de données}

\niveauChapitre{Expert}{Chapitres~\ref{chap:preuve}
et~\ref{chap:applications}.}

\begin{boxobjectifs}
\begin{itemize}
    \item Expliquer le risque de fraude au sens du COSO (principe~8) et le
          triangle de la fraude.
    \item Conduire sur un fichier de paie les tests d'existence, d'unicité
          et de cohérence.
    \item Repérer le fractionnement des commandes sous un seuil.
    \item Appliquer la loi de Benford à une série de montants et en
          interpréter l'écart avec prudence.
    \item Passer d'un indice à un constat, sans confondre anomalie et
          fraude.
\end{itemize}
\end{boxobjectifs}

\bigskip

% ============================================================================
\section*{Pourquoi ce chapitre}
% ============================================================================

Les chapitres précédents ont testé des contrôles. Celui-ci teste les
\emph{données} elles-mêmes, sur toute la population, pour y chercher ce
que les contrôles auraient dû empêcher. C'est l'outil le plus puissant de
l'auditeur public d'aujourd'hui~: un tableur suffit, et quelques
formules révèlent ce que des semaines de contrôle sur pièces ne
voient pas.

\separateur

% ============================================================================
\section{Le risque de fraude}
\label{sec:c12-risque}
% ============================================================================

\begin{boxdef}
\termecle{Fraude}\index{fraude}~: acte intentionnel destiné à obtenir un
avantage indu, par la tromperie ou la dissimulation. Elle se distingue de
l'erreur par l'intention, que l'auditeur ne peut pas établir seul.
\end{boxdef}

Le COSO demande à l'organisation d'évaluer le risque de fraude
(principe~8)~\cite{coso2013icfr}, et le guide commun du COSO et de l'ACFE
en détaille la démarche~\cite{coso2016fraude}. L'analyse classique de
Cressey réunit trois conditions, dites \termecle{triangle de la
fraude}\index{triangle de la fraude}~: une \emph{pression}, une
\emph{occasion}, une \emph{rationalisation}~\cite{cressey1953}.
L'auditeur ne peut agir que sur l'occasion~: c'est exactement ce que
créent les déficiences du contrôle interne.

\begin{tableaucours}{L{3.6cm}L{8.6cm}}{%
    \tableauHeader{Schéma} &
    \tableauHeader{Occasion dans le SI, et test de données}
}
Agent fictif & Un profil peut créer un agent et valider la paie~; tester
l'existence au fichier du personnel \\
Double paiement & Pas de blocage des doublons~; tester fournisseur,
montant, facture normalisée \\
Changement de bénéficiaire & RIB modifiable sans seconde validation~;
tester les RIB modifiés avant un paiement \\
Fractionnement & Seuil de mise en concurrence~; tester les montants juste
sous le seuil \\
Annulation de recette & Annulation sans second agent~; tester la
concentration des annulations \\
\end{tableaucours}
\vspace{-0.8em}
\begin{center}{\small\itshape Tableau~12.1. Schémas de fraude, occasions
et tests}\end{center}

\separateur

% ============================================================================
\section{Analyser un fichier de paie}
\label{sec:c12-paie}
% ============================================================================

\subsection{Trois questions}

\begin{enumerate}
    \item \textbf{Existence}~: chaque agent payé existe-t-il au fichier du
          personnel, en position d'activité~?
    \item \textbf{Unicité}~: chaque agent n'est-il payé qu'une fois, sur un
          compte qui n'est qu'à lui~?
    \item \textbf{Cohérence}~: l'âge, le grade, l'indice et les éléments de
          solde sont-ils compatibles entre eux et avec les règles~?
\end{enumerate}

\begin{boxoutils}{Tests de paie dans un tableur}
Feuille \texttt{PAIE}~: A matricule, B nom, C date de naissance, D compte
bancaire, E net. Feuille \texttt{PERSONNEL}~: A matricule.
\begin{lstlisting}
F2 (existence) : =SI(NB.SI(PERSONNEL!A:A;A2)=0;"ABSENT";"")
G2 (compte)    : =NB.SI(D:D;D2)
H2 (age)       : =DATEDIF(C2;DATE(2025;12;31);"y")
\end{lstlisting}
Un filtre sur \texttt{ABSENT}, sur \texttt{G~>~1} et sur \texttt{H~>~60}
donne les trois listes à examiner.
\end{boxoutils}

\begin{boxscenario}{Dossier CAPNO. Les trois listes}
\textbf{Constatation.} Sur les 4\,312 agents payés en décembre~2025~:
37 ne figurent pas au fichier du personnel (net cumulé
9\,842\,500~FCFA par mois)~; 12 comptes bancaires sont partagés par 27
agents~; 9 agents ont plus de 60 ans et sont toujours payés.

\medskip
\textbf{Analyse} (confiance moyenne). Les 37 absents sont des agents
fictifs \emph{probables}, pas certains~: un retard de mise à jour du
fichier du personnel peut en expliquer une partie. Les comptes partagés
peuvent être des comptes de couples. Chaque cas est examiné avec sa
pièce~; seuls les cas inexpliqués deviennent des constats.
\end{boxscenario}

\begin{boxalert}
\textbf{Violation typique.} Écrire au rapport «~37 agents fictifs~» à la
sortie du tableur.

\smallskip
\textbf{Règle~:} l'analyse de données produit des \emph{indices}. Le
constat vient après l'examen de chaque cas~: dossier individuel,
présence physique, acte de recrutement.
\end{boxalert}

\separateur

% ============================================================================
\section{Le fractionnement des commandes}
\label{sec:c12-fractionnement}
% ============================================================================

Quand un seuil déclenche une procédure plus lourde (mise en concurrence,
visa), la tentation est de découper une commande en plusieurs commandes
sous le seuil. Le test est simple~: compter les montants dans une bande
juste sous le seuil, par fournisseur et par période.

\begin{boxscenario}{Dossier MINÉQUIR. Juste sous le seuil}
\textbf{Constatation.} La DAF fixe à 5\,000\,000~FCFA le seuil de mise en
concurrence de ses bons de commande\footnote{Seuil interne fictif, propre
au cas~; il ne reproduit pas les seuils de la réglementation des marchés
publics.}. En 2025, 23 bons de commande sont compris entre 4\,800\,000 et
4\,999\,999~FCFA~; ils sont adressés à trois fournisseurs seulement.

\medskip
\textbf{Analyse} (confiance moyenne). La concentration dans une bande
étroite sous le seuil, sur peu de fournisseurs, est un indice classique
de fractionnement. Il faut examiner, pour chaque fournisseur, l'objet et
la date des bons~: des commandes de même objet à quelques jours
d'intervalle confirmeraient l'indice.
\end{boxscenario}

\separateur

% ============================================================================
\section{La loi de Benford}
\label{sec:c12-benford}
% ============================================================================

Dans de nombreuses séries de montants naturels, le premier chiffre n'est
pas uniformément réparti~: le 1 apparaît environ trois fois sur dix, le 9
moins d'une fois sur vingt~\cite{benford1938,nigrini2012benford}. La
proportion attendue du premier chiffre $d$ est $\log_{10}(1 + 1/d)$.

\begin{tableaucours}{L{3.0cm}ccccccccc}{%
    \tableauHeader{Premier chiffre} &
    \tableauHeader{1} & \tableauHeader{2} & \tableauHeader{3} &
    \tableauHeader{4} & \tableauHeader{5} & \tableauHeader{6} &
    \tableauHeader{7} & \tableauHeader{8} & \tableauHeader{9}
}
Attendu (\%) & 30,1 & 17,6 & 12,5 & 9,7 & 7,9 & 6,7 & 5,8 & 5,1 & 4,6 \\
\end{tableaucours}
\vspace{-0.8em}
\begin{center}{\small\itshape Tableau~12.2. Distribution de Benford du
premier chiffre}\end{center}

\begin{boxoutils}{Premier chiffre et proportion attendue}
\begin{lstlisting}
B2 (premier chiffre) : =GAUCHE(TEXTE(A2;"0");1)*1
Attendu pour d       : =LOG10(1+1/d)
Observe pour d       : =NB.SI(B:B;d)/NB(B:B)
\end{lstlisting}
\end{boxoutils}

\begin{boxinfo}
\textbf{Prudence.} La loi de Benford vaut pour des séries nombreuses,
étalées sur plusieurs ordres de grandeur, et non bornées par des seuils ou
des barèmes. Une paie (barème) ou des montants plafonnés ne la suivent
pas, sans aucune fraude. L'écart désigne où regarder~; il ne prouve rien.
\end{boxinfo}

\begin{boxscenario}{Dossier MINÉQUIR. Le chiffre 4}
\textbf{Constatation.} Parmi les montants des bons de commande de 2025, le
premier chiffre 4 apparaît dans 15,8~\% des cas, contre 9,7~\% attendus.

\medskip
\textbf{Analyse} (confiance faible à moyenne). L'excès de 4 est cohérent
avec les 23 bons compris entre 4,8 et 5~millions. Benford ne dit rien de
plus que le test de bande~; il le confirme par une autre voie.
\end{boxscenario}

\separateur

% ============================================================================
\section{De l'indice au constat}
\label{sec:c12-constat}
% ============================================================================

\begin{boxprocedure}
\textbf{Feuille d'analyse de données}~:
\begin{enumerate}
    \item fichier analysé (source, date, empreinte, totaux rapprochés)~;
    \item question posée et formule ou requête utilisée, recopiée
          intégralement~;
    \item résultat brut (nombre de cas, liste en annexe)~;
    \item examen de chaque cas, avec la pièce~: expliqué, inexpliqué, en
          attente~;
    \item constat sur les seuls cas inexpliqués~; \NEK{} sur les cas en
          attente.
\end{enumerate}
\end{boxprocedure}

\begin{boxmetier}{Vérificateur}
Gardez la formule exacte de chaque test dans le dossier. Lors du
contradictoire, l'entité refera votre calcul~; elle doit trouver le même
résultat.
\end{boxmetier}

\begin{boxmetier}{Juge des comptes}
Un paiement à un agent absent du fichier du personnel engage la
responsabilité de celui qui a payé sans vérifier l'existence du
créancier. Demandez quel contrôle d'application devait l'en empêcher.
\end{boxmetier}

\begin{boxmetier}{Auditeur interne}
Faites tourner chaque mois les trois tests de paie et le test de
fractionnement~: c'est un contrôle continu peu coûteux, qui relève du
pilotage (principe~16 du COSO).
\end{boxmetier}

\separateur

\begin{boxresume}
\begin{itemize}
    \item \textbf{La fraude suppose une intention}, que l'auditeur ne peut
          établir seul~; il agit sur l'occasion.
    \item \textbf{Chaque déficience ouvre un schéma}, et chaque schéma
          appelle un test de données.
    \item \textbf{Paie~:} existence, unicité, cohérence.
    \item \textbf{Fractionnement~:} une bande étroite sous le seuil, peu de
          fournisseurs.
    \item \textbf{Benford désigne, ne prouve pas}~; il ne s'applique pas à
          des séries bornées.
    \item \textbf{Un indice n'est pas un constat.} Chaque cas est examiné
          avec sa pièce.
\end{itemize}
\end{boxresume}

\begin{boxexo}[title={\ding{45}\quad Exercice~12.1. Le triangle \niveaubase}]
Pour le double paiement du MINÉQUIR, dites ce que l'auditeur peut
établir de chaque côté du triangle de la fraude, et ce qui ne relève pas
de lui.
\end{boxexo}

\begin{boxexo}[title={\ding{45}\quad Exercice~12.2. Benford \niveaubase}]
Calculez la proportion attendue du premier chiffre 1 et du premier
chiffre 9. Dites pourquoi on n'applique pas le test aux indemnités de
logement de la CAPNO, qui suivent un barème.
\end{boxexo}

\begin{boxexo}[title={\ding{45}\quad Exercice~12.3. Les comptes partagés \niveaumoyen}]
Proposez la démarche d'examen des 12 comptes partagés par 27 agents de la
CAPNO~: pièces, ordre d'examen, critères pour classer un cas comme
expliqué.
\end{boxexo}

\begin{boxexo}[title={\ding{45}\quad Exercice~12.4. Du tableur au rapport \niveauexpert}]
Après examen, 31 des 37 agents absents restent inexpliqués. Rédigez le
constat en cinq éléments, avec l'effet chiffré sur l'exercice 2025, et
dites ce que vous écrivez sur les 6 autres.
\end{boxexo}

\begin{boxinfo}
\textbf{Transition.} Il reste à regarder le système d'information de
plus haut~: qui le pilote, comment il est acheté, et à qui il est confié.
C'est l'objet du Chapitre~\ref{chap:gouvernance}.
\end{boxinfo}
}{}
        \IfFileExists{partie4/P04_C13_gouvernance_projets_prestataires.tex}{% ============================================================================
% Fichier  : partie4/P04_C13_gouvernance_projets_prestataires.tex
% Titre    : Gouvernance, Projets Numériques et Prestataires
% Auteur   : MINKA MI NGUIDJOÏ Thierry Emmanuel
% Version  : 0.1 -- novembre 2026
% ============================================================================
% Sources : INF3226 C17, C19, C24 (resserrés) ; COBIT 2019 (gouvernance,
% projets) ; GUID 5100 ; loi n° 2024/017 (décrite sans citation d'article).
% Faits e-RECETTES établis ici : marché de 640 000 000 FCFA attribué à
% NUMÉRIKA SARL en mars 2024 ; mise en service le 2 janvier 2025 ; PV de
% réception du 15 janvier 2025 signé avec 47 anomalies ouvertes, dont 9
% bloquantes ; aucun schéma directeur ; pas de clause de réversibilité ni
% de propriété des données ; hébergement hors du pays ; compte du
% prestataire avec droits d'administration en production ; 12 400
% contribuables enregistrés (données personnelles).
% ============================================================================

\chapter{Gouvernance, Projets Numériques et Prestataires}
\label{chap:gouvernance}
\chaptermark{Gouvernance, projets, prestataires}

\epigraph{%
    « On peut confier un service~; on ne confie pas sa responsabilité. »%
}{Principe du cours}

\niveauChapitre{Intermédiaire}{Chapitre~\ref{chap:sci}.}

\begin{boxobjectifs}
\begin{itemize}
    \item Évaluer la gouvernance de la fonction informatique d'une
          administration~: schéma directeur, organisation, pilotage.
    \item Auditer un projet numérique et le marché qui le porte, de
          l'expression du besoin à la réception.
    \item Évaluer les contrats avec les prestataires~: propriété des
          données, réversibilité, accès, niveaux de service.
    \item Repérer les risques de l'hébergement externe.
    \item Situer les obligations de protection des données personnelles.
\end{itemize}
\end{boxobjectifs}

\bigskip

% ============================================================================
\section*{Pourquoi ce chapitre}
% ============================================================================

Les déficiences des chapitres précédents ont souvent une cause commune,
plus haut~: un système acheté sans besoin écrit, reçu sans tests, confié
à un prestataire sans contrat solide, piloté par personne. Ce chapitre
remonte à cette cause, sur le dossier e-RECETTES.

\separateur

% ============================================================================
\section{La gouvernance de la fonction informatique}
\label{sec:c13-gouvernance}
% ============================================================================

\begin{tableaucours}{L{3.2cm}L{4.4cm}L{4.6cm}}{%
    \tableauHeader{Élément} &
    \tableauHeader{Question} &
    \tableauHeader{Pièce}
}
Schéma directeur & Les projets répondent-ils à des objectifs écrits~? &
Schéma directeur adopté, suivi annuel \\
Organisation & Qui décide, qui exécute, qui contrôle~? & Organigramme,
décisions, propriétaires d'application \\
Pilotage & Un comité suit-il les projets et les risques~? & Comptes rendus
\\
Budget & Le coût complet du SI est-il connu~? & Budget informatique,
contrats \\
Compétences & L'entité peut-elle contrôler ses prestataires~? & Effectifs,
formations \\
\end{tableaucours}
\vspace{-0.8em}
\begin{center}{\small\itshape Tableau~13.1. Les éléments de la gouvernance
de la fonction informatique}\end{center}

COBIT~2019 sépare la \emph{gouvernance} (évaluer, diriger, surveiller),
qui relève de la direction, de la \emph{gestion}, qui relève des
services~\cite{isaca2018cobitintro}. Dans une administration, la
gouvernance revient à l'ordonnateur et à la direction~; elle ne peut pas
être laissée au prestataire.

\separateur

% ============================================================================
\section{Auditer un projet numérique}
\label{sec:c13-projet}
% ============================================================================

\subsection{Le cycle du projet et ses contrôles}

\begin{tableaucours}{L{2.6cm}L{5.0cm}L{4.6cm}}{%
    \tableauHeader{Phase} &
    \tableauHeader{Contrôle attendu} &
    \tableauHeader{Pièce}
}
Besoin & Besoin écrit par le métier, rattaché au schéma directeur & Note
de besoin, étude d'opportunité \\
Marché & Cahier des charges qui inclut les contrôles internes attendus,
la sécurité, la réversibilité & Dossier d'appel d'offres, contrat \\
Réalisation & Suivi par un comité, avec les risques & Comptes rendus \\
Recette & Tests par le métier, anomalies bloquantes levées avant
réception & Cahier de recette, PV de réception \\
Mise en service & Reprise des données contrôlée, formation & PV de reprise,
totaux de contrôle \\
\end{tableaucours}
\vspace{-0.8em}
\begin{center}{\small\itshape Tableau~13.2. Les contrôles du cycle d'un
projet}\end{center}

\begin{boxscenario}{Dossier e-RECETTES. Une réception trop tôt}
\textbf{Constatation.} Le marché de 640\,000\,000~FCFA a été attribué à
NUMÉRIKA~SARL en mars~2024. e-RECETTES est en service depuis le
2~janvier~2025. Le procès-verbal de réception, signé le 15~janvier~2025,
mentionne 47 anomalies ouvertes, dont 9 bloquantes. Le cahier des charges
ne prévoit ni contrôle des annulations, ni numérotation continue des
quittances. Il n'existe pas de schéma directeur.

\medskip
\textbf{Analyse} (confiance élevée). Les déficiences constatées aux
Chapitres~\ref{chap:exploitation} et~\ref{chap:applications} étaient
inscrites dans le projet~: ce qui n'est pas demandé au cahier des
charges n'est pas livré, et ce qui est reçu avec des anomalies
bloquantes est payé sans être conforme.
\end{boxscenario}

\begin{boxalert}
\textbf{Violation typique.} Prononcer la réception d'un système pour
respecter le calendrier budgétaire, avec des anomalies bloquantes
ouvertes.

\smallskip
\textbf{Règle~:} une réception avec réserves exige un plan de levée daté
et une retenue sur paiement~; une réception sans réserve d'un système
défaillant engage la responsabilité de celui qui la signe.
\end{boxalert}

\separateur

% ============================================================================
\section{Les prestataires et l'hébergement externe}
\label{sec:c13-prestataires}
% ============================================================================

\begin{tableaucours}{L{3.4cm}L{8.8cm}}{%
    \tableauHeader{Clause} &
    \tableauHeader{Ce que l'auditeur vérifie}
}
Propriété des données & Les données appartiennent à l'administration, qui
peut les obtenir à tout moment dans un format exploitable \\
Réversibilité & Les conditions de reprise du service par l'administration
ou un autre prestataire sont écrites \\
Accès du prestataire & Ses accès en production sont nominatifs, limités,
journalisés, et ouverts sur demande \\
Niveaux de service & Disponibilité, délais de correction, pénalités \\
Localisation & Lieu d'hébergement connu et compatible avec les obligations
légales \\
Droit d'audit & L'administration et ses contrôleurs peuvent auditer le
prestataire \\
\end{tableaucours}
\vspace{-0.8em}
\begin{center}{\small\itshape Tableau~13.3. Les clauses à vérifier dans un
contrat de prestation informatique}\end{center}

\begin{boxscenario}{Dossier e-RECETTES. Le contrat de NUMÉRIKA}
\textbf{Constatation.} Le contrat ne contient ni clause de propriété des
données, ni clause de réversibilité, ni droit d'audit. Les données sont
hébergées hors du pays. Le prestataire dispose d'un compte
d'administration permanent en production (Chapitre~\ref{chap:exploitation}).

\medskip
\textbf{Analyse} (confiance élevée). La recette des finances dépend
entièrement de NUMÉRIKA pour accéder à ses propres données. Le décret
n°~2013/287 oblige les responsables de l'entité auditée à présenter les
pièces~; s'ils ne peuvent pas les obtenir de leur prestataire, c'est
l'entité qui est en défaut.
\end{boxscenario}

\separateur

% ============================================================================
\section{Les données personnelles}
\label{sec:c13-donnees}
% ============================================================================

Les systèmes financiers publics traitent des données personnelles~: agents
(paie), contribuables, usagers, fournisseurs personnes physiques. La loi
n°~2024/017 du 23~décembre~2024 encadre ces traitements au Cameroun. Le
cours n'en cite pas les articles~; l'auditeur les lit avant de s'y
référer. Il vérifie au moins que les traitements sont recensés, que
l'accès aux données est limité aux personnes qui en ont besoin, et que
les transferts vers un hébergeur, en particulier hors du pays, sont
encadrés.

\begin{boxinfo}
\textbf{Les extractions de l'auditeur sont aussi des données
personnelles.} Le fichier de paie de la CAPNO, avec ses 4\,312 noms et
comptes bancaires, doit être protégé par le vérificateur avec le même
soin qu'il l'exige de l'entité (Chapitre~\ref{chap:normes}).
\end{boxinfo}

\begin{boxlecture}{Gouvernance et projet e-RECETTES}
\textbf{Constatation.} Sur huit contrôles attendus (schéma directeur,
cahier des charges complet, comité de suivi, réception sans anomalie
bloquante, propriété des données, réversibilité, accès du prestataire
encadré, droit d'audit), un seul fonctionne~: un comité de suivi s'est
réuni trois fois pendant le projet.

\textbf{Analyse} (confiance élevée). La cause première des déficiences
d'e-RECETTES est dans la gouvernance du projet.

\medskip
\hfill\conformite{1}{8}{0}
\end{boxlecture}

\begin{boxmetier}{Vérificateur}
Pour tout système acquis par marché, demandez le cahier des charges et le
procès-verbal de réception. Comparez les contrôles demandés aux contrôles
livrés~: l'écart explique souvent les déficiences.
\end{boxmetier}

\begin{boxmetier}{Inspecteur}
Une réception signée avec des anomalies bloquantes ouvertes est un fait
précis, daté et signé. Elle mérite une question écrite à son signataire.
\end{boxmetier}

\begin{boxmetier}{Contrôleur financier}
Avant de viser le paiement d'un solde de marché informatique, demandez le
procès-verbal de réception et la liste des réserves levées.
\end{boxmetier}

\separateur

\begin{boxresume}
\begin{itemize}
    \item \textbf{La gouvernance reste à l'administration}~: schéma
          directeur, propriétaires, comité, budget, compétences.
    \item \textbf{Le projet porte les contrôles}~: ce qui n'est pas au
          cahier des charges n'est pas livré.
    \item \textbf{La réception engage}~; des anomalies bloquantes ouvertes
          imposent réserves, plan et retenue.
    \item \textbf{Six clauses de contrat}~: données, réversibilité, accès,
          niveaux de service, localisation, droit d'audit.
    \item \textbf{Les données personnelles} des systèmes financiers sont
          encadrées par la loi n°~2024/017, y compris chez l'auditeur.
\end{itemize}
\end{boxresume}

\begin{boxexo}[title={\ding{45}\quad Exercice~13.1. Gouvernance ou gestion~? \niveaubase}]
Classez~: (a) adopter un schéma directeur~; (b) faire les sauvegardes~;
(c) décider d'externaliser l'hébergement~; (d) attribuer un profil.
\end{boxexo}

\begin{boxexo}[title={\ding{45}\quad Exercice~13.2. Un cahier des charges \niveaumoyen}]
Rédigez les six exigences de contrôle interne que le cahier des charges
d'e-RECETTES aurait dû contenir, à partir des
Chapitres~\ref{chap:acces} à~\ref{chap:applications}.
\end{boxexo}

\begin{boxexo}[title={\ding{45}\quad Exercice~13.3. Le contrat \niveaumoyen}]
Rédigez la recommandation sur le contrat de NUMÉRIKA, avec les clauses à
ajouter par avenant et l'ordre de priorité.
\end{boxexo}

\begin{boxexo}[title={\ding{45}\quad Exercice~13.4. Les causes \niveauexpert}]
Reliez chaque déficience d'e-RECETTES relevée aux Chapitres~\ref{chap:exploitation}
et~\ref{chap:applications} à une défaillance du projet ou du contrat, et
dites quelle recommandation de gouvernance les traiterait toutes.
\end{boxexo}

\begin{boxinfo}
\textbf{Transition.} Toutes les pièces sont en place. Le
Chapitre~\ref{chap:cas-integrateur} les réunit dans une mission
complète, conduite en équipe en deux heures.
\end{boxinfo}
}{}

    % ======================================================================
    % PARTIE 5 — SYNTHÈSE (J08)
    % « Une mission se juge à sa conclusion. »
    % Standards de référence : l'ensemble du cours.
    % Niveau : expert
    % ======================================================================
    \part{Synthèse}
    \label{part:synthese}

        \IfFileExists{partie5/P05_C14_cas_integrateur.tex}{% ============================================================================
% Fichier  : partie5/P05_C14_cas_integrateur.tex
% Titre    : Cas Intégrateur : la Mission e-RECETTES
% Auteur   : MINKA MI NGUIDJOÏ Thierry Emmanuel
% Version  : 0.1 -- novembre 2026
% ============================================================================
% J08, 2 h en salle, par équipes de 4 à 5 ; les 2 heures suivantes sont
% offertes. Reprend les faits e-RECETTES de C10, C11, C13. Faits ajoutés :
% profils e-RECETTES : FOMEKONG (ENC, ANN), TAGNE (ENC, ANN, VAL),
% MBIDA (VER), numerika_adm (ADM, prestataire) ; recette totale 2025 :
% 2 184 600 000 FCFA ; versements au Trésor : 2 171 950 000 FCFA (écart
% 12 650 000 FCFA, non expliqué à la date de la mission).
% Corrigé détaillé : livret (corriges/Corr_C14.tex).
% ============================================================================

\chapter{Cas Intégrateur : la Mission e-RECETTES}
\label{chap:cas-integrateur}
\chaptermark{Cas intégrateur}

\epigraph{%
    « Une mission se juge à sa conclusion. »%
}{Règle du cours}

\niveauChapitre{Expert}{L'ensemble du cours.}

\begin{boxobjectifs}
\begin{itemize}
    \item Conduire, en équipe et en temps limité, une mission complète
          d'évaluation du contrôle interne d'un système d'information
          public.
    \item Mobiliser le mandat, la matrice des risques et des contrôles, le
          questionnaire, les tests et le classement des déficiences.
    \item Rédiger au moins deux constats en cinq éléments.
    \item Formuler et défendre oralement une conclusion sur l'efficacité du
          contrôle interne.
\end{itemize}
\end{boxobjectifs}

\bigskip

% ============================================================================
\section*{Mode d'emploi}
% ============================================================================

Le cas se traite en salle, le dernier jour, en deux heures, par équipes de
quatre ou cinq. Chaque équipe est une équipe de vérification du Contrôle
Supérieur de l'État. Elle reçoit le dossier de la section~\ref{sec:c14-dossier},
suit le déroulé de la section~\ref{sec:c14-deroule}, et restitue ses
résultats oralement en dix minutes. Le formateur dispose du corrigé dans
le livret de correction.

\separateur

% ============================================================================
\section{Le mandat}
\label{sec:c14-mandat}
% ============================================================================

\begin{tableaucours}{L{3.0cm}L{9.2cm}}{%
    \tableauHeader{Rubrique} &
    \tableauHeader{Ordre de mission}
}
Fondement & Décret n°~2013/287, art.~2, al.~3 et~4 \\
Objectif & Le contrôle interne d'e-RECETTES garantit-il l'exhaustivité et
la régularité des recettes de la recette des finances de Bafoussam~? \\
Périmètre & e-RECETTES~; le contrat avec NUMÉRIKA~SARL~; les versements
au Trésor \\
Période & Exercice 2025 \\
Livrables & Matrice des risques et des contrôles, cotations, tableau des
déficiences, deux constats, conclusion \\
\end{tableaucours}
\vspace{-0.8em}
\begin{center}{\small\itshape Tableau~14.1. Ordre de mission du cas
intégrateur}\end{center}

\separateur

% ============================================================================
\section{Le dossier}
\label{sec:c14-dossier}
% ============================================================================

\subsection*{Pièce 1. Fiche de l'entité}

La recette des finances de Bafoussam encaisse les recettes de l'État de
son ressort. Depuis le 2~janvier~2025, les encaissements passent par
e-RECETTES, livré par NUMÉRIKA~SARL au titre d'un marché de
640\,000\,000~FCFA. 12\,400 contribuables y sont enregistrés. Trois agents
utilisent l'application.

\subsection*{Pièce 2. Profils d'e-RECETTES}

\begin{tableaucours}{L{3.0cm}L{4.4cm}L{4.8cm}}{%
    \tableauHeader{Compte} &
    \tableauHeader{Fonction} &
    \tableauHeader{Profils}
}
\texttt{fomekong.a} & Caissier & ENC (encaisser), ANN (annuler) \\
\texttt{tagne.m} & Chef de la recette & ENC, ANN, VAL (valider la
journée) \\
\texttt{mbida.s} & Comptable & VER (verser au Trésor) \\
\texttt{numerika\_adm} & Prestataire & ADM (administration) \\
\end{tableaucours}

\subsection*{Pièce 3. Réponses au questionnaire de contrôle interne}

\begin{tableaucours}{L{6.0cm}L{6.2cm}}{%
    \tableauHeader{Question} &
    \tableauHeader{Réponse déclarée}
}
Les quittances sont-elles numérotées automatiquement et sans trou~? &
«~Oui, l'application numérote.~» \\
Une annulation exige-t-elle une validation~? & «~Le chef valide la
journée.~» \\
Les encaissements sont-ils rapprochés des versements~? & «~Chaque fin de
mois.~» \\
Les changements de l'application sont-ils approuvés~? & «~NUMÉRIKA nous
prévient.~» \\
Les sauvegardes sont-elles testées~? & «~NUMÉRIKA s'en occupe.~» \\
\end{tableaucours}

\subsection*{Pièce 4. Résultats d'extractions (faites devant l'équipe)}

\begin{itemize}
    \item Quittances 2025~: Q-000001 à Q-018640~; 37 numéros absents.
    \item Annulations~: 112, dont 61 par \texttt{fomekong.a}~; aucune
          validation d'un second agent n'est enregistrée.
    \item Recettes encaissées en 2025~: 2\,184\,600\,000~FCFA~; versements
          au Trésor~: 2\,171\,950\,000~FCFA.
    \item Aucune trace d'un rapprochement mensuel.
\end{itemize}

\subsection*{Pièce 5. Journaux techniques}

\begin{itemize}
    \item 14 mises en production en 2025, dont 11 sans demande écrite.
    \item 9~septembre~2025~: modification directe du taux de pénalité par
          \texttt{numerika\_adm}.
    \item Sauvegardes~: 23 échecs, dont 6 consécutifs~; aucune
          restauration testée.
\end{itemize}

\subsection*{Pièce 6. Contrat et réception}

\begin{itemize}
    \item Ni propriété des données, ni réversibilité, ni droit d'audit~;
          hébergement hors du pays.
    \item Procès-verbal de réception du 15~janvier~2025~: 47 anomalies
          ouvertes, dont 9 bloquantes.
    \item Aucun audit de sécurité de l'ANTIC.
\end{itemize}

\subsection*{Pièce 7. Réponses des répondants (instrument OUI / NON / NEK)}

Cinq questions par domaine, posées à quatre fonctions (chef de la recette,
caissier, comptable, correspondant de NUMÉRIKA), soit vingt réponses par
domaine.

\begin{tableaucours}{L{5.0cm}ccc}{%
    \tableauHeader{Domaine} &
    \tableauHeader{\OUI{}} & \tableauHeader{\NON{}} & \tableauHeader{\NEK{}}
}
R1 Encaissement, quittances, annulations & 9 & 4 & 7 \\
R2 Accès et changements & 6 & 7 & 7 \\
R3 Sauvegardes et continuité & 3 & 9 & 8 \\
\end{tableaucours}

\separateur

% ============================================================================
\section{Le déroulé}
\label{sec:c14-deroule}
% ============================================================================

\begin{tableaucours}{L{2.0cm}L{5.4cm}L{4.8cm}}{%
    \tableauHeader{Temps} &
    \tableauHeader{Travail de l'équipe} &
    \tableauHeader{Chapitres mobilisés}
}
0 -- 10 min & Lire le mandat~; répartir les rôles (chef de mission,
rédacteur, deux testeurs) & \ref{chap:institutions},
\ref{chap:preparer} \\
10 -- 30 min & Construire la matrice des risques et des contrôles
(au moins 10 lignes) & \ref{chap:eval-sci}, \ref{chap:applications} \\
30 -- 50 min & Confronter le QCI aux pièces~; coter chaque contrôle~;
calculer $\mu$, $U$, $H$ (pièce~7) et situer les domaines dans le quadrant
& \ref{chap:preuve}, \ref{chap:constater} \\
50 -- 65 min & Regrouper et classer les déficiences & \ref{chap:eval-sci} \\
65 -- 85 min & Rédiger deux observations et la conclusion &
\ref{chap:constater} \\
85 -- 105 min & Restitution~: 5 à 7 minutes par équipe, questions du jury
& tout le cours \\
105 -- 120 min & Ouverture sur la recherche~: quel sujet de mémoire ce
cas suggère-t-il~? & \ref{chap:recherche} \\
\end{tableaucours}
\vspace{-0.8em}
\begin{center}{\small\itshape Tableau~14.2. Déroulé du cas en deux
heures}\end{center}

\separateur

% ============================================================================
\section{La grille de restitution}
\label{sec:c14-grille}
% ============================================================================

\begin{tableaucours}{L{7.8cm}c}{%
    \tableauHeader{Critère} &
    \tableauHeader{Points}
}
Le mandat est cité et l'objectif reformulé correctement & 2 \\
La matrice part des risques et couvre accès, changements, exploitation,
application, versements, prestataire & 4 \\
Les cotes distinguent les déclarations (QCI) des preuves~; les trois
domaines sont situés dans le quadrant & 3 \\
Le classement identifie au moins une déficience majeure, et la justifie &
4 \\
Les deux observations ont leurs sept éléments et parlent de comptes, pas
de personnes & 4 \\
La conclusion porte sur un objectif, une période et un périmètre & 3 \\
\textbf{Total} & \textbf{20} \\
\end{tableaucours}
\vspace{-0.8em}
\begin{center}{\small\itshape Tableau~14.3. Grille d'évaluation de la
restitution}\end{center}

\separateur

% ============================================================================
\section{Les erreurs fréquentes}
\label{sec:c14-erreurs}
% ============================================================================

\begin{tableaucours}{L{5.8cm}L{6.4cm}}{%
    \tableauHeader{Erreur} &
    \tableauHeader{Ce qu'il fallait faire}
}
Conclure « conforme à 45~\% » sur R1 & $U = 0{,}35 \geq 0{,}20$~: le
domaine est opaque~; aucune conclusion de conformité (règle~1) \\
Coter \OUI{} la numérotation parce que l'application «~numérote~» &
Partir des 37 numéros absents~: \NON{} \\
Coter \NEK{} le rapprochement mensuel & Aucune trace~: l'absence d'une
trace exigée est un \NON{} \\
Écrire «~M.~FOMEKONG détourne des recettes~» & Constater la
concentration des annulations sur un compte, sans validation \\
Conclure sur la sécurité d'ensemble & Rester dans l'objectif~; signaler
l'absence d'audit de l'ANTIC \\
Additionner les \NON{} comme autant de déficiences & Regrouper par cause
(accès, changements, projet) \\
Oublier l'écart de 12\,650\,000~FCFA & C'est l'effet chiffré le plus
direct sur l'objectif d'exhaustivité \\
\end{tableaucours}
\vspace{-0.8em}
\begin{center}{\small\itshape Tableau~14.4. Grille des erreurs
fréquentes}\end{center}

\begin{boxmetier}{Vérificateur}
Dans un vrai dossier, l'écart entre recettes encaissées et versements
appellerait une note de faits distincte du rapport
(Chapitre~\ref{chap:constater})~: c'est une suite possible devant le CDBF.
\end{boxmetier}

\begin{boxmetier}{Juge des comptes}
L'écart entre encaissements et versements concerne directement le compte
du comptable~: il sera examiné, par extension, à travers les comptes
produits par e-RECETTES.
\end{boxmetier}

\separateur

\begin{boxresume}
\begin{itemize}
    \item \textbf{Une mission complète} tient en six gestes~: mandat,
          matrice, cotation, classement, constats, conclusion.
    \item \textbf{Le QCI se confronte aux pièces}~; une déclaration ne vaut
          pas preuve.
    \item \textbf{Les déficiences se regroupent par cause.}
    \item \textbf{La conclusion se défend}~: objectif, période, périmètre,
          travaux.
\end{itemize}
\end{boxresume}

\begin{boxexo}[title={\ding{45}\quad Exercice~14.1. La matrice \niveaumoyen}]
Construisez la matrice des risques et des contrôles d'e-RECETTES (au moins
dix lignes), à partir des pièces du dossier.
\end{boxexo}

\begin{boxexo}[title={\ding{45}\quad Exercice~14.2. Les déficiences \niveauexpert}]
Regroupez les contrôles cotés \NON{} en déficiences, classez-les avec la
grille du Chapitre~\ref{chap:eval-sci}, et justifiez chaque déficience
majeure.
\end{boxexo}

\begin{boxexo}[title={\ding{45}\quad Exercice~14.3. La conclusion \niveauexpert}]
Rédigez la conclusion de la mission et le premier paragraphe de la
synthèse du rapport.
\end{boxexo}

\begin{boxinfo}
\textbf{Transition.} La mission est complète. Le
Chapitre~\ref{chap:recherche} remonte aux fondements de la méthode, le
programme de recherche de l'auditologie, et ouvre des pistes de mémoire
pour les auditeurs qui voudront contribuer à le tester.
\end{boxinfo}
}{}

    % ======================================================================
    % PARTIE 6 — OUVERTURE SUR LA RECHERCHE (J08, fin de journée, et lecture)
    % « Un programme qui énonce sa falsifiabilité n'est pas une idéologie. »
    % Standards de référence : articles A0 (hal-05701007) et A1
    %   (hal-05701270) ; travaux publiés.
    % Niveau : recherche
    % ======================================================================
    \part{Ouverture sur la Recherche}
    \label{part:recherche}

        \IfFileExists{partie6/P06_C15_recherche_auditologie.tex}{% ============================================================================
% Fichier  : partie6/P06_C15_recherche_auditologie.tex
% Titre    : Ouvrir sur la Recherche : l'Auditologie
% Auteur   : MINKA MI NGUIDJOÏ Thierry Emmanuel
% Version  : 0.1 -- novembre 2026
% ============================================================================
% Ajouté à la demande de l'auteur (07/11/2026) : ultime chapitre qui ouvre le
% cours sur la recherche, avec des axes et des exemples de sujets.
% Sources : A0 « Vers l'auditologie » (hal-05701007) ; A1 « Modelling
% Epistemic Uncertainty in IS Audit » (hal-05701270) ; INF3226 C30 et C31.
% Règle de l'auteur (06/11/2026) : seuls A0 et A1 (HAL) et des travaux
% publiés sont cités ; aucun travail non publié n'est annoncé.
% Axiomes cités avec la numérotation d'A0 (1 à 6), pas avec l'environnement
% auto-numéroté.
% ============================================================================

\chapter[Ouvrir sur la recherche : l'auditologie]{Ouvrir sur la Recherche : l'Auditologie}
\label{chap:recherche}
\chaptermark{Ouvrir sur la recherche}

\epigraph{%
    « Ce que l'organisation ignore de ses propres systèmes conditionne la
    qualité des décisions qu'elle peut prendre à leur sujet. »%
}{\textit{Vers l'auditologie}~\cite{minka2026auditologie}}

\niveauChapitre{Recherche}{L'ensemble du cours, en particulier les
Chapitres~\ref{chap:preparer} à~\ref{chap:constater}.}

\begin{boxobjectifs}
\begin{itemize}
    \item Définir l'auditologie et la mission auditologique.
    \item Énoncer les six axiomes de l'article A0 et relier chacun à une
          pratique du cours.
    \item Présenter les cinq piliers théoriques et les prédictions
          falsifiables qu'ils engendrent.
    \item Distinguer ce qui est établi par des travaux publiés, ce qui est
          déposé en prépublication, et ce qui reste à valider.
    \item Identifier des axes de recherche à partir de la pratique de
          l'audit public, et formuler un sujet de mémoire falsifiable.
\end{itemize}
\end{boxobjectifs}

\bigskip

% ============================================================================
\section*{Pourquoi ce chapitre}
% ============================================================================

Tout au long du cours, des choix de méthode ont été justifiés par une
recherche~: la réponse \NEK{}, la décomposition en huit niveaux, les mesures
$\mu$, $U$ et $H$, le quadrant, la comparaison des cycles. Ces outils
viennent d'un programme de recherche que l'auteur a nommé
\emph{auditologie}~\cite{minka2026auditologie}, et dont l'instrument de
réponse est étudié dans un second article~\cite{minka2026nek}. Ce chapitre
le présente pour lui-même, puis montre comment un auditeur des finances
publiques peut y contribuer, par un mémoire ou par sa pratique.

\begin{boiteRemarque}{Ce que ce chapitre est, et n'est pas}
Il expose un programme de recherche déposé en prépublication, non un savoir
établi. Les axiomes et les piliers forment un cadre~; les prédictions qui en
découlent restent à tester. C'est précisément là que les auditeurs publics,
qui disposent de missions réelles et répétées, peuvent apporter ce qui
manque~: des données longitudinales.
\end{boiteRemarque}

\separateur

% ============================================================================
\section{L'objet : la mission auditologique}
\label{sec:c15-objet}
% ============================================================================

L'article A0 part d'un constat~: l'audit des systèmes d'information dispose
de méthodes, de référentiels et d'une abondante littérature instrumentale,
mais d'aucun socle axiomatique ni d'aucun espace de mesure unifié pour son
objet central, la réduction de l'incertitude. Il propose de le constituer en
une science appliquée.

\begin{boxdef}
\termecle{Mission auditologique}\index{mission auditologique} (A0)~:
quintuplet
$\mathcal{M} = \langle \mathcal{S}, \mathcal{A}, \mathcal{I},
\mathcal{G}, \mathcal{T} \rangle$, où $\mathcal{S}$ est le système audité
(ensemble sociotechnique d'agents, d'artefacts, de flux et de dispositifs),
$\mathcal{A}$ l'architecture cognitive de l'auditeur ou de l'équipe,
$\mathcal{I}$ le canal d'information par lequel les preuves parviennent à
$\mathcal{A}$, $\mathcal{G}$ la structure stratégique de l'interaction
entre l'auditeur, l'audité et les tiers, et $\mathcal{T}$ la trajectoire
temporelle du système au fil des missions.

\medskip
\termecle{Auditologie}\index{auditologie}~: science appliquée de la
réduction de l'incertitude sur les missions auditologiques.
\end{boxdef}

\begin{tableaucours}{cL{3.6cm}L{8.4cm}}{%
    \tableauHeader{} &
    \tableauHeader{Composante} &
    \tableauHeader{Dans la mission MINÉQUIR}
}
$\mathcal{S}$ & Système audité & La DAF, SIDÉPENSE, ses 47 comptes, ses
procédures, le poste comptable \\
$\mathcal{A}$ & Architecture cognitive & L'équipe du CONSUPE, sa méthode,
ses papiers de travail \\
$\mathcal{I}$ & Canal d'information & Entretiens, pièces des trois strates,
extractions faites devant la vérificatrice \\
$\mathcal{G}$ & Structure stratégique & DAF, cellule informatique, DRH,
contrôle financier, vérificateurs, CDBF \\
$\mathcal{T}$ & Trajectoire & Mars~2026, novembre~2026, et les cycles
suivants \\
\end{tableaucours}
\vspace{-0.8em}
\begin{center}{\small\itshape Tableau~15.1. La mission MINÉQUIR comme
mission auditologique}\end{center}

\separateur

% ============================================================================
\section{Six axiomes}
\label{sec:c15-axiomes}
% ============================================================================

L'article A0 pose six axiomes, qui contraignent toute mission
auditologique~\cite{minka2026auditologie}. Le cours les a rencontrés
chapitre après chapitre~; leur numérotation est celle de l'article.

\begin{tableaucours}{cL{3.4cm}L{5.0cm}L{3.6cm}}{%
    \tableauHeader{N°} &
    \tableauHeader{Axiome} &
    \tableauHeader{Énoncé résumé} &
    \tableauHeader{Dans le cours}
}
1 & Observabilité partielle & L'auditeur n'accède qu'à une partie des états
du système, elle-même partiellement opaque & Partir du système~; la réponse
\NEK{} (Ch.~\ref{chap:preuve}) \\
2 & Cognition limitée & L'auditeur et l'audité sont des agents
cognitivement limités & Décomposition bornée à sept éléments
(Ch.~\ref{chap:preparer}) \\
3 & Mesurabilité & L'incertitude épistémique se mesure par l'entropie et
ses généralisations & $\mu$, $U$, $H$ (Ch.~\ref{chap:constater}) \\
4 & Interaction stratégique & Auditeur et audité ont des intérêts
distincts, dans un jeu à information imparfaite & Biais du choix forcé,
indice $\rho$~; indépendance \\
5 & Individuation temporelle & Le système évolue, et son état dépend des
missions antérieures & Suivi, divergence entre cycles \\
6 & Opposabilité & Toute conclusion doit pouvoir être défendue devant un
tiers légitime & Papiers de travail, contradictoire, observation en sept
éléments \\
\end{tableaucours}
\vspace{-0.8em}
\begin{center}{\small\itshape Tableau~15.2. Les six axiomes de
l'auditologie (d'après A0)}\end{center}

A0 écarte explicitement trois candidats~: l'\emph{indépendance} n'est pas un
axiome, mais une conséquence des axiomes~4 et~6~; la \emph{conformité à un
référentiel} est un choix de fonction objectif~; la \emph{matérialité} est
un paramètre dérivé de l'opposabilité.

\subsection{Quatre propositions}

A0 en tire quatre propositions, que ce cours a mises en pratique.
\begin{enumerate}
    \item \textbf{Modalité d'ignorance explicite} (axiomes~1 à~4)~: tout
          instrument de réponse valide inclut \NEK{}.
    \item \textbf{Décomposition hiérarchique} (axiomes~2 et~6)~: tout
          instrument couvrant un espace d'audit non trivial est organisé en
          niveaux, bornés par la mémoire de travail.
    \item \textbf{Métriques longitudinales} (axiome~5)~: les conclusions ne
          se testent que par l'observation de plusieurs cycles.
    \item \textbf{Réflexivité} (axiomes~4 et~6)~: l'audit fait partie de la
          trajectoire du système qu'il observe.
\end{enumerate}

\separateur

% ============================================================================
\section{Cinq piliers}
\label{sec:c15-piliers}
% ============================================================================

Les axiomes disent ce qu'est une mission~; les piliers apportent les théories
qui permettent de la modéliser. A0 en retient cinq, chacun avec une
prédiction falsifiable.

\begin{tableaucours}{L{2.6cm}L{3.6cm}L{6.0cm}}{%
    \tableauHeader{Pilier} &
    \tableauHeader{Ancrage} &
    \tableauHeader{Prédiction falsifiable}
}
1. Cognition limitée & Miller, Simon,
Kahneman~\cite{miller1956,simon1957,kahneman2011} & Les grilles de plus de
neuf dimensions simultanées produisent plus d'incohérences et d'oublis \\
2. Théorie de l'information & Shannon, Kullback et
Leibler~\cite{shannon1948,kullback1951} & Le diagnostic conformité et
entropie prédit mieux que le seul taux de conformité \\
3. Théorie des jeux & von~Neumann, Nash,
Harsanyi~\cite{vonneumann1944,nash1951,harsanyi1967} & Les instruments
compatibles avec les incitations réduisent la dissimulation \\
4. Systèmes dynamiques & Chaînes de Markov~\cite{norris1998} & Les
prévisions par matrice de transition entre niveaux de maturité battent la
persistance naïve \\
5. Ontologie critique-réaliste & Bhaskar, Simondon,
Latour~\cite{bhaskar1975,simondon1958,latour2005} & Les cadres qui tiennent
compte des personnes, des procédures et des artefacts font mieux appliquer
les recommandations \\
\end{tableaucours}
\vspace{-0.8em}
\begin{center}{\small\itshape Tableau~15.3. Les cinq piliers et leurs
prédictions (d'après A0)}\end{center}

\separateur

% ============================================================================
\section{La méthode : science du design et falsifiabilité}
\label{sec:c15-methode}
% ============================================================================

\subsection{Une science du design}

L'auditologie produit des artefacts~: instruments, métriques, grilles,
outils. A0 l'inscrit dans la recherche en science du design~: on construit
un artefact, on l'évalue dans son environnement, on l'appuie sur la base de
connaissances, et l'on recommence~\cite{hevner2004design}. Les outils de ce
cours (l'instrument \OUI{}~/ \NON{}~/ \NEK{}, le gabarit canonique, les
labs de NDAP SUKULU) sont de tels artefacts, et chaque promotion
d'auditeurs participe à leur évaluation.

\subsection{Un noyau dur et une ceinture protectrice}

A0 adopte le falsificationnisme sophistiqué de Lakatos~\cite{lakatos1970}.
Les axiomes et les piliers forment un \emph{noyau dur}, qui n'est pas
directement testé~; les hypothèses auxiliaires, les procédures de mesure et
les prédictions forment la \emph{ceinture protectrice}, qui l'est. La
réfutation d'une prédiction n'emporte pas le programme~: elle oblige à
réviser la ceinture.

\subsection{Ce que l'auditologie n'est pas}

A0 précise quatre refus~: l'auditologie n'est pas \emph{normative} (elle ne
prescrit pas le bon audit), ni \emph{critique} (elle n'invalide pas les
cadres existants), ni \emph{un référentiel de plus}, ni \emph{une idéologie}
(ses propositions sont falsifiables). C'est pourquoi le cours a enseigné
ensemble les normes du métier (Partie~II) et leur lecture auditologique
(Partie~III).

\separateur

% ============================================================================
\section{Établi, déposé, à valider}
\label{sec:c15-statut}
% ============================================================================

\begin{tableaucours}{L{3.0cm}L{5.0cm}L{4.2cm}}{%
    \tableauHeader{Statut} &
    \tableauHeader{Contenu} &
    \tableauHeader{Références}
}
Établi par des travaux publiés & Borne de la mémoire de travail, entropie
de Shannon, divergence de Kullback-Leibler, théorie des jeux, contrôle
interne & Miller, Shannon, Kullback, Nash, COSO \\
Déposé en prépublication & Axiomes, piliers, mission auditologique,
instrument \OUI{}~/ \NON{}~/ \NEK{}, mesures, quadrant & A0 et A1, sur
HAL~\cite{minka2026auditologie,minka2026nek} \\
À valider & Les prédictions des cinq piliers, sur des missions réelles et
répétées & Les recherches à venir, dont les vôtres \\
\end{tableaucours}
\vspace{-0.8em}
\begin{center}{\small\itshape Tableau~15.4. Le statut des travaux sur
lesquels repose le cours}\end{center}

\separateur

% ============================================================================
\section{Axes de recherche pour l'auditeur public}
\label{sec:c15-axes}
% ============================================================================

Les auditeurs des finances publiques disposent de ce qui manque le plus à la
recherche en audit~: des missions réelles, sur des systèmes réels, répétées
dans le temps. Huit axes s'ouvrent à partir de ce cours. Pour chacun, un
exemple de sujet de mémoire est donné, avec sa question falsifiable.

\begin{tableaucours}{cL{3.2cm}L{5.0cm}L{3.0cm}}{%
    \tableauHeader{Axe} &
    \tableauHeader{Thème} &
    \tableauHeader{Question de recherche} &
    \tableauHeader{Pilier, axiome}
}
1 & Le biais du choix forcé & De combien les grilles binaires des
inspections surestiment-elles la conformité~? & Pilier~3~; axiome~4 \\
2 & L'incertitude dans les rapports publics & Quelle part des conclusions
des rapports de contrôle repose sur des déclarations sans preuve~? &
Pilier~2~; axiome~1 \\
3 & La valeur prédictive du quadrant & Les domaines en conformité opaque
révèlent-ils plus d'irrégularités aux missions suivantes~? & Pilier~2~;
axiome~3 \\
4 & Trajectoires des recommandations & Le suivi des recommandations
suit-il une chaîne de Markov prévisible~? & Pilier~4~; axiome~5 \\
5 & L'usage stratégique de \NEK{} & Les audités répondent-ils davantage
\NEK{} en entretien formel qu'en auto-évaluation~? & Pilier~3~; axiome~4 \\
6 & La décomposition et l'erreur & Les grilles qui dépassent la borne de
sept produisent-elles plus d'oublis~? & Pilier~1~; axiome~2 \\
7 & Contrôle interne des applications de l'État & Quelles déficiences
reviennent dans les applications financières publiques, et pourquoi~? &
Pilier~5~; axiome~1 \\
8 & L'opposabilité de la preuve électronique & Quelles conditions rendent
une extraction recevable devant le CDBF et la Chambre des Comptes~? &
Axiome~6 \\
\end{tableaucours}
\vspace{-0.8em}
\begin{center}{\small\itshape Tableau~15.5. Huit axes de recherche ouverts
par le cours}\end{center}

\subsection{Exemples de sujets}

\begin{boxscenario}{Axe 1. Mesurer le biais des grilles binaires}
\textbf{Sujet.} \emph{Le biais du choix forcé dans les questionnaires de
contrôle interne des administrations financières camerounaises~: mesure par
l'indice $\rho$.}

\textbf{Hypothèse falsifiable.} Sur les mêmes questions et les mêmes
répondants, la part de \OUI{} d'une grille binaire dépasse d'au moins dix
points celle de l'instrument trichotomique.

\textbf{Données.} Deux passages d'un même questionnaire (binaire puis
trichotomique) dans trois à cinq services, avec l'accord de la hiérarchie.

\textbf{Méthode.} Calcul de $\rho$ par question et par domaine~; test
statistique de l'écart~; analyse des questions où le biais est le plus fort.

\textbf{Réfutation.} Un $\rho$ proche de zéro sur l'ensemble des services.
\end{boxscenario}

\begin{boxscenario}{Axe 3. La valeur prédictive du quadrant}
\textbf{Sujet.} \emph{Conformité opaque et irrégularités ultérieures~: étude
longitudinale de missions de vérification.}

\textbf{Hypothèse falsifiable.} Les domaines classés en conformité opaque à
une mission présentent, à la mission suivante, plus d'irrégularités établies
que les domaines en conformité claire.

\textbf{Données.} Les mesures $\mu$, $U$, $H$ de domaines audités sur deux
cycles au moins, et les constats de la mission suivante.

\textbf{Méthode.} Comparaison des taux d'irrégularités par région~; contrôle
de la taille des domaines.

\textbf{Réfutation.} Aucune différence entre régions, ou une différence
inverse.
\end{boxscenario}

\begin{boxscenario}{Axe 4. Trajectoires des recommandations}
\textbf{Sujet.} \emph{Le suivi des recommandations de contrôle comme
processus de Markov~: estimation d'une matrice de transition des niveaux de
maturité.}

\textbf{Hypothèse falsifiable.} Une matrice de transition estimée sur trois
cycles prévoit le niveau de maturité au cycle suivant mieux que la simple
reconduction du niveau actuel.

\textbf{Données.} Les niveaux de maturité (échelle du
Chapitre~\ref{chap:eval-sci}) de plusieurs domaines, sur au moins quatre
cycles.

\textbf{Méthode.} Estimation de la matrice, prévision, comparaison à la
persistance naïve.

\textbf{Réfutation.} Aucun avantage prédictif de la matrice.
\end{boxscenario}

\begin{boxscenario}{Axe 8. L'opposabilité de la preuve électronique}
\textbf{Sujet.} \emph{Les conditions de recevabilité des extractions
informatiques dans les procédures du CDBF et de la Chambre des Comptes.}

\textbf{Question.} Quelles garanties (origine, exhaustivité, exactitude,
intégrité) sont effectivement exigées, et lesquelles manquent le plus
souvent~?

\textbf{Données.} Un échantillon de dossiers, anonymisés, avec l'autorisation
des institutions~; entretiens avec des rapporteurs.

\textbf{Méthode.} Grille d'analyse fondée sur les quatre garanties du
Chapitre~\ref{chap:preuve}~; étude des cas où une pièce a été écartée.
\end{boxscenario}

\separateur

% ============================================================================
\section{Conduire un mémoire de recherche appliquée}
\label{sec:c15-memoire}
% ============================================================================

\begin{boxprocedure}
\textbf{Les étapes d'un mémoire en auditologie appliquée}~:
\begin{enumerate}
    \item partir d'une difficulté réelle de mission, et la rattacher à un
          axiome ou à un pilier~;
    \item formuler une hypothèse \emph{falsifiable}~: dire à l'avance quel
          résultat la réfuterait~;
    \item obtenir l'autorisation écrite de la hiérarchie et des entités
          concernées~;
    \item construire le protocole~: instrument, échantillon, mesures
          ($\mu$, $U$, $H$, $\rho$, $D_{\mathrm{KL}}$), seuils déclarés~;
    \item anonymiser les données personnelles et protéger les extractions~;
    \item analyser, puis séparer ce qui est établi, ce qui est probable et ce
          qui reste inconnu~;
    \item discuter les limites, et proposer la révision de la ceinture
          protectrice si la prédiction est réfutée.
\end{enumerate}
\end{boxprocedure}

\begin{boxalert}
\textbf{Violation typique.} Formuler un sujet du type «~l'importance de
l'audit des SI dans l'administration camerounaise~».

\smallskip
\textbf{Règle~:} un sujet de recherche pose une question à laquelle des
données peuvent répondre non. «~Importance~» ne se réfute pas~; «~les
domaines en conformité opaque révèlent plus d'irrégularités au cycle
suivant~» se réfute.
\end{boxalert}

\begin{boxmetier}{Vérificateur}
Vos missions répétées sur les mêmes entités sont des données longitudinales
rares. Conservez, d'un cycle à l'autre, les mêmes questions et les mêmes
fonctions répondantes~: vous rendez possibles les axes~3 et~4.
\end{boxmetier}

\begin{boxmetier}{Auditeur interne}
Mesurer $\rho$ sur vos propres questionnaires est un premier travail de
recherche, utile dès la première année.
\end{boxmetier}

\separateur

\begin{boxresume}
\begin{itemize}
    \item \textbf{L'auditologie} est la science appliquée de la réduction de
          l'incertitude~; son objet est la mission
          $\langle \mathcal{S}, \mathcal{A}, \mathcal{I}, \mathcal{G},
          \mathcal{T} \rangle$.
    \item \textbf{Six axiomes}~: observabilité partielle, cognition limitée,
          mesurabilité, interaction stratégique, individuation temporelle,
          opposabilité.
    \item \textbf{Quatre propositions} fondent les pratiques du cours~:
          \NEK{}, décomposition, mesures longitudinales, réflexivité.
    \item \textbf{Cinq piliers}, chacun avec une prédiction falsifiable.
    \item \textbf{Une science du design, à la Lakatos}~: un noyau dur, une
          ceinture protectrice qui se teste.
    \item \textbf{Un statut clair}~: A0 et A1 sont des prépublications~; les
          prédictions restent à valider.
    \item \textbf{Huit axes} s'ouvrent aux auditeurs publics~; un bon sujet
          dit à l'avance ce qui le réfuterait.
\end{itemize}
\end{boxresume}

\begin{boxexo}[title={\ding{45}\quad Exercice~15.1. Axiomes et pratiques \niveaubase}]
Pour chacune des pratiques suivantes, citez l'axiome d'A0 qui la fonde~:
(a) coter \NEK{} une pièce non encore remise~; (b) borner chaque niveau de
la décomposition à sept éléments~; (c) comparer deux cycles par la
divergence~; (d) soumettre le rapport provisoire à l'entité.
\end{boxexo}

\begin{boxexo}[title={\ding{45}\quad Exercice~15.2. Rendre un sujet falsifiable \niveaumoyen}]
Réécrivez les sujets suivants pour qu'ils deviennent falsifiables, en
précisant l'hypothèse, les données et la condition de réfutation~: (a)
«~Le rôle du CONSUPE dans la bonne gouvernance numérique~»~; (b) «~Les
enjeux de la dématérialisation des recettes~».
\end{boxexo}

\begin{boxexo}[title={\ding{45}\quad Exercice~15.3. Un protocole \niveauexpert}]
Choisissez l'un des huit axes du Tableau~15.5 et rédigez le protocole
complet d'un mémoire~: question, hypothèse, données et autorisations,
instrument, mesures, seuils, analyse, condition de réfutation, risques
éthiques et parade.
\end{boxexo}

\begin{boxinfo}
\textbf{Pour finir.} Le cours s'arrête ici~; le métier, et peut-être la
recherche, commencent. Chaque mission refera les mêmes gestes~: partir du
système, écrire ses critères avant de tester, coter honnêtement, y compris
\NEK{}, conclure sur ce que l'on a prouvé, et mesurer, d'un cycle à
l'autre, ce que l'on a appris.
\end{boxinfo}
}{}

\appendix

    %% ============================================================================
%% Fichier  : annexes/Ann_A_glossaire.tex
%% Titre    : Glossaire
%% Version  : 0.1 -- novembre 2026
%% ============================================================================
%% Définitions abrégées des termes marqués \termecle dans les chapitres.
%% ============================================================================

\chapter{Glossaire}
\label{ann:glossaire}

Ce glossaire reprend, sous une forme abrégée, les définitions données dans
les chapitres. Le renvoi indique le chapitre où le terme est défini.

\begin{description}
    \item[\NEK{} (\textit{Not Enough Knowledge})] Cote réservée au contrôle
          dont l'auditeur ne peut pas établir l'état malgré des diligences
          normales (Ch.~\ref{chap:prologue}, \ref{chap:eval-sci}).
    \item[Analyse] Interprétation tirée de constatations au regard d'un
          critère, avec un niveau de confiance déclaré
          (Ch.~\ref{chap:prologue}).
    \item[Assurance limitée] Confiance plus faible, exprimée sous forme
          négative, après des travaux réduits (Ch.~\ref{chap:audit-public}).
    \item[Assurance raisonnable] Confiance élevée mais non absolue, exprimée
          sous forme positive (Ch.~\ref{chap:audit-public}).
    \item[Audit] Examen méthodique, indépendant et documenté, qui compare
          une situation à des critères et conduit à une opinion motivée
          (Ch.~\ref{chap:audit-public}).
    \item[Changement] Toute modification d'un programme, d'un paramètre ou
          d'une base en production (Ch.~\ref{chap:exploitation}).
    \item[Constatation] Fait vérifiable, établi par l'auditeur et appuyé
          sur une preuve conservée (Ch.~\ref{chap:prologue}).
    \item[Contrôle] Dispositif ou action de l'entité qui empêche ou détecte
          une anomalie (Ch.~\ref{chap:audit-public}).
    \item[Contrôle compensatoire] Revue indépendante qui couvre un cumul de
          profils que l'on ne peut pas éviter (Ch.~\ref{chap:sci}).
    \item[Contrôle correctif, détectif, préventif] Contrôle qui corrige,
          repère après coup, ou empêche une anomalie (Ch.~\ref{chap:sci}).
    \item[Contrôle interne] Processus qui fournit une assurance raisonnable
          sur les opérations, l'information et la conformité, selon le
          COSO~2013 (Ch.~\ref{chap:sci}).
    \item[Contrôles d'application] Contrôles intégrés à une application~:
          saisie, traitement, interface, sortie (Ch.~\ref{chap:sci},
          \ref{chap:applications}).
    \item[Contrôles généraux informatiques] Contrôles de l'environnement des
          applications~: accès, changements, exploitation, acquisition
          (Ch.~\ref{chap:sci}).
    \item[Critère d'audit] Exigence de référence, choisie avant le test
          (Ch.~\ref{chap:audit-public}, \ref{chap:normes}).
    \item[Déficience] Insuffisance d'une composante ou d'un principe qui
          réduit la probabilité d'atteindre les objectifs
          (Ch.~\ref{chap:eval-sci}).
    \item[Déficience majeure] Déficience, ou combinaison, qui réduit
          fortement cette probabilité~; elle interdit de conclure à
          l'efficacité (Ch.~\ref{chap:eval-sci}).
    \item[Échantillon] Partie de la population examinée en détail,
          aléatoire ou raisonnée (Ch.~\ref{chap:preuve}).
    \item[Empreinte] Valeur calculée par une fonction de hachage (SHA-256)
          qui change au moindre changement du fichier (Ch.~\ref{chap:preuve}).
    \item[Fraude] Acte intentionnel destiné à obtenir un avantage indu
          (Ch.~\ref{chap:donnees-fraude}).
    \item[Gestion des changements] Contrôles qui garantissent qu'un
          changement est demandé, testé, approuvé, mis en production par un
          tiers et tracé (Ch.~\ref{chap:exploitation}).
    \item[Inspection] Vérification ponctuelle au nom de l'autorité
          hiérarchique (Ch.~\ref{chap:audit-public}).
    \item[Limitation d'étendue] Partie de l'objectif que la mission n'a pas
          pu établir, déclarée au rapport (Ch.~\ref{chap:constater}).
    \item[Matrice des risques et des contrôles] Liste des risques d'un
          circuit et des contrôles attendus en face, établie avant les tests
          (Ch.~\ref{chap:eval-sci}).
    \item[Normalisation] Transformation d'une donnée avant comparaison
          (Ch.~\ref{chap:applications}).
    \item[Norme professionnelle] Règle qui encadre la conduite d'une mission
          (Ch.~\ref{chap:normes}).
    \item[Plan de continuité d'activité] Document qui fixe comment l'entité
          poursuit ses missions pendant une interruption
          (Ch.~\ref{chap:exploitation}).
    \item[Population] Ensemble des opérations sur lesquelles porte la
          conclusion, prise dans le système (Ch.~\ref{chap:preuve}).
    \item[Profil] Ensemble de droits attribué dans une application
          (Ch.~\ref{chap:sci}).
    \item[Propriétaire d'application] Responsable métier qui décide des
          droits dans une application (Ch.~\ref{chap:acces}).
    \item[Questionnaire de contrôle interne] Questions fermées, tirées de la
          matrice, qui recueillent les contrôles déclarés
          (Ch.~\ref{chap:eval-sci}).
    \item[Référentiel] Ensemble structuré de bonnes pratiques où l'on puise
          des critères (Ch.~\ref{chap:normes}).
    \item[Revue des habilitations] Examen périodique et signé des comptes et
          profils par le propriétaire de l'application (Ch.~\ref{chap:acces}).
    \item[Risque inhérent, risque lié au contrôle] Gravité et probabilité
          d'une anomalie sans contrôle~; probabilité que le contrôle ne
          l'empêche ni ne la détecte (Ch.~\ref{chap:preparer}).
    \item[Séparation des tâches] Répartition des étapes d'une opération
          entre des personnes différentes (Ch.~\ref{chap:sci}).
    \item[Système de contrôle interne] Ensemble organisé des dispositifs de
          contrôle interne d'une entité, évalué comme un tout
          (Ch.~\ref{chap:sci}).
    \item[Test de cheminement] Suivi d'une opération réelle du début à la
          fin du circuit (Ch.~\ref{chap:eval-sci}).
    \item[Totaux de contrôle] Nombre de lignes et sommes rapprochés d'une
          source indépendante pour prouver l'exhaustivité
          (Ch.~\ref{chap:preuve}).
    \item[Triangle de la fraude] Pression, occasion, rationalisation
          (Ch.~\ref{chap:donnees-fraude}).
    \item[Vérification] Activité des vérificateurs du CONSUPE, sur pièces et
          sur place (Ch.~\ref{chap:audit-public}).
    \item[Auditologie] Science appliquée de la réduction de l'incertitude
          sur les missions d'audit, selon l'article A0
          (Ch.~\ref{chap:recherche}).
    \item[Divergence de Kullback-Leibler] Mesure, en bits, de l'écart entre
          les distributions de réponses de deux cycles
          (Ch.~\ref{chap:constater}).
    \item[Entropie] Dispersion des réponses \OUI{}, \NON{}, \NEK{}, en bits,
          au plus $\log_2 3$ (Ch.~\ref{chap:constater}).
    \item[Espace d'audit] Ensemble des objets sur lesquels porte une
          mission (Ch.~\ref{chap:preparer}).
    \item[Fiche d'assertion] Objectif, critère, question, risque et
          conséquence d'un contrôle (Ch.~\ref{chap:preparer}).
    \item[Gabarit canonique] Décomposition en huit niveaux, du domaine à la
          trace (Ch.~\ref{chap:preparer}).
    \item[Mission auditologique] Quintuplet $\langle \mathcal{S},
          \mathcal{A}, \mathcal{I}, \mathcal{G}, \mathcal{T} \rangle$ de
          l'article A0 (Ch.~\ref{chap:recherche}).
    \item[Observation] Développement du constat en sept éléments~: norme,
          constatation, risques, contradictoire, analyse, conclusion,
          recommandations (Ch.~\ref{chap:constater}).
    \item[Quadrant diagnostique] Lecture croisée de $\mu$ et de $U$ en quatre
          régions~; la conformité opaque est la plus dangereuse
          (Ch.~\ref{chap:constater}).
    \item[Taux de conformité ($\mu$), taux d'incertitude ($U$)] Parts des
          réponses \OUI{} et \NEK{} (Ch.~\ref{chap:constater}).
\end{description}

    %% ============================================================================
%% Fichier  : annexes/Ann_B_outils.tex
%% Titre    : Référentiel des Outils
%% Version  : 0.1 -- novembre 2026
%% ============================================================================

\chapter{Référentiel des Outils}
\label{ann:outils}

Le cours ne demande qu'un tableur (Excel ou LibreOffice Calc) et l'outil de
calcul d'empreinte du système. Les formules sont données avec les noms
français des fonctions et le point-virgule comme séparateur.

\section*{Intégrité des fichiers}

\begin{tableaucours}{L{3.6cm}L{8.6cm}}{%
    \tableauHeader{Besoin} &
    \tableauHeader{Commande}
}
Empreinte sous Windows & \texttt{certutil -hashfile fichier.csv SHA256} \\
Empreinte sous Linux & \texttt{sha256sum fichier.csv} \\
Vérifier un manifeste & \texttt{sha256sum -c MANIFESTE\_SHA256.txt} \\
\end{tableaucours}

\section*{Formules de tableur}

\begin{tableaucours}{L{3.8cm}L{5.8cm}L{2.6cm}}{%
    \tableauHeader{Test} &
    \tableauHeader{Formule} &
    \tableauHeader{Chapitre}
}
Doublons sur une clé & \texttt{=NB.SI(\$F\$2:\$F\$9612;F2)} &
\ref{chap:prologue} \\
Tirage aléatoire & \texttt{=ALEA()} puis tri, valeurs figées &
\ref{chap:preuve} \\
Existence dans une autre liste & \texttt{=SIERREUR(RECHERCHEV(A2;
PERSONNEL!\$A:\$D;4;FAUX);"ABSENT")} & \ref{chap:acces} \\
Conflit de profils & \texttt{=SI(ET(C2=1;D2=1);"CONFLIT";"")} &
\ref{chap:acces} \\
Suite d'échecs & \texttt{=SI(B2="ECHEC";C1+1;0)} puis
\texttt{=MAX(C:C)} & \ref{chap:exploitation} \\
Seuil sans visa & \texttt{=SI(ET(F2>50000000;G2<>"OUI");"SANS VISA";"")} &
\ref{chap:applications} \\
Âge & \texttt{=DATEDIF(C2;DATE(2025;12;31);"y")} &
\ref{chap:donnees-fraude} \\
Premier chiffre & \texttt{=GAUCHE(TEXTE(A2;"0");1)*1} &
\ref{chap:donnees-fraude} \\
Proportion de Benford & \texttt{=LOG10(1+1/d)} & \ref{chap:donnees-fraude}
\\
\end{tableaucours}

\section*{NDAP SUKULU}

Les labs se font dans le bureau d'analyse de NDAP SUKULU, ouvert dans le
navigateur. Il contient un tableur, un éditeur de texte et un terminal.
Les pièces sont en lecture seule dans \texttt{preuves}~; on travaille sur
la copie du dossier \texttt{travail}.

    %% ============================================================================
%% Fichier  : annexes/Ann_C_templates.tex
%% Titre    : Papiers de Travail Types
%% Version  : 0.1 -- novembre 2026
%% ============================================================================

\chapter{Papiers de Travail Types}
\label{ann:templates}

Les modèles ci-dessous reprennent les encadrés «~Papier de travail~» des
chapitres. Ils sont à recopier dans le tableur ou le traitement de texte
de l'institution.

\section*{C.1 Ordre de mission}

\begin{tableaucours}{L{3.6cm}L{8.6cm}}{%
    \tableauHeader{Rubrique} &
    \tableauHeader{À remplir}
}
Fondement & Texte, article, alinéa \\
Objectif & Une question à laquelle le rapport répondra \\
Périmètre & Entité, applications, ce qui est exclu \\
Période & Exercice ou dates \\
Durée et équipe & Dates, noms, chef de mission \\
Livrables & Rapport provisoire, contradictoire, rapport définitif \\
\end{tableaucours}

\section*{C.2 Demande de pièces}

\begin{tableaucours}{cL{5.6cm}L{3.0cm}L{1.6cm}L{1.6cm}}{%
    \tableauHeader{N°} &
    \tableauHeader{Pièce (extraite devant l'auditeur~?)} &
    \tableauHeader{Destinataire} &
    \tableauHeader{Délai} &
    \tableauHeader{Reçue le}
}
1 & & & & \\
2 & & & & \\
\end{tableaucours}

\section*{C.3 Questionnaire de contrôle interne}

\begin{tableaucours}{cL{4.6cm}L{2.6cm}L{2.4cm}L{2.0cm}}{%
    \tableauHeader{N°} &
    \tableauHeader{Question (qui, quand, quelle trace~?)} &
    \tableauHeader{Réponse déclarée} &
    \tableauHeader{Répondant} &
    \tableauHeader{Pièce}
}
1 & & & & \\
2 & & & & \\
\end{tableaucours}

\section*{C.4 Feuille de test}

\begin{tableaucours}{L{3.6cm}L{8.6cm}}{%
    \tableauHeader{Rubrique} &
    \tableauHeader{À remplir}
}
Référence & Domaine, contrôle, ligne du programme \\
Critère & Énoncé, source, force \\
Population & Source, nombre, totaux, empreinte \\
Échantillon & Méthode, taille, liste \\
Résultats & Élément par élément, avec la pièce \\
Cote & \OUI{}~/ \NON{}~/ \NEK{}, et constatation \\
Visa & Auteur, date, chef de mission \\
\end{tableaucours}

\section*{C.5 Fiche de constat}

\begin{tableaucours}{L{3.0cm}L{9.2cm}}{%
    \tableauHeader{Élément} &
    \tableauHeader{À remplir}
}
Condition & Fait chiffré, daté, sourcé \\
Critère & Règle citée et sa force \\
Cause & Ce qui, dans le dispositif, explique l'écart \\
Effet & Risque réalisé ou possible, chiffré \\
Recommandation & Action, responsable, échéance \\
Classement & Déficience ou déficience majeure, et pourquoi \\
\end{tableaucours}

\section*{C.6 Grille de classement des déficiences}

\begin{tableaucours}{L{2.6cm}L{2.2cm}L{2.2cm}L{2.2cm}L{1.6cm}L{1.4cm}}{%
    \tableauHeader{Déficience} &
    \tableauHeader{Perte possible} &
    \tableauHeader{Probabilité} &
    \tableauHeader{Compensation prouvée} &
    \tableauHeader{Fraude} &
    \tableauHeader{Classe}
}
A & & & & & \\
B & & & & & \\
\end{tableaucours}

\section*{C.7 Plan d'action et suivi}

\begin{tableaucours}{cL{4.6cm}L{2.4cm}L{1.8cm}L{1.6cm}L{1.6cm}}{%
    \tableauHeader{N°} &
    \tableauHeader{Recommandation} &
    \tableauHeader{Responsable} &
    \tableauHeader{Échéance} &
    \tableauHeader{Statut} &
    \tableauHeader{Preuve}
}
R-01 & & & & & \\
R-02 & & & & & \\
\end{tableaucours}

    %% ============================================================================
%% Fichier  : annexes/Ann_D_textes.tex
%% Titre    : Textes Juridiques et Normatifs
%% Version  : 0.1 -- novembre 2026
%% ============================================================================
%% Textes lus dans le dossier Normes : décret n° 2013/287
%% (texte_organique_consupe.pdf) ; loi n° 2010/012. Les autres textes sont
%% décrits sans citation.
%% ============================================================================

\chapter{Textes Juridiques et Normatifs}
\label{ann:textes}

Les extraits ci-dessous sont recopiés de la version des textes disponible
dans le dossier des normes du cours. L'auditeur les vérifie sur la version
officielle avant de les citer dans un rapport.

\section*{D.1 Décret n° 2013/287 du 4 septembre 2013 portant organisation
des services du Contrôle supérieur de l'État}

\begin{boxjuris}
\textbf{Article~1er}~: \textit{« (1) Les services du Contrôle supérieur de
l'Etat relèvent de l'autorité directe du Président de la République, dont
ils reçoivent les instructions et à qui ils rendent compte. (2) Ils sont
dirigés par un ministre délégué à la présidence de la République. »}

\medskip
\textbf{Article~2, alinéa~1} (début)~: \textit{« Les services du Contrôle
supérieur de l'Etat constituent l'institution supérieure de contrôle des
finances publiques (ISC) du Cameroun. Ils sont chargés de l'audit
externe. »}

\medskip
\textbf{Article~2, alinéa~3}~: \textit{« Dans l'exercice de leurs
attributions, les services du Contrôle supérieur de l'Etat effectuent~:
un contrôle de conformité et de régularité~; un contrôle financier~; un
contrôle de performance~; un contrôle de l'environnement~; un contrôle des
systèmes d'information. »}

\medskip
\textbf{Article~2, alinéa~4}~: \textit{« Les contrôles sont effectués par
les vérificateurs du Contrôle supérieur de l'Etat sur pièces et sur place.
A cet effet, et nonobstant les contrôles effectués par les autres organes
de contrôle, les responsables des structures auditées sont tenus de
présenter les pièces relatives à la période auditée. »}

\medskip
\textbf{Article~3, alinéa~1}~: \textit{« Les services du Contrôle supérieur
de l'Etat concourent à la sanction des ordonnateurs et des gestionnaires
des deniers publics, dans les conditions prévues par les lois et
règlements en vigueur. A ce titre, le ministre délégué en charge du
Contrôle supérieur de l'Etat préside le conseil de discipline budgétaire
et financière, dont l'organisation et le fonctionnement sont fixés par un
texte particulier. »}

\medskip
\textbf{Article~5}~: \textit{« (1) Les structures d'audit interne
transmettent aux services du Contrôle supérieur de l'Etat, à la diligence
des chefs de départements ministériels concernés, un exemplaire de chacun
des rapports rédigés à la suite de toutes les interventions menées au
cours de l'exercice, ainsi qu'un exemplaire du rapport annuel d'activités.
(2) Après examen des rapports visés à l'alinéa 1er ci-dessus, les services
du Contrôle supérieur de l'Etat peuvent, le cas échéant, saisir le
conseil de discipline budgétaire et financière. »}

\medskip
\textbf{Article~27, alinéa~2} (extrait)~: la Division des inspections et
de contrôle des administrations publiques comprend notamment \textit{« la
section des administrations financières, économiques, sociales et des
systèmes d'informations »}.

\medskip
\textbf{Article~28, alinéa~1}~: les sections sont chargées, chacune dans
son domaine de compétence, \textit{« de la préparation et de l'exécution
des opérations de vérification~; du suivi et de l'évaluation des
activités des structures de contrôle interne~; de la tenue de la
documentation technique et du fichier. »}
\end{boxjuris}

\section*{D.2 Loi n° 2010/012 du 21 décembre 2010 relative à la
cybersécurité et à la cybercriminalité}

\begin{boxjuris}
\textbf{Article~4, point~9}~: \textit{« Audit de sécurité~: examen
méthodique des composantes et des acteurs de la sécurité, de la
politique, des mesures, des solutions, des procédures et des moyens mis
en oeuvre par une organisation, pour sécuriser son environnement,
effectuer des contrôles de conformité, des contrôles d'évaluation de
l'adéquation des moyens (organisationnels, techniques, humains,
financiers) investis au regard des risques encourus, d'optimisation, de
rationalité et de performance. »}

\medskip
\textbf{Article~7, alinéa~2} (extrait)~: l'Agence est notamment chargée
\textit{« de s'assurer de la régularité, de l'effectivité des audits de
sécurité des systèmes d'information suivant les normes en la matière, des
organismes publics et des autorités de certification »}.

\medskip
\textbf{Article~13, alinéa~1}~: \textit{« Sont soumis à un audit de
sécurité obligatoire, les réseaux de communications électroniques et les
systèmes d'information, des opérateurs, des autorités de certification et
des fournisseurs de services de communications électroniques. »}

\medskip
\textbf{Article~32}~: \textit{« (1) Les réseaux de communications
électroniques et les systèmes d'information sont soumis à un régime
d'audit de sécurité obligatoire et périodique de leurs systèmes de
sécurité par l'Agence. (2) L'audit de sécurité et les mesures d'impact de
gravité sont effectués chaque année ou lorsque les circonstances
l'exigent. (3) Les rapports d'audit sont confidentiels et adressés au
Ministre chargé des Télécommunications. »}
\end{boxjuris}

\section*{D.3 Textes décrits sans citation}

\begin{tableaucours}{L{4.8cm}L{7.4cm}}{%
    \tableauHeader{Texte} &
    \tableauHeader{Usage dans le cours}
}
Loi n°~2018/012 du 11~juillet~2018 portant régime financier de l'État et
des autres entités publiques & Exécution de la dépense, séparation
ordonnateur / comptable (Ch.~\ref{chap:sci}) \\
Texte organisant la Chambre des Comptes de la Cour suprême & Contrôle des
comptes~; les SI par extension (Ch.~\ref{chap:institutions}) \\
Loi n°~2024/017 du 23~décembre~2024 relative à la protection des données
personnelles & Données des agents, contribuables, usagers
(Ch.~\ref{chap:gouvernance}) \\
\end{tableaucours}

\section*{D.4 Normes}

\begin{tableaucours}{L{4.4cm}L{7.8cm}}{%
    \tableauHeader{Norme} &
    \tableauHeader{Passages utilisés}
}
COSO~2013 (traduction IFACI) & Définition, composantes et principes,
efficacité, déficiences (§§~138 à~140) \\
INTOSAI GUID~5100 & Audit des SI autonome ou intégré~; treize domaines \\
Manuel d'audit de conformité de l'IDI & Relation tripartite, critères,
documentation \\
Normes de l'IIA (2017, 2024) & Exigences de l'audit interne \\
COBIT~2019, ISO/IEC~27001 et 27002 & Sources de critères (résumés, aucune
reproduction) \\
\end{tableaucours}

    %% ============================================================================
%% Fichier  : annexes/Ann_E_bibliographie.tex
%% Titre    : Bibliographie Commentée
%% Version  : 0.1 -- novembre 2026
%% ============================================================================

\chapter{Bibliographie Commentée}
\label{ann:biblio}

Cette annexe dit, pour chaque référence principale, ce qu'elle apporte à
l'auditeur public et par où commencer. La liste complète figure dans les
références bibliographiques à la fin du livre.

\section*{Pour l'évaluation du contrôle interne}

\begin{itemize}
    \item \textbf{COSO, \emph{Référentiel intégré de contrôle interne}}
          (2013, traduction IFACI 2014)~\cite{coso2013icfr}. Le texte de
          base des Chapitres~\ref{chap:sci} et~\ref{chap:eval-sci}. Lire
          d'abord la synthèse, puis le chapitre sur l'efficacité du
          contrôle interne.
    \item \textbf{COSO et ACFE, guide de gestion du risque de fraude}
          (2016)~\cite{coso2016fraude}. La démarche du principe~8, utile
          pour le Chapitre~\ref{chap:donnees-fraude}.
\end{itemize}

\section*{Pour l'audit public}

\begin{itemize}
    \item \textbf{INTOSAI, GUID~5100}~\cite{intosai2019guid5100}. Les
          lignes directrices de la vérification des SI par les
          institutions supérieures de contrôle~; les treize domaines
          d'une mission.
    \item \textbf{IDI, manuel d'audit de conformité}~\cite{idi2021conformite}.
          Le plus pratique des textes de l'INTOSAI~: critères, preuve,
          documentation, rapport.
    \item \textbf{Décret n°~2013/287}~\cite{decret2013287consupe}. À
          connaître par cœur pour ses articles~2, 3, 5 et~28.
\end{itemize}

\section*{Pour l'audit interne et les référentiels}

\begin{itemize}
    \item \textbf{IIA, normes}~\cite{iia2017cripp,iia2024gias}. Les
          exigences de l'auditeur interne~; l'édition 2024 remplace celle
          de 2017.
    \item \textbf{ISACA, COBIT~2019}~\cite{isaca2018cobitintro,isaca2018cobitobjectifs}.
          À consulter comme catalogue de critères, objectif par objectif.
    \item \textbf{ISO/IEC~27001 et 27002}~\cite{iso27001_2022,iso27002_2022}.
          Les mesures de sécurité, à citer quand l'entité les a adoptées.
\end{itemize}

\section*{Pour l'analyse de données}

\begin{itemize}
    \item \textbf{Nigrini, \emph{Benford's Law}}~\cite{nigrini2012benford}.
          La référence pour appliquer la loi de Benford en audit, avec ses
          limites.
    \item \textbf{Cressey}~\cite{cressey1953}. L'origine du triangle de la
          fraude.
\end{itemize}

\section*{Pour aller plus loin}

\begin{itemize}
    \item \textbf{Minka, \emph{Vers l'auditologie}} et l'article sur la
          réponse \NEK{}~\cite{minka2026auditologie,minka2026nek}. Le
          programme de recherche dont le cours tire la cotation
          \OUI{}~/ \NON{}~/ \NEK{}~; prépublications sur HAL.
    \item \textbf{Le cours INF3226} (ENSPY), version complète de ce cours,
          pour qui veut approfondir chaque domaine.
\end{itemize}

    %% ============================================================================
%% Fichier  : annexes/Ann_F_contacts.tex
%% Titre    : Contacts et Ressources
%% Version  : 0.1 -- novembre 2026
%% ============================================================================
%% Sites donnés à titre indicatif ; seuls les sites connus avec certitude
%% sont indiqués, les autres sont à compléter par l'auteur.
%% ============================================================================

\chapter{Contacts et Ressources}
\label{ann:contacts}

Les adresses ci-dessous sont celles des sites institutionnels à la date de
rédaction~; elles peuvent changer et sont à vérifier avant usage.

\section*{Institutions camerounaises}

\begin{tableaucours}{L{4.4cm}L{4.8cm}L{3.0cm}}{%
    \tableauHeader{Institution} &
    \tableauHeader{Rôle dans le cours} &
    \tableauHeader{Site}
}
Contrôle Supérieur de l'État & Institution supérieure de contrôle,
contrôle des SI (Ch.~\ref{chap:institutions}) & · \\
Chambre des Comptes de la Cour suprême & Contrôle des comptes~; SI par
extension & · \\
ANTIC & Audit de sécurité des SI publics (loi n°~2010/012) &
\texttt{www.antic.cm} \\
Ministère des Finances & Régime financier, applications financières &
\texttt{www.minfi.gov.cm} \\
PSSFP, Université de Yaoundé~II & Formation des auditeurs des finances
publiques & · \\
\end{tableaucours}

\section*{Organisations internationales et professionnelles}

\begin{tableaucours}{L{4.4cm}L{4.8cm}L{3.0cm}}{%
    \tableauHeader{Organisation} &
    \tableauHeader{Apport} &
    \tableauHeader{Site}
}
INTOSAI & Normes ISSAI, GUID~5100 & \texttt{www.intosai.org} \\
Initiative de développement de l'INTOSAI (IDI) & Manuels de mise en
œuvre & \texttt{www.idi.no} \\
COSO & Référentiels de contrôle interne & \texttt{www.coso.org} \\
IIA & Normes de l'audit interne & \texttt{www.theiia.org} \\
ISACA & COBIT, certification CISA & \texttt{www.isaca.org} \\
\end{tableaucours}

\section*{Le cours}

\begin{itemize}
    \item NDAP SUKULU~: portail des fiches et des labs, code de cours
          \codeUE.
    \item Contact de l'enseignant~:
          \href{mailto:maletyon@proton.me}{maletyon@proton.me}.
\end{itemize}


\backmatter

    \printbibliography[heading=bibintoc, title={Références bibliographiques}]

    \printindex

% ==============================================================================
% 4ÈME DE COUVERTURE
% ==============================================================================
\clearpage
\thispagestyle{empty}
\pagecolor{CouleurPrimaire}

\begin{tikzpicture}[remember picture, overlay]
    \fill[CouleurAccent]
        (current page.south west) rectangle
        ([yshift=1.8cm] current page.south east);
    \fill[CouleurAccent]
        ([yshift=-1.4cm] current page.north west) rectangle
        ([yshift=-1.7cm] current page.north east);
\end{tikzpicture}

\vspace*{2.5cm}

\begin{center}
    {\color{CouleurAccent}\rule{0.6\linewidth}{1pt}}\\[1.5ex]
    {\color{white}\Large\bfseries À propos de ce cours}\\[1.5ex]
    {\color{CouleurAccent}\rule{0.6\linewidth}{1pt}}
\end{center}

\vspace{1cm}

\begin{center}
\begin{minipage}{0.75\textwidth}
\color{white!85!gray}\setstretch{1.4}
Quand l'administration financière s'informatise, le contrôle interne ne
disparaît pas~: il passe du dossier papier aux droits d'accès, aux
contrôles automatiques et aux journaux. Ce cours apprend à l'auditeur
public à le suivre, à l'évaluer et à conclure sur son efficacité.

\medskip

Il progresse en six parties, précédées d'un prologue~: le contrôle
interne et son évaluation, les institutions et les normes, la démarche de
mission et la mesure de l'incertitude, l'évaluation domaine par domaine
(accès, exploitation, applications, données, projets), un cas intégrateur,
et une ouverture sur la recherche en auditologie.

\medskip

Chaque notion est ancrée dans les normes internationales (INTOSAI, COSO,
IIA, ISACA, ISO) et dans le droit camerounais, à commencer par le décret
qui confie au Contrôle Supérieur de l'État le contrôle des systèmes
d'information. Le cours s'adresse aux vérificateurs, juges des comptes,
inspecteurs, contrôleurs financiers et auditeurs internes du secteur
public.
\end{minipage}
\end{center}

\vfill

\begin{center}
    {\color{CouleurAccent}\rule{0.7\linewidth}{1.2pt}}\\[1.2ex]
    {\color{white}\large\bfseries MINKA MI NGUIDJOÏ Thierry Emmanuel}\\[0.4ex]
    {\color{white!80!gray}\small
        CISA~·~CISM~·~CGEIT~·~CRISC~·~CDPSE~·~ISO 27001~·~ISO 22301~·~ISO 9001}\\[0.3ex]
    {\color{white!70!gray}\footnotesize\itshape
        Chercheur au laboratoire d'Ingénierie Mathématiques et des Systèmes d'Information}\\[0.3ex]
    {\color{white!70!gray}\footnotesize\itshape
        Chercheur~|~Cryptologue~|~Auditeur Senior~|~Expert Judiciaire~|~GRC \& Sécurité de l'Information}\\[1.2ex]
    {\color{CouleurAccent}\rule{0.7\linewidth}{1.2pt}}
\end{center}

\vspace{1cm}

\afterpage{\nopagecolor}
\end{document}
